\documentclass[11pt,reqno]{amsart}

% ================================================================
% Packages
% ================================================================
\usepackage{amsmath,amssymb,amsthm}
\usepackage{mathrsfs}
\usepackage{enumerate}
\usepackage[shortlabels]{enumitem}
\usepackage{booktabs}
\usepackage{array}
\usepackage{microtype}
\usepackage{tikz-cd}
\usepackage[dvipsnames]{xcolor}
\usepackage[colorlinks=true,
  linkcolor=blue!60!black,
  citecolor=green!40!black,
  urlcolor=blue!60!black]{hyperref}

% ================================================================
% Theorem environments
% ================================================================
\newtheorem{theorem}{Theorem}[section]
\newtheorem{proposition}[theorem]{Proposition}
\newtheorem{lemma}[theorem]{Lemma}
\newtheorem{corollary}[theorem]{Corollary}
\newtheorem{conjecture}[theorem]{Conjecture}
\theoremstyle{definition}
\newtheorem{definition}[theorem]{Definition}
\newtheorem{construction}[theorem]{Construction}
\newtheorem{example}[theorem]{Example}
\newtheorem{warning}[theorem]{Warning}
\theoremstyle{remark}
\newtheorem{remark}[theorem]{Remark}

% ================================================================
% Macros
% ================================================================
\newcommand{\cA}{\mathcal{A}}
\newcommand{\cB}{\mathcal{B}}
\newcommand{\cC}{\mathcal{C}}
\newcommand{\cD}{\mathcal{D}}
\newcommand{\cF}{\mathcal{F}}
\newcommand{\cH}{\mathcal{H}}
\newcommand{\cM}{\mathcal{M}}
\newcommand{\cO}{\mathcal{O}}
\newcommand{\cL}{\mathcal{L}}
\newcommand{\cW}{\mathcal{W}}
\newcommand{\cP}{\mathcal{P}}

\newcommand{\Ran}{\mathrm{Ran}}
\newcommand{\MC}{\mathrm{MC}}
\newcommand{\HH}{\mathrm{HH}}
\newcommand{\CE}{\mathrm{CE}}

\newcommand{\Ainf}{A_{\infty}}
\newcommand{\Eone}{\mathsf{E}_{1}}
\newcommand{\Etwo}{\mathsf{E}_{2}}
\newcommand{\Ethree}{\mathsf{E}_{3}}
\newcommand{\En}{\mathsf{E}_{n}}
\newcommand{\Einf}{\mathsf{E}_{\infty}}
\newcommand{\Pthree}{\mathsf{P}_{3}}

\newcommand{\Sym}{\mathrm{Sym}}
\newcommand{\Ass}{\mathrm{Ass}}
\newcommand{\Com}{\mathrm{Com}}
\newcommand{\Lie}{\mathrm{Lie}}

\DeclareMathOperator{\Hom}{Hom}
\DeclareMathOperator{\RHom}{RHom}
\DeclareMathOperator{\Res}{Res}
\DeclareMathOperator{\Aut}{Aut}
\DeclareMathOperator{\End}{End}
\DeclareMathOperator{\Conf}{Conf}
\DeclareMathOperator{\FM}{FM}
\DeclareMathOperator{\Spec}{Spec}
\DeclareMathOperator{\av}{av}
\DeclareMathOperator{\rk}{rk}
\DeclareMathOperator{\GRT}{GRT}
\DeclareMathOperator{\PBr}{PBr}

\newcommand{\fg}{\mathfrak{g}}
\newcommand{\fh}{\mathfrak{h}}

\newcommand{\gmod}{\mathfrak{g}^{\mathrm{mod}}_{\cA}}
\newcommand{\gEone}{\mathfrak{g}^{\Eone}_{\cA}}

\newcommand{\Barord}{\overline{B}^{\mathrm{ord}}}
\newcommand{\BarSig}{\overline{B}^{\Sigma}}
\newcommand{\Barch}{\overline{B}^{\mathrm{ch}}}

\newcommand{\CC}{\mathbb{C}}
\newcommand{\RR}{\mathbb{R}}
\newcommand{\ZZ}{\mathbb{Z}}
\newcommand{\PP}{\mathbb{P}}
\newcommand{\id}{\mathrm{id}}
\newcommand{\op}{\mathrm{op}}

\DeclareMathOperator{\Perf}{Perf}
\DeclareMathOperator{\Rep}{Rep}
\DeclareMathOperator{\Cl}{Cl}
\DeclareMathOperator{\Mod}{Mod}
\DeclareMathOperator{\Hol}{Hol}
\renewcommand{\SS}{\mathbb{S}}
\newcommand{\HHH}{\mathbb{H}}

\numberwithin{equation}{section}

% ================================================================
\begin{document}

\title[Chiral Yangians]
{On the Integrable Theory of Chiral Yangians\\
and Related Chiral Quantum Groups}

\author{Raeez Lorgat}
\address{Perimeter Institute for Theoretical Physics,
  Waterloo, ON N2L 2Y5, Canada \\
  Department of Physics, University of Waterloo,
  Waterloo, ON N2L 3G1, Canada}
\email{rlorgat@perimeterinstitute.ca}

\date{\today}

\begin{abstract}
For a chiral algebra $\cA$ on a smooth algebraic
curve~$X$, we define the \emph{ordered chiral homology}
$\int_X^{\mathrm{ord}} \cA$ as the derived pushforward of
the ordered factorisation $\cD$-module
$\cF^{\mathrm{ord}}(\cA)$ on the ordered Ran space
$\Ran^{\mathrm{ord}}(X) = \coprod_n U_n(X)$, the moduli
of ordered tuples of distinct points. The
Beilinson--Drinfeld Ran space $\Ran(X)$ is the
$\Sigma_\bullet$-coinvariant quotient, and the averaging
map
$\av \colon R\Gamma(\Ran^{\mathrm{ord}}(X),
\cF^{\mathrm{ord}}) \to R\Gamma(\Ran(X), \cF^{\Sigma})$
recovers the Beilinson--Drinfeld chiral homology.
For $\Einf$-chiral algebras (all standard vertex algebras),
the Kontsevich--Tamarkin formality of the little-disks operad
gives a quasi-isomorphism between the ordered and symmetric
theories; for genuinely $\Eone$-chiral algebras (Yangians,
Etingof--Kazhdan quantum vertex algebras), this formality
fails and the descent is lossy, depending on the choice of a
Drinfeld associator.

The main structural result is the chiral quantum group
equivalence: on the Koszul locus, three structures on an
$\Eone$-chiral algebra (the vertex $R$-matrix, the chiral
$\Ainf$-structure, and the chiral coproduct) determine each
other, forming an equivalence triangle mediated by the
ordered bar complex
$\Barord(\cA) = T^c(s^{-1}\bar{\cA})$ and its
deconcatenation coproduct.

Five computations span the $G/L/C/M$ shadow classification:
Heisenberg (class~$G$, regular KZ, ordered $=$ symmetric),
$V_k(\mathfrak{sl}_2)$ (class~$L$, KZ conformal blocks),
$\beta\gamma$ (class~$C$, nilpotent $R$-matrix),
Virasoro (class~$M$, irregular KZ connection with Stokes
data), and the Yangian $Y_\hbar(\mathfrak{sl}_2)$, which
is the first example where ordered chiral homology is
strictly richer than symmetric: the kernel $\ker(\av)$ is
non-trivial at every arity $n \geq 2$.

For $\Einf$-chiral algebras, the symmetric bar
$\BarSig(\cA)$ carries an $\Etwo$-coalgebra structure from
the curve geometry, and the Higher Deligne Conjecture
(Lurie, Francis) equips the derived chiral centre
$Z^{\mathrm{der}}_{\mathrm{ch}}(\cA)$ with an
$\Ethree$-algebra structure ($\Etwo + 1 = \Ethree$).
This $\Ethree$ structure does not extend to genuinely
$\Eone$-chiral algebras, where the symmetric bar does not
exist and the ordered bar yields only $\Etwo$ on the
derived centre.
For $V_k(\fg)$ at non-critical level, we construct a
filtered $\Ethree$-algebra
$\CE^{\mathrm{ch}}_\hbar(\fg)$ in $\cD$-modules on
$\Ran(X)$ whose $\Ethree$-deformation space is
$\hbar H^3(\fg)[[\hbar]]$. The same space classifies
the Costello--Francis--Gwilliam perturbative
Chern--Simons $\Ethree$-algebra; formal disk restriction
recovers the CFG algebra at the associated-graded and
$\Pthree$-bracket level. We conjecture that the two
formal deformation \emph{families} are isomorphic as
$\Ethree$-algebras; the matching of deformation spaces
is proved, but the chain-level identification is open.

The construction extends to genus~$1$: the ordered
chiral homology
$\int_{E_\tau}^{\mathrm{ord}} Y_\hbar(\mathfrak{sl}_2)$
on an elliptic curve $E_\tau$ is computed arity by arity,
with the KZ connection replaced by the
Knizhnik--Zamolodchikov--Bernard connection
(Bernard~\cite{Bernard88}, Felder~\cite{Felder94}) and the
$B$-cycle monodromy producing the quantum group parameter
$q = e^{2\pi i \hbar}$.
The genus-$1$ kernel $\ker(\av)$ is strictly larger than
at genus~$0$.
\end{abstract}

\maketitle
\tableofcontents

% ================================================================
\section{Introduction}
\label{sec:intro}

% ----------------------------------------------------------------
% CG DEFICIENCY OPENING
% ----------------------------------------------------------------

Beilinson and Drinfeld's chiral homology~\cite{BD04} is the
fundamental invariant of a chiral algebra on an algebraic curve.
For a chiral algebra $\cA$ on a smooth proper curve~$X$,
the chiral homology $\int_X \cA$ computes the space of
conformal blocks: the derived global sections of a factorisation
$\cD$-module on the Ran space $\Ran(X)$. The construction
requires that $\cA$ be a \emph{commutative} factorisation
algebra: the factorisation product, which merges local
observables at disjoint collections of points, must be
symmetric under permutation of the inputs.

Every vertex algebra in the standard landscape (Heisenberg,
affine Kac--Moody, Virasoro, $\cW$-algebras, $\beta\gamma$
systems, lattice algebras) satisfies this commutativity
hypothesis. The OPE $Y(a,z)Y(b,w)$ and $Y(b,w)Y(a,z)$ agree
(after clearing poles) by the locality axiom, and the
factorisation product descends to the symmetric quotient
$X^{(n)} = X^n / \Sigma_n$.

The hypothesis fails for the class of algebras that the
first volume of this series~\cite{Lorgat26I} identifies as
the \emph{genuinely $\Eone$-chiral} objects: the Yangians
$Y_\hbar(\fg)$, the Etingof--Kazhdan quantum vertex
algebras~\cite{EK00}, and more broadly any chiral algebra
whose operator product is governed by a non-trivial
vertex $R$-matrix $S(z)$ satisfying the quantum Yang--Baxter
equation. For these algebras, the factorisation product is
\emph{braided}, not commutative: swapping two groups of
inputs introduces a twist by $S(z)$.

The deficiency is not in the algebras but in the available
integration theory. There is no $\cD$-module obstruction
to working on ordered spaces: $\cD$-modules on $X^n$, on
the ordered Fulton--MacPherson compactification
$\overline{\FM}_n^{\mathrm{ord}}(X)$, and on Mok's
log-smooth compactifications are perfectly standard objects,
studied since Kashiwara and Malgrange. The correct
foundational object is the \emph{ordered Ran space}
$\Ran^{\mathrm{ord}}(X) = \coprod_n U_n(X)$, the moduli of
ordered tuples of distinct points. The Beilinson--Drinfeld
Ran space is the $\Sigma_\bullet$-coinvariant quotient
$\Ran(X) = \Ran^{\mathrm{ord}}(X) / \Sigma_\bullet$, and
the commutative factorisation structure on $\Ran(X)$ is the
equivariant shadow of the ordered factorisation structure on
$\Ran^{\mathrm{ord}}(X)$ that exists for any chiral algebra,
commutative or not.

This paper constructs the ordered factorisation $\cD$-module
$\cF^{\mathrm{ord}}(\cA)$ on $\Ran^{\mathrm{ord}}(X)$, defines
ordered chiral homology as the derived pushforward of this
single object, and establishes the basic properties of the
resulting invariant.

% ----------------------------------------------------------------
\subsection{Main results}

Three results emerge from the ordered theory.
The first constructs the $\cD$-module and defines ordered
chiral homology; the second separates it from topological
Hochschild; the third identifies the $\Ethree$ structure
that only the ordered picture sees.

\begin{theorem}[Ordered chiral homology]\label{thm:main-ord-ch}
Let $\cA$ be a chiral algebra on a smooth proper curve $X$
of genus~$g$, and let $r(z)$ denote the classical $r$-matrix
(the arity-$2$ projection of the Maurer--Cartan element
$\Theta_\cA$ in the $\Eone$ convolution algebra).
\begin{enumerate}[label=\textup{(\roman*)}]
\item \textup{(Ordered factorisation $\cD$-module on
  $\Ran^{\mathrm{ord}}(X)$.)}
  The OPE of $\cA$, extended meromorphically across the
  diagonals, defines a factorisation $\cD$-module
  $\cF^{\mathrm{ord}}(\cA)$ on the ordered Ran space
  $\Ran^{\mathrm{ord}}(X) = \coprod_n U_n(X)$. Its
  restriction to the $n$-th stratum is the holonomic
  $\cD$-module
  \begin{equation}\label{eq:fact-dmod}
    \cF^{\mathrm{ord}}(\cA)\big|_{U_n}
    = \cA^{\boxtimes n}\big|_{U_n},
    \qquad
    \cF_n^{\mathrm{ord}}(\cA)
    := j_* j^* (\cA^{\boxtimes n})
    \in \cD\text{-}\mathrm{mod}(X^n),
  \end{equation}
  where $j \colon U_n = X^n \setminus
  \bigcup_{1 \leq a < b \leq n}
  \Delta_{ab} \hookrightarrow X^n$ is the inclusion of
  the ordered configuration space. No $\Sigma_n$-equivariance
  is imposed. The factorisation product, the KZ connection,
  and the chiral differential are encoded in the single
  object $\cF^{\mathrm{ord}}(\cA)$ on
  $\Ran^{\mathrm{ord}}(X)$.
\item \textup{(Ordered chiral homology.)}
  The \emph{ordered chiral homology} of $\cA$ over $X$
  is the derived pushforward
  \begin{equation}\label{eq:ord-ch-hom}
    \int_X^{\mathrm{ord}} \cA
    \;:=\;
    R\Gamma_{\mathrm{dR}}\bigl(
    \Ran^{\mathrm{ord}}(X),\;
    \cF^{\mathrm{ord}}(\cA)\bigr).
  \end{equation}
  The arity-by-arity decomposition
  $\bigoplus_{n \geq 0}
  R\Gamma_{\mathrm{dR}}(X^n,\,
  \cF_n^{\mathrm{ord}}(\cA))$
  computes this pushforward stratum by stratum.
\item \textup{(Symmetric descent.)}
  The quotient
  $q \colon \Ran^{\mathrm{ord}}(X) \to \Ran(X)$
  induces the averaging map
  \begin{equation}\label{eq:descent}
    \int_X^{\mathrm{ord}} \cA
    = R\Gamma\bigl(\Ran^{\mathrm{ord}}(X),\,
    \cF^{\mathrm{ord}}\bigr)
    \;\xrightarrow{\;\av\;}
    R\Gamma\bigl(\Ran(X),\, \cF^{\Sigma}\bigr)
    = \int_X \cA,
  \end{equation}
  which is a quasi-isomorphism for $\Einf$-chiral $\cA$.
  For $\Eone$-chiral $\cA$, the descent is lossy:
  the kernel of $\av \colon
  R\Gamma_{\mathrm{dR}}(X^n,
  \cF_n^{\mathrm{ord}})
  \twoheadrightarrow
  R\Gamma_{\mathrm{dR}}(X^n,
  \cF_n^{\mathrm{ord}})_{\Sigma_n}$
  is non-trivial at $n \geq 3$, carrying $\GRT_1$-dependent
  data from the Drinfeld associator.
\end{enumerate}
\end{theorem}

\begin{theorem}[Topological Hochschild vs.\ ordered
chiral homology of $D^\times$]\label{thm:hh-vs-ch}
Let $\cA$ be a chiral algebra on a curve $X$, and let
$D^\times = \Spec\, \CC((z))$ be the punctured formal disk.
\begin{enumerate}[label=\textup{(\roman*)}]
\item \textup{(Topological Hochschild.)}
  Viewing $\cA$ as an $\Eone$-topological algebra by
  forgetting the holomorphic structure, the factorisation
  homology $\int_{S^1} \cA$ is the Hochschild homology
  $\HH_*(\cA)$, computed by the cyclic bar complex with the
  Hochschild differential using the associative product
  $m_2 \colon \cA \otimes \cA \to \cA$ and the higher
  $\Ainf$ operations $m_k$.
\item \textup{(Ordered chiral homology of $D^\times$.)}
  The ordered chiral chain complex at arity~$n$ on $D^\times$ is
  \begin{equation}\label{eq:ch-cx-Dx}
    \cC_n^{\mathrm{ord}}(D^\times, \cA)
    = \Bigl(
    \Omega^*_{\log}\bigl(
    \overline{\FM}_n^{\mathrm{ord}}(D^\times)\bigr)
    \otimes (s^{-1}\bar{\cA})^{\otimes n},
    \;\; d_{\mathrm{total}}
    \Bigr)
  \end{equation}
  where $\Omega^*_{\log}$ denotes logarithmic differential
  forms on the Fulton--MacPherson compactification,
  $s^{-1}$ is the desuspension
  \textup{(}$|s^{-1}v| = |v| - 1$\textup{)},
  $\bar{\cA} = \ker(\varepsilon)$ is the augmentation ideal,
  and
  \begin{equation}\label{eq:d-total}
    d_{\mathrm{total}}
    = d_{\mathrm{dR}}
    + d_{\mathrm{bar}}
    + \nabla_{\mathrm{KZ}}
  \end{equation}
  with $d_{\mathrm{dR}}$ the de~Rham differential,
  $d_{\mathrm{bar}}$ the bar differential encoding the OPE, and
  \begin{equation}\label{eq:kz}
    \nabla_{\mathrm{KZ}}
    = \sum_{i < j} r_{ij}(z_i - z_j) \, dz_{ij},
    \qquad
    dz_{ij} = d(z_i - z_j),
  \end{equation}
  the Knizhnik--Zamolodchikov connection.
  The Arnold forms
  $\omega_{ij} = d\log(z_i - z_j)$ are generators of the
  cohomology algebra $H^*(\Conf_n(\CC))$ and appear as
  coefficients in the bar differential $d_{\mathrm{bar}}$;
  they are \emph{not} the connection $1$-form of $\nabla_{\mathrm{KZ}}$.
\item \textup{(Formality bridge.)} For $\Einf$-chiral $\cA$,
  the formality bridge
  \textup{(}Theorem~\textup{\ref{thm:formality-bridge}}%
  \textup{)} gives a quasi-isomorphism
  $\cC_n^{\mathrm{ord}}(D^\times, \cA)
  \xrightarrow{\sim} \HH_*(\cA)$:
  ordered and symmetric chiral homologies agree as chain
  complexes, though the transferred $\Ainf$-structure may
  carry non-trivial higher operations for class $L/C/M$.
  For genuinely $\Eone$-chiral $\cA$ \textup{(}where
  $S(z) \neq \id$\textup{)}, this quasi-isomorphism does
  not exist.
\end{enumerate}
\end{theorem}

\begin{theorem}[The $\Ethree$-algebra and Chern--Simons]
\label{thm:e3-cs}
Let $\fg$ be a simple finite-dimensional Lie algebra and
$V_k(\fg)$ the affine Kac--Moody vertex algebra at level~$k$.
\begin{enumerate}[label=\textup{(\roman*)}]
\item \textup{(The $\Ethree$ structure.)}
  The curve geometry provides the $\Etwo$ structure on
  $\BarSig(V_k(\fg))$
  \textup{(}Warning~\textup{\ref{warn:e1-vs-e2-source}}\textup{)}.
  By the Higher Deligne Conjecture
  \textup{(}Lurie~\cite{HA}, Francis~\cite{Francis2013}\textup{)},
  the Hochschild cohomology of this $\Etwo$-coalgebra
  \begin{equation}\label{eq:e3-center}
    Z^{\mathrm{der}}_{\mathrm{ch}}(V_k(\fg))
    = \HH^*\bigl(\BarSig(V_k(\fg)),\;
    \BarSig(V_k(\fg))\bigr)
  \end{equation}
  carries a natural $\Ethree$-algebra structure
  \textup{(}$\Etwo + 1 = \Ethree$\textup{)}.
\item \textup{(Deformation classification.)}
  The space of $\Ethree$-deformations of
  $Z^{\mathrm{der}}_{\mathrm{ch}}(V_k(\fg))$
  is non-canonically isomorphic to
  $\hbar H^3(\fg)[[\hbar]]$. For simple~$\fg$,
  $H^3(\fg)$ is one-dimensional, so there is a single
  deformation parameter at each order in $\hbar$.
\item \textup{(Comparison.)}
  The same $\hbar H^3(\fg)[[\hbar]]$ classifies, by
  Costello--Francis--Gwilliam~\cite{CFG26}:
  perturbative Chern--Simons quantisations,
  filtered $\Ethree$-algebras deforming $C^*(\fg)$,
  braided monoidal deformations of
  $\mathrm{Rep}_{\mathrm{fin}}(\fg)$, and
  quasi-triangular quasi-Hopf deformations of $U(\fg)$;
  and by Etingof--Kazhdan~\cite{EK96,EK00}:
  quantum vertex $R$-matrices on $V(\fg, K)$, forming
  a $\GRT_1$-torsor.
  All six classification spaces are in natural bijection.
\item \textup{(Topological enhancement.)}
  By Khan--Zeng~\cite{KhanZeng25}, the Sugawara Virasoro
  element at non-critical level upgrades the
  holomorphic-topological symmetry of the $3$d gauge theory
  to fully topological symmetry.
  Khan--Zeng~\cite{KhanZeng25} prove this at the
  classical \textup{(}PVA\textup{)} level; the quantum
  extension holds because the Sugawara element exists as
  an exact operator in $V_k(\fg)$ at every non-critical
  level, satisfying the full Virasoro commutation relations
  without quantum corrections to the BRST-exactness
  argument \textup{(}see
  Proposition~\ref{prop:khan-zeng-topological}\textup{)}.
  At the critical level $k = -h^\vee$, the
  Sugawara element is undefined and the topological
  enhancement does not apply; the topological nature of
  the perturbative $\Ethree$ at
  $k = -h^\vee + \hbar\lambda_1 + \cdots$ holds
  order-by-order in $\hbar$
  \textup{(}since the Sugawara element exists at any
  nonzero departure from critical level\textup{)}.
\end{enumerate}
\end{theorem}

% ----------------------------------------------------------------
\subsection{Architecture of the paper}

The logical sequence is forced by the construction.
One needs the $\Eone/\Einf$ hierarchy and the ordered Ran
space $\Ran^{\mathrm{ord}}(X)$
(\S\ref{sec:background}) before building the ordered
factorisation $\cD$-module
$\cF^{\mathrm{ord}}(\cA)$ on $\Ran^{\mathrm{ord}}(X)$
(\S\ref{sec:dmod-ordered}),
this single $\cD$-module before defining ordered chiral
homology as its derived pushforward and establishing
symmetric descent to the Beilinson--Drinfeld theory
(\S\ref{sec:ord-ch-hom}), the
definition before separating it from topological Hochschild
through four levels of formality (\S\ref{sec:hh-vs-ch}),
and the formality analysis before identifying the
$\Ethree$ structure via Chern--Simons theory
(\S\ref{sec:e3}). The examples in \S\ref{sec:examples}
cover all four shadow classes $G/L/C/M$.
Section~\ref{sec:comparison} places the construction
against Ben-Zvi--Brochier--Jordan and Latyntsev.
Section~\ref{sec:chiral-quantum-groups} establishes the
chiral quantum group equivalence: vertex $R$-matrices,
chiral $\Ainf$-structures, and chiral coproducts are
three projections of the ordered bar complex.
Section~\ref{sec:elliptic-ordered} extends the construction
to genus~$1$: the ordered chiral homology
$\int_{E_\tau}^{\mathrm{ord}} Y_\hbar(\mathfrak{sl}_2)$
is computed arity by arity on an elliptic curve, where
the KZB connection replaces the KZ connection and the
$B$-cycle monodromy produces the quantum group parameter
$q = e^{2\pi i \hbar}$.
Section~\ref{sec:conclusion} collects open problems.


% ================================================================
\section{Background: chiral algebras and the $\Eone/\Einf$
hierarchy}
\label{sec:background}

The ordered theory rests on a single structural observation:
every chiral algebra has two independent invariants (chirality
hierarchy and topological ladder), and the standard
Beilinson--Drinfeld theory sees only their $\Einf$-equivariant
shadow. This section recalls the two hierarchies and the
bridge theorem that identifies $\Eone$-chiral algebras with
$\Etwo$-topological data.

% ----------------------------------------------------------------
\subsection{Two orthogonal axes}
\label{subsec:two-axes}

A chiral algebra carries two independent pieces of structure:
how commutative its OPE is (the chirality hierarchy,
$\Einf \subset \cP_\infty \subset \Eone$) and how many
real dimensions its configuration spaces see (the
topological ladder, $\Eone \subset \Etwo \subset \cdots
\subset \Einf$). Conflating the two is the source of most
errors in the subject; separating them is the starting
point of the ordered theory.

\begin{definition}[Chirality hierarchy]
\label{def:chirality-hierarchy}
For chiral algebras on a fixed algebraic curve $X$,
the \emph{chirality hierarchy}
$\Einf\text{-chiral} \subset
 \cP_\infty\text{-chiral} \subset
 \Eone\text{-chiral}$
records the commutativity of the OPE:
\begin{enumerate}[label=\textup{(\roman*)}]
\item An $\Einf$\emph{-chiral algebra} has symmetric OPE:
  the vertex operators $Y(a,z)$ and $Y(b,w)$ commute
  after clearing poles. Every standard vertex algebra
  (Heisenberg, affine Kac--Moody, Virasoro,
  $\cW$-algebras, lattice algebras) is
  $\Einf$-chiral. The symmetry means the factorisation
  product descends to $X^{(n)} = X^n / \Sigma_n$.
\item An $\Eone$\emph{-chiral algebra} has braided OPE:
  the vertex operators satisfy $S$-locality
  \begin{equation}\label{eq:s-locality}
    (z - w)^N Y(a,z) Y(b,w)
    = (z - w)^N Y(b,w) Y(a,z) \cdot S(z - w)
  \end{equation}
  where $S(z)$ is a vertex $R$-matrix satisfying the
  quantum Yang--Baxter equation, unitarity, the shift
  condition, and the hexagon axiom
  \textup{(}Etingof--Kazhdan~\cite{EK00}\textup{)}.
  The factorisation product lives on ordered
  configurations $X^n$, not on $X^{(n)}$.
\end{enumerate}
\end{definition}

\begin{warning}[Multiple notions of $\Eone$-chiral]
\label{warn:multiple-e1-notions}
At least five notions coexist in the literature:
\textup{(A)}~strict ChirAss-algebra,
\textup{(B)}~$\Ainf$-algebra in the chiral endomorphism operad
\textup{(}structure maps meromorphic in spectral parameters\textup{)},
\textup{(C)}~Etingof--Kazhdan quantum vertex algebra
\textup{(}five axioms including $S$-locality\textup{)},
\textup{(D)}~$\Ainf$-algebra object in $\Eone$-chiral algebras
\textup{(}double $\Ainf$\textup{)}, and
\textup{(E)}~factorisation $\cD$-module on
$\Ran^{\mathrm{ord}}(X)$.
On the Koszul locus, \textup{(B)} and~\textup{(C)} determine
each other via a Drinfeld associator $\Phi$
\textup{(}Theorem~\textup{\ref{thm:chiral-qg-equiv}}\textup{)};
\textup{(A)} is a special case of~\textup{(B)};
\textup{(E)} is the $\cD$-module realisation of~\textup{(B)}.
Each notion admits its own derived centre construction, and
the output operad depends on which construction is used
\textup{(}see~\cite[Warning~2.21]{Lorgat26I}\textup{)}.
In this paper, Definition~\textup{\ref{def:chirality-hierarchy}}
uses Notion~\textup{(C)} for the algebraic definition and
Notion~\textup{(E)} for the geometric realisation.
\end{warning}

\begin{definition}[The ordered Ran space]
\label{def:ran-ord}
For a smooth algebraic curve $X$, the \emph{ordered Ran
space} is
\[
  \Ran^{\mathrm{ord}}(X) = \coprod_{n \geq 1} U_n(X), \qquad
  U_n(X) = X^n \setminus \bigcup_{i<j} \Delta_{ij},
\]
the moduli of ordered finite subsets of $X$: a point of
$\Ran^{\mathrm{ord}}(X)$ is an ordered tuple
$(x_1, \ldots, x_n)$ of distinct points of $X$.
The \emph{ordered factorisation structure} on
$\Ran^{\mathrm{ord}}(X)$ encodes the merging of ordered
subsets: for $(x_1, \ldots, x_k)$ and
$(x_{k+1}, \ldots, x_n)$ with disjoint underlying sets,
the concatenation $(x_1, \ldots, x_n)$ defines a map
\[
  \mathrm{union}_{k,n-k} \colon
  \Ran^{\mathrm{ord}}(X)_k
  \times_{\mathrm{disj}}
  \Ran^{\mathrm{ord}}(X)_{n-k}
  \to \Ran^{\mathrm{ord}}(X)_n.
\]
\end{definition}

\begin{definition}[The Ran space as coinvariant quotient]
\label{def:ran-space}
The Beilinson--Drinfeld \emph{Ran space}
$\Ran(X) = \coprod_{n \geq 1} X^{(n)}$, where
$X^{(n)} = X^n / \Sigma_n$ is the $n$-th symmetric product,
is the $\Sigma_\bullet$-coinvariant quotient:
\[
  \Ran(X) = \Ran^{\mathrm{ord}}(X) / \Sigma_\bullet.
\]
The factorisation structure on $\Ran(X)$ is the
$\Sigma_\bullet$-equivariant shadow of the ordered
factorisation structure on $\Ran^{\mathrm{ord}}(X)$: for
disjoint finite subsets $I, J \subset X$, the merging map
$\Ran(X)_I \times \Ran(X)_J \to \Ran(X)_{I \cup J}$ is
the descent of $\mathrm{union}_{k,n-k}$ through the
symmetric group quotient.
\end{definition}

\begin{remark}[Chiral algebras, $\Etwo$, and the role of the curve]
\label{rem:chiral-e2}
An $\Einf$-chiral algebra $\cA$ on a complex curve $X$ is a
factorisation algebra on $\Ran(X)$; for an $\Eone$-chiral
algebra, the factorisation $\cD$-module
$\cF^{\mathrm{ord}}(\cA)$ lives on the ordered Ran space
$\Ran^{\mathrm{ord}}(X)$
(Definition~\ref{def:ran-ord}), and descends to
$\Ran(X)$ only when the braiding is trivial.
In either case, $\cA$ is \emph{not} an
$\Etwo$-algebra in the operadic sense. The $\Etwo$
structure lives on the symmetric bar $\BarSig(\cA)$,
provided by the curve geometry
(Warning~\ref{warn:e1-vs-e2-source}).
The holomorphic data
(the spectral parameter $z$, the $r$-matrix $r(z)$,
the KZ connection) refines this topological $\Etwo$
structure; forgetting the complex structure of $X$
recovers the bare $\Etwo$-coalgebra.
\end{remark}

\begin{warning}[Bar-cobar gives $\Eone$; the curve gives $\Etwo$]
\label{warn:e1-vs-e2-source}
The abstract bar-cobar adjunction for $\Eone$-algebras is:
\[
\mathrm{Bar} \colon \Eone\text{-}\mathrm{Alg}
\rightleftarrows
\Eone\text{-}\mathrm{CoAlg}
\colon \mathrm{Cobar}.
\]
The output is an $\Eone$-coalgebra with the
\emph{deconcatenation} coproduct:
$\Delta(a_1 \otimes \cdots \otimes a_n)
= \sum_{k=0}^{n} (a_1 \otimes \cdots \otimes a_k)
\otimes (a_{k+1} \otimes \cdots \otimes a_n)$.
This is $\Eone$ \textup{(}it respects the linear
ordering of tensor factors\textup{)}.

The symmetric bar $\BarSig(\cA)$ of a chiral algebra
on a curve $X$ carries the \emph{coshuffle} coproduct:
$\Delta(a_1 \cdots a_n)
= \sum_{I \sqcup J = \{1,\ldots,n\}} a_I \otimes a_J$
\textup{(}sum over all unordered bipartitions\textup{)}.
This is $\Etwo$/$\Einf$: it uses the
$\Sigma_n$-symmetry from unordered configuration spaces
$\Conf_k(X)$ of the curve, whose real $2k$-dimensional
topology carries the $\FM_2 \simeq \Etwo$ operad.

The passage from $\Barord$ \textup{(}$\Eone$,
deconcatenation\textup{)} to $\BarSig$
\textup{(}$\Etwo$, coshuffle\textup{)} is the
$\Sigma_n$-coinvariant quotient. This quotient
\emph{introduces} the curve geometry: it replaces
ordered configurations with unordered ones, and
the deconcatenation with the coshuffle.

Without the curve, bar-cobar gives $\Eone$, and
Higher Deligne gives the centre $\Etwo$
\textup{(}$\Eone + 1$; this is the classical Deligne
conjecture for associative algebras\textup{)}.
With the curve, the bar is $\Etwo$, and
Higher Deligne gives the centre $\Ethree$
\textup{(}$\Etwo + 1$\textup{)}.
The extra operadic level is the curve.
\end{warning}

\begin{definition}[Topological ladder]
\label{def:top-ladder}
The \emph{topological ladder}
$\Eone \subset \Etwo \subset \Ethree \subset \cdots
 \subset \Einf$
is indexed by real manifold dimension:
\begin{enumerate}[label=\textup{(\roman*)}]
\item $n = 1$: operations on $\Conf_k(\RR)$.
  $\Eone$-algebras $= \Ainf$-algebras $=$ homotopy
  associative algebras.
\item $n = 2$: operations on
  $\Conf_k(\RR^2) \cong \Conf_k(\CC)$.
  $\Etwo$-algebras govern braided monoidal categories.
\item $n = 3$: operations on $\Conf_k(\RR^3)$.
  $\Ethree$-algebras arise from perturbative Chern--Simons
  theory (Costello--Francis--Gwilliam~\cite{CFG26}).
\end{enumerate}
\end{definition}

\begin{theorem}[Bridge; {\cite[Theorem~1.1(ii)]{Lorgat26I}}]
\label{thm:bridge}
The chiral bar-cobar adjunction on a complex curve $X$
is the holomorphic refinement of the $n = 2$ chapter of
$\En$ Koszul duality. The symmetric bar $\BarSig(\cA)$
of a chiral algebra $\cA$ on $X$ is an
$\Etwo$-coalgebra \textup{(}from the coshuffle coproduct
on $\Conf_k(X) \cong \Conf_k(\RR^2)$\textup{)}, not
merely an $\Eone$-coalgebra \textup{(}which is what the
abstract bar-cobar adjunction would give\textup{)}.
For an $\Eone$-chiral algebra, the braid-group action
on $\Conf_k(\CC)$ is non-trivial \textup{(}the
$R$-matrix braiding\textup{)}; for an $\Einf$-chiral
algebra, it factors through the symmetric group.
In both cases, the $\Etwo$ structure on $\BarSig$
comes from the curve, not from bar-cobar.
\end{theorem}

\begin{remark}[The bridge at the level of operads]
\label{rem:bridge-operads}
The real Fulton--MacPherson compactification
$\overline{\FM}_n^{\mathrm{ord}}(\RR)$ is isomorphic
to the Stasheff associahedron $K_n$
(Sinha~\cite{Sinha04}).
The inclusion $\RR \hookrightarrow \CC$ induces an
embedding
\begin{equation}\label{eq:assoc-embed}
  K_n = \overline{\FM}_n^{\mathrm{ord}}(\RR)
  \hookrightarrow
  \overline{\FM}_n^{\mathrm{ord}}(\CC)
\end{equation}
of the $\Ainf$ operad into the ordered Fulton--MacPherson
operad. An $\Eone$-chiral algebra structure is an
operadic map
$\Ainf \to \End^{\mathrm{ch,ord}}_\cA$
factoring through~\eqref{eq:assoc-embed}.
The image of $K_n$ in $\overline{\FM}_n^{\mathrm{ord}}(\CC)$
is real-codimension~$n$: the $\Ainf$ operations are a
thin slice of the full chiral operations.
\end{remark}


% ================================================================
\section{$\cD$-modules on ordered configuration spaces}
\label{sec:dmod-ordered}

The integration theory for $\Einf$-chiral algebras lives on
symmetric products $X^{(n)} = X^n/\Sigma_n$, but the
$\Eone$ theory requires ordered configurations $X^n$.
What are the $\cD$-modules on $X^n$ that encode the OPE,
and what happens to them at the diagonals? The answer
is the ordered factorisation $\cD$-module: a holonomic
$\cD$-module on $X^n$ whose singularities along the
diagonals are precisely the OPE poles.

% ----------------------------------------------------------------
\subsection{Ordered configuration spaces and their
compactifications}
\label{subsec:config}

The configuration space $U_n(X) = X^n \setminus
\bigcup \Delta_{ij}$ and its Fulton--MacPherson
compactification are the geometric substrate for
ordered chiral homology. The compactification resolves
the diagonals by blowup, producing boundary strata
indexed by ordered collision patterns: rooted planar trees
with $n$ labelled leaves.

\begin{definition}[Ordered configuration space]
\label{def:ord-conf}
Let $X$ be a smooth algebraic curve over $\CC$.
For $n \geq 0$, the \emph{ordered configuration space} is
the smooth quasi-projective variety
\begin{equation}\label{eq:ord-conf}
  U_n = U_n(X)
  := X^n \setminus \bigcup_{1 \leq i < j \leq n} \Delta_{ij}
\end{equation}
where $\Delta_{ij} = \{(x_1, \ldots, x_n) \in X^n :
x_i = x_j\}$ is the $(i,j)$-diagonal.
For $n = 0$, $U_0 = \{\varnothing\} = \Spec\,\CC$.
For $n = 1$, $U_1 = X$.
\end{definition}

\begin{definition}[Ordered FM compactification]
\label{def:ord-fm}
The \emph{ordered Fulton--MacPherson compactification}
$\overline{\FM}_n^{\mathrm{ord}}(X)$ is the real oriented
blowup of $X^n$ along all diagonals $\Delta_{ij}$
($i < j$), iterated in order of increasing codimension.
It is a smooth manifold with corners whose interior is
$U_n(X)$ and whose boundary strata are indexed by
\emph{ordered collision patterns}: rooted planar trees
$T$ with $n$ labelled leaves, where each internal vertex
records a group of points colliding simultaneously.

For each ordered collision pattern $T$ with internal
vertices $v_1, \ldots, v_r$ of valences $k_1, \ldots, k_r$,
the corresponding boundary stratum is
\begin{equation}\label{eq:fm-stratum}
  \partial_T \overline{\FM}_n^{\mathrm{ord}}(X)
  \cong
  \prod_{i=1}^r
  \overline{\FM}_{k_i}^{\mathrm{ord}}(T_{x_i} X)
  \;\times\;
  \overline{\FM}_{n'}^{\mathrm{ord}}(X)
\end{equation}
where $T_{x_i} X \cong \CC$ is the tangent space at the
collision point $x_i$, and $n' = n - \sum_i (k_i - 1)$ is
the number of surviving points after collisions.
\end{definition}

\begin{remark}[$\cD$-modules on ordered spaces]
\label{rem:dmod-standard}
The category $\cD\text{-}\mathrm{mod}(X^n)$ of
$\cD$-modules on $X^n$ is the standard abelian category of
quasi-coherent sheaves with flat connection
(Kashiwara~\cite{Kashiwara03},
Hotta--Takeuchi--Tanisaki~\cite{HTT08}).
The diagonals $\Delta_{ij} \subset X^n$ are smooth
divisors; the $*$-extension $j_*$ and the $!$-extension
$j_!$ along $j \colon U_n \hookrightarrow X^n$ are the
standard meromorphic and distributional extensions.
The classical Riemann--Hilbert correspondence applies
to regular holonomic $\cD$-modules on $X^n$; for
irregular holonomic $\cD$-modules, one needs the
irregular Riemann--Hilbert correspondence of
D'Agnolo--Kashiwara~\cite{DAgnoloKashiwara16}.
Both cases arise in the chiral setting
(see Proposition~\ref{prop:fact-dmod-props} and
Warning~\ref{warn:regularity}).
\end{remark}


% ----------------------------------------------------------------
\subsection{The ordered factorisation $\cD$-module}
\label{subsec:fact-dmod}

Given the configuration spaces and their compactifications,
one extends the external product $\cA^{\boxtimes n}$
meromorphically across the diagonals. The resulting
$*$-extension $j_* j^* (\cA^{\boxtimes n})$ is the ordered
factorisation $\cD$-module: a holonomic $\cD$-module on
$X^n$ whose singularity type along each diagonal $\Delta_{ij}$
is dictated by the OPE pole order of $\cA$.

\begin{construction}[Ordered factorisation $\cD$-module]
\label{constr:fact-dmod}
Let $\cA$ be a chiral algebra on $X$.
The \emph{external product} $\cA^{\boxtimes n}$ is the
$\cD$-module on $X^n$ defined by
\begin{equation}\label{eq:external-prod}
  \cA^{\boxtimes n}
  := \cA \boxtimes \cA \boxtimes \cdots \boxtimes \cA
  = \mathrm{pr}_1^* \cA \otimes_{\cO_{X^n}}
    \mathrm{pr}_2^* \cA \otimes_{\cO_{X^n}} \cdots
    \otimes_{\cO_{X^n}} \mathrm{pr}_n^* \cA
\end{equation}
where $\mathrm{pr}_i \colon X^n \to X$ is the
$i$-th projection.

The \emph{ordered factorisation $\cD$-module} at arity $n$
is the $*$-extension:
\begin{equation}\label{eq:fact-dmod-def}
  \cF_n^{\mathrm{ord}}(\cA)
  := j_* j^* (\cA^{\boxtimes n})
  \in \cD\text{-}\mathrm{mod}(X^n).
\end{equation}
Concretely, sections of $\cF_n^{\mathrm{ord}}$ over an
open $V \subset X^n$ are sections of $\cA^{\boxtimes n}$
over $V \cap U_n$ that extend meromorphically across the
diagonals.  The poles along $\Delta_{ij}$ are the OPE poles:
if the OPE $Y(a, z) b = \sum_{k \geq -N} a_{(k)} b \,
z^{-k-1}$ has poles of order at most $N$ for all
$a, b \in \cA$, then sections of $\cF_n^{\mathrm{ord}}$
have poles of order at most $N$ along each $\Delta_{ij}$.
\end{construction}

\begin{definition}[Ordered factorisation $\cD$-module
on $\Ran^{\mathrm{ord}}(X)$]
\label{def:fact-dmod-ran}
For a chiral algebra $\cA$ on $X$, the \emph{ordered
factorisation $\cD$-module} is the $\cD$-module on
$\Ran^{\mathrm{ord}}(X)$
(Definition~\ref{def:ran-ord}) whose restriction to the
$n$-th stratum $U_n(X)$ is
\[
  \cF^{\mathrm{ord}}(\cA)\big|_{U_n}
  = \cA^{\boxtimes n}\big|_{U_n}
\]
with the factorisation product: for the merging map
$\mathrm{union}_{k,n-k}$, the factorisation isomorphism
\[
  \mathrm{union}_{k,n-k}^* \cF^{\mathrm{ord}}\big|_{U_n}
  \;\simeq\;
  \cF^{\mathrm{ord}}\big|_{U_k}
  \boxtimes
  \cF^{\mathrm{ord}}\big|_{U_{n-k}}
\]
is the OPE (the meromorphic extension across the diagonals).
The KZ connection endows $\cF^{\mathrm{ord}}$ with a flat
connection on each stratum, compatible with the factorisation
maps. The arity-by-arity $\cD$-modules
$\cF_n^{\mathrm{ord}}(\cA)$ of
Construction~\ref{constr:fact-dmod} are the restrictions
of this single object to the strata of
$\Ran^{\mathrm{ord}}(X)$.
\end{definition}

\begin{proposition}[Properties of the ordered
factorisation $\cD$-module]
\label{prop:fact-dmod-props}
Let $\cA$ be a chiral algebra on a smooth proper curve $X$.
\begin{enumerate}[label=\textup{(\roman*)}]
\item \textup{(Holonomicity and singularity type.)}
  $\cF_n^{\mathrm{ord}}(\cA)$ is a holonomic
  $\cD$-module on $X^n$.
  The singularity type along $\bigcup \Delta_{ij}$ depends
  on the pole order of the $r$-matrix:
  \begin{itemize}
  \item \textup{Class $G$ \textup{(}Heisenberg\textup{)}:}
    $r(z) = k/z$, pole order $1$. Regular singularities.
  \item \textup{Class $L$ \textup{(}affine KM\textup{)}:}
    $r(z) = k\Omega/z$ \textup{(}trace-form\textup{)},
    pole order $1$. Regular singularities.
  \item \textup{Class $C$ \textup{(}$\beta\gamma$\textup{)}:}
    $r$ is a constant nilpotent endomorphism
    \textup{(}$d$-log absorbs the simple OPE pole\textup{)}.
    Regular singularities.
  \item \textup{Class $M$ \textup{(}Virasoro, $\cW_N$\textup{)}:}
    $r^{\mathrm{Vir}}(z) = (c/2)/z^3 + 2T/z$,
    pole order $3$. \emph{Irregular} singularities
    of Poincar\'e rank~$2$.
  \end{itemize}
\item \textup{(The KZ connection.)}
  The flat connection on $\cF_n^{\mathrm{ord}}(\cA)$
  restricts on $U_n$ to the standard product connection
  on $\cA^{\boxtimes n}|_{U_n}$ plus the
  Knizhnik--Zamolodchikov correction:
  \begin{equation}\label{eq:kz-conn}
    \nabla_n
    = \nabla_{\cA}^{\boxtimes n}
    + \sum_{i < j}
    r_{ij}(z_i - z_j) \, dz_{ij}
  \end{equation}
  where $r_{ij}(z) \in \End(\cA \otimes \cA)((z))$
  is the classical $r$-matrix \textup{(}arity-$2$
  projection of $\Theta_\cA$\textup{)} acting on the
  $(i,j)$ tensor factors, and
  $dz_{ij} = d(z_i - z_j)$ is the connection $1$-form.

  For affine Kac--Moody at level $k$, two equivalent
  conventions coexist. In the trace-form convention,
  with $\Omega_{\mathrm{tr}}$ the Casimir built from the
  normalised trace form
  $(X,Y)_{\mathrm{tr}}
  = \mathrm{tr}(\mathrm{ad}_X \circ \mathrm{ad}_Y)
  /(2h^\vee)$:
  \begin{equation}\label{eq:kz-km}
    \nabla_n^{\mathrm{KM}}
    = \nabla^{\boxtimes n} + \sum_{i < j}
    \frac{k\,(\Omega_{\mathrm{tr}})_{ij}}{z_i - z_j}
    \, d(z_i - z_j).
  \end{equation}
  In the KZ normalisation, with $\Omega$ the inverse
  Killing form Casimir:
  \begin{equation}\label{eq:kz-km-kz}
    \nabla_n^{\mathrm{KM}}
    = \nabla^{\boxtimes n} + \sum_{i < j}
    \frac{\Omega_{ij}}{(k + h^\vee)(z_i - z_j)}
    \, d(z_i - z_j).
  \end{equation}
  The bridge identity
  $k\,(\Omega_{\mathrm{tr}})_{ij}
  = \Omega_{ij}/(k + h^\vee)$
  equates the two at generic~$k$
  (see~\cite[\S III.2]{Lorgat26I}).
  Both have a \emph{simple pole} along $\Delta_{ij}$:
  the connection is Fuchsian.
\item \textup{(Flatness.)} $\nabla_n^2 = 0$ on $U_n$,
  following from the classical Yang--Baxter equation
  for $r(z)$.
\item \textup{(Equivariant descent.)}
  For $\Einf$-chiral $\cA$ \textup{(}where $S(z) = \id$
  up to Koszul signs\textup{)}:
  $\cF_n^{\mathrm{ord}}$
  is $\Sigma_n$-equivariant, and
  $(\cF_n^{\mathrm{ord}})^{\Sigma_n}
  \simeq \cF_n^{\Sigma}$
  is the Beilinson--Drinfeld factorisation $\cD$-module
  on $X^{(n)}$.
\item \textup{(Non-equivariance for $\Eone$.)}
  For $\Eone$-chiral $\cA$:
  $\cF_n^{\mathrm{ord}}$ is \emph{not}
  $\Sigma_n$-equivariant. The braiding $S(z)$ on the
  $(i,j)$ factor introduces monodromy as the $i$-th and
  $j$-th points circle each other.
\end{enumerate}
\end{proposition}

\begin{proof}
Part~(i): holonomicity follows from
$j_*$ preserving the holonomic property
(Kashiwara's theorem~\cite{Kashiwara03}).
The external product $\cA^{\boxtimes n}$ is holonomic on
$X^n$ (being an external tensor of holonomic $\cD$-modules
on curves). The singularity type is determined by the
pole structure of the connection. For class $G$ and $L$,
the $r$-matrix $r(z)$ has a simple pole: writing the
connection locally as $d + r(z)\,dz$ with
$r(z) = A/z + O(1)$, the singularity is regular
(Fuchsian). For class $C$, the $d$-log absorption
of the OPE simple pole produces a connection with at
worst logarithmic singularity: regular.
For class $M$, the Virasoro $r$-matrix has a cubic pole
$r^{\mathrm{Vir}}(z) = (c/2)/z^3 + 2T/z$; the connection
$d + r(z)\,dz$ has a pole of order $3$ in $dz$, giving
an irregular singularity of Poincar\'e rank $3 - 1 = 2$.
Holonomic $\cD$-modules on proper varieties have
finite-dimensional de~Rham cohomology regardless of the
regularity of singularities
(Mebkhout~\cite{Mebkhout84}).

Part~(ii): the KZ connection arises from the Maurer--Cartan
element. In the convolution algebra
$\gEone = \bigoplus_n
\End((s^{-1}\bar{\cA})^{\otimes n})$,
the MC element $\Theta_\cA$ at arity $2$ gives
$r(z) = \Theta_\cA|_{\text{arity } 2}
\in \End(\cA \otimes \cA)((z))$.
The connection $1$-form is $r_{ij}(z_i - z_j)\,dz_{ij}$
where $dz_{ij} = d(z_i - z_j)$. Note: the Arnold forms
$\omega_{ij} = d\log(z_i - z_j) = dz_{ij}/(z_i - z_j)$
are generators of $H^*(\Conf_n(\CC))$ and appear as
coefficients in the bar differential $d_{\mathrm{bar}}$;
they are distinct from the connection $1$-form.

Part~(iii): the flatness condition $\nabla_n^2 = 0$
splits into three contributions:
$(\nabla^{\boxtimes n})^2 = 0$ (individual connections);
the cross-term
$[\nabla^{\boxtimes n}, \sum r_{ij}\,dz_{ij}] = 0$
(from the shift condition $[\partial_z, r(z)] = -r'(z)$
paired with $d(dz_{ij}) = 0$);
and $(\sum r_{ij}\,dz_{ij})^2 = 0$, which expands as
$\sum_{i<j<k} ([r_{12}, r_{13}] + [r_{12}, r_{23}]
+ [r_{13}, r_{23}]) \, dz_{ij} \wedge dz_{jk} = 0$,
where the coefficient vanishes by the classical
Yang--Baxter equation.

Part~(iv): for $\Einf$-chiral $\cA$, the locality axiom
gives $S(z) = \mathrm{id}$ (up to signs from the Koszul
convention), so the KZ monodromy is trivial and
$\cF_n^{\mathrm{ord}}$ descends to $X^{(n)}$.

Part~(v): for $\Eone$-chiral $\cA$, the KZ monodromy around
$\Delta_{ij}$ is $\exp(2\pi i \cdot r_{ij})$, which acts
by the quantum $R$-matrix $S_{ij}$
(Drinfeld--Kohno theorem). Since $S \neq \mathrm{id}$,
$\cF_n^{\mathrm{ord}}$ does not descend.
\end{proof}

\begin{warning}[Regularity dichotomy: gauge theory vs.\ gravity]
\label{warn:regularity}
The singularity type of the factorisation $\cD$-module
encodes a fundamental physical distinction.
\begin{itemize}
\item \textup{Regular singularities \textup{(}class $G/L/C$\textup{)}:}
  the connection is Fuchsian; the classical
  Riemann--Hilbert correspondence of Deligne applies.
  The $3$d theory is a gauge theory
  \textup{(}Chern--Simons for class $L$; abelian for
  class $G$; free matter for class $C$\textup{)}.
\item \textup{Irregular singularities
  \textup{(}class $M$\textup{)}:}
  the connection has a higher-order pole; the classical
  Riemann--Hilbert correspondence does not apply, and one
  needs the irregular Riemann--Hilbert
  correspondence of D'Agnolo--Kashiwara~\cite{DAgnoloKashiwara16}.
  The $3$d theory involves gravity
  \textup{(}holomorphic-topological gravity in the
  Beltrami parametrisation for Virasoro\textup{)}.
  The irregular singularity carries Stokes data and
  wall-crossing phenomena \textup{(}for the Virasoro
  algebra, at $c = -22/5$ the Stokes multipliers
  exhibit a wall-crossing\textup{)}.
\end{itemize}
The regularity type depends on the pole order of $r(z)$,
which is the OPE pole order minus one \textup{(}$d$-log
absorption\textup{)}: class $G/L/C$ have $r$-matrix pole
order $\leq 1$, class $M$ has pole order $\geq 2$.
\end{warning}

% ================================================================
\section{Ordered chiral homology}
\label{sec:ord-ch-hom}

The ordered factorisation $\cD$-module
$\cF^{\mathrm{ord}}(\cA)$ is a holonomic $\cD$-module with
flat connection. Its de~Rham cohomology is a
finite-dimensional graded vector space at each arity.
Assembling these across arities, with the chiral differential
from FM boundary strata, produces the ordered chiral
homology of $\cA$ over $X$.

% ----------------------------------------------------------------
\subsection{Definition}
\label{subsec:ord-ch-def}

The de~Rham cohomology of a holonomic $\cD$-module on a
proper variety is finite-dimensional (Mebkhout).
Applied to $\cF^{\mathrm{ord}}(\cA)$ on
$\Ran^{\mathrm{ord}}(X)$, this gives a graded vector space
at each arity; the cross-arity chiral differential
assembles these into a chain complex.

\begin{definition}[Ordered chiral homology via
$\Ran^{\mathrm{ord}}(X)$]
\label{def:ord-ch-hom}
Let $\cA$ be a chiral algebra on a smooth proper curve $X$,
and let $\cF^{\mathrm{ord}}(\cA)$ be the ordered
factorisation $\cD$-module on $\Ran^{\mathrm{ord}}(X)$
(Definition~\ref{def:fact-dmod-ran}).
The \emph{ordered chiral homology} is the derived
pushforward to the point:
\begin{equation}\label{eq:ord-ch-hom-ran}
  \int_X^{\mathrm{ord}} \cA
  \;:=\;
  R\Gamma_{\mathrm{dR}}\bigl(
  \Ran^{\mathrm{ord}}(X),\;
  \cF^{\mathrm{ord}}(\cA)\bigr)
  \;=\;
  \pi_* \cF^{\mathrm{ord}}(\cA)
\end{equation}
where $\pi \colon \Ran^{\mathrm{ord}}(X) \to \mathrm{pt}$
is the projection. The arity-by-arity decomposition
\begin{equation}\label{eq:ord-ch-hom-def}
  \int_X^{\mathrm{ord}} \cA
  \;=\;
  \Bigl(
  \bigoplus_{n \geq 0}
  R\Gamma_{\mathrm{dR}}\bigl(X^n,\;
  \cF_n^{\mathrm{ord}}(\cA)\bigr),
  \;\; d_{\mathrm{ch}}
  \Bigr)
\end{equation}
is the computation of this pushforward stratum by stratum,
where $\mathrm{dR}(\cF) = \cF \otimes_{\cD}
\Omega^{\bullet}_{X^n}$ is the de~Rham complex of the
$\cD$-module $\cF$ and $d_{\mathrm{ch}}$ is the
\emph{chiral differential}: the alternating sum of the
restriction maps to the codimension-$1$ boundary strata of
$\overline{\FM}_n^{\mathrm{ord}}(X)$, corresponding to
pairwise collisions of adjacent points. At each
arity, the \emph{ordered chiral homology at arity $n$} is
\begin{equation}\label{eq:total-ord-ch}
  \int_{X,n}^{\mathrm{ord}} \cA
  \;:=\;
  R\Gamma_{\mathrm{dR}}\bigl(X^n,\;
  \cF_n^{\mathrm{ord}}(\cA)\bigr).
\end{equation}
\end{definition}

% ----------------------------------------------------------------
\subsection{Symmetric descent}
\label{subsec:sym-descent}

The quotient map
$q \colon \Ran^{\mathrm{ord}}(X) \to \Ran(X)$
(Definition~\ref{def:ran-space}) induces the natural map
from ordered to symmetric chiral homology.

\begin{proposition}[Symmetric descent recovers
Beilinson--Drinfeld]
\label{prop:sym-descent}
Let $\cA$ be an $\Einf$-chiral algebra on a smooth
proper curve $X$. The quotient
$q \colon \Ran^{\mathrm{ord}}(X) \to \Ran(X)$
induces the averaging map
\begin{equation}\label{eq:sym-descent-ran}
  \int_X^{\mathrm{ord}} \cA
  = R\Gamma\bigl(\Ran^{\mathrm{ord}}(X),\,
  \cF^{\mathrm{ord}}\bigr)
  \;\xrightarrow{\;\;\av\;\;}
  R\Gamma\bigl(\Ran(X),\,
  \cF^{\Sigma}\bigr)
  = \int_X \cA,
\end{equation}
which is a quasi-isomorphism. At each arity, the
$\Sigma_n$-coinvariants of the ordered chiral homology
recover the Beilinson--Drinfeld chiral homology:
\begin{equation}\label{eq:sym-descent}
  \Bigl(\int_{X,n}^{\mathrm{ord}} \cA\Bigr)_{\Sigma_n}
  \;\simeq\;
  R\Gamma_{\mathrm{dR}}\bigl(X^{(n)},\;
  \cF_n^{\Sigma}(\cA)\bigr)
  \;=:\;
  \int_{X,n} \cA.
\end{equation}
\end{proposition}

\begin{proof}
For $\Einf$-chiral $\cA$,
$\cF_n^{\mathrm{ord}}$ is $\Sigma_n$-equivariant
(Proposition~\ref{prop:fact-dmod-props}(iv)).
The equivariant descent for $\cD$-modules works through the
quotient stack $[X^n / \Sigma_n]$: the category of
$\Sigma_n$-equivariant $\cD$-modules on $X^n$ is equivalent
to the category of $\cD$-modules on $[X^n / \Sigma_n]$.
Since $\Sigma_n$ is a finite group acting in characteristic
zero, Maschke's theorem ensures that representations of
$\Sigma_n$ are semisimple; the quotient stack
$[X^n / \Sigma_n]$ and the coarse moduli space
$X^{(n)} = X^n / \Sigma_n$ have equivalent $\cD$-module
categories. The quotient $\pi \colon X^n \to X^{(n)}$ is
ramified along the diagonals, but equivariant descent
is not affected by ramification in characteristic zero.

The de~Rham functor commutes with
$\Sigma_n$-coinvariants (which are exact for the symmetric
group acting on vector spaces over $\CC$ of
characteristic zero), giving:
$R\Gamma_{\mathrm{dR}}(X^n,
\cF_n^{\mathrm{ord}})_{\Sigma_n}
\simeq
R\Gamma_{\mathrm{dR}}(X^{(n)},
(\cF_n^{\mathrm{ord}})^{\Sigma_n})
\simeq
R\Gamma_{\mathrm{dR}}(X^{(n)},
\cF_n^{\Sigma})$.
The right side is the Beilinson--Drinfeld chiral
homology at arity $n$.

The cross-arity chiral differential $d_{\mathrm{ch}}$
is $\Sigma$-equivariant for $\Einf$-chiral algebras
because the factorisation maps (collision/OPE) are
symmetric: the OPE $Y(a,z)b$ and $Y(b,z)a$ agree up to
signs. This ensures that the total complex, not just
the fixed-arity pieces, descends to the symmetric quotient.
\end{proof}

\begin{proposition}[Lossy descent for $\Eone$-chiral algebras]
\label{prop:lossy-descent}
For an $\Eone$-chiral algebra $\cA$, the natural map
\begin{equation}\label{eq:lossy}
  \av_n \colon
  \int_{X,n}^{\mathrm{ord}} \cA
  \;\longrightarrow\;
  \Bigl(\int_{X,n}^{\mathrm{ord}} \cA\Bigr)_{\Sigma_n}
\end{equation}
is surjective but not injective for $n \geq 3$. The kernel
$\ker(\av_n)$ carries ordered data invisible to the
Beilinson--Drinfeld theory:
\begin{enumerate}[label=\textup{(\roman*)}]
\item At $n = 2$: the full spectral $r$-matrix $r(z)$
  versus the scalar $\kappa(\cA)$
  (for abelian algebras, $\kappa = \av(r(z))$;
  for non-abelian KM, $\kappa = \av(r(z)) + \dim(\fg)/2$,
  the Sugawara shift).
\item At $n = 3$: the Drinfeld associator $\Phi$ appears
  in $\ker(\av_3)$. The space of sections of $\av_3$
  compatible with the Maurer--Cartan equation is a
  $\GRT_1$-torsor.
\item At $n \geq 4$: the kernel carries higher-arity
  ordered shadow data from the $\Eone$ obstruction
  tower.
\end{enumerate}
\end{proposition}

\begin{proof}
By \cite[the non-splitting proposition]{Lorgat26I},
the short exact sequence
$0 \to \ker(\av) \to \gEone \to \gmod \to 0$
does not split as dg Lie algebras. The cross-arity
obstruction (the bar differential is
$\Sigma_{n-2}$-equivariant but not
$\Sigma_n$-equivariant) forces non-trivial elements
of $\ker(\av_n)$ at $n \geq 3$. The $\GRT_1$-torsor
structure at $n = 3$ follows from Drinfeld's theorem
on associators combined with the Etingof--Kazhdan
quantisation theorem.
\end{proof}


% ----------------------------------------------------------------
\subsection{Genus $g \geq 1$ curvature}
\label{subsec:genus-curvature}

At genus zero the bar differential squares to zero and
ordered chiral homology is an ordinary chain complex.
At genus $g \geq 1$, this fails: the fiberwise bar
differential acquires curvature proportional to the Koszul
invariant $\kappa(\cA)$ times the Hodge class
$\omega_g = c_1(\lambda)$ on $\overline{\cM}_g$.
The curvature is a \emph{base} class (on moduli),
not a \emph{fiber} form (on the curve).

\begin{remark}[Curvature at higher genus]
\label{rem:higher-genus-curvature}
At genus $g \geq 1$, the fiberwise bar differential acquires
curvature:
\begin{equation}\label{eq:fib-curvature}
  d_{\mathrm{fib}}^2 = \kappa(\cA) \cdot \omega_g
  \qquad\text{(UNIFORM-WEIGHT)}
\end{equation}
where $\kappa(\cA)$ is the Koszul invariant of $\cA$ and
$\omega_g = c_1(\lambda)$ is the first Chern class of the
Hodge line bundle on $\overline{\cM}_g$; $\omega_g$ is a
class on the \emph{moduli space} $\overline{\cM}_g$, not
a form on the curve. The scalar form of the curvature
holds in the uniform-weight sector; at multi-weight,
a cross-channel correction $\delta d^2_{\mathrm{cross}}$
must be included.

Three points require emphasis.
\begin{enumerate}[label=\textup{(\alph*)}]
\item The $\cD$-module connection at \emph{fixed} arity $n$
  remains flat: $\nabla_n^2 = 0$ on $U_n$.
  The curvature~\eqref{eq:fib-curvature} appears in the
  \emph{cross-arity} chiral differential $d_{\mathrm{ch}}$,
  which connects different arities.
\item The curvature $\kappa(\cA) \cdot \omega_g$ is a
  \emph{scalar} (it acts on all tensor factors uniformly),
  so it is $\Sigma_n$-invariant: it survives descent to
  the symmetric quotient unchanged.
\item The period-corrected differential $D^{(g)}$ of
  \cite[\S II.3]{Lorgat26I} restores $d^2 = 0$ at the
  total level by incorporating the Hodge correction.
\end{enumerate}
\end{remark}


% ================================================================
\section{Topological Hochschild versus ordered chiral
homology of $D^\times$}
\label{sec:hh-vs-ch}

Two chain complexes compute invariants of a chiral algebra
on a punctured disk: the cyclic bar complex
(topological Hochschild, from the circle
$S^1 = \partial D$) and the ordered chiral chain complex
(from the punctured disk $D^\times$ itself, with its
complex structure). For $\Einf$-chiral algebras, operad
formality identifies the two. For $\Eone$-chiral algebras,
they diverge: the ordered complex carries holomorphic data
from the spectral parameter that the cyclic bar complex
cannot see.

% ----------------------------------------------------------------
\subsection{Two computations}
\label{subsec:two-computations}

Let $\cA$ be a chiral algebra on a curve $X$. Fix a point
$p \in X$ and let $D = \Spec\, \CC[[z]]$ be the formal disk
at $p$, with $D^\times = \Spec\, \CC((z))$ the punctured
formal disk. The two chain complexes are the following.

\begin{definition}[Topological Hochschild]
\label{def:top-hh}
The \emph{topological Hochschild homology} of $\cA$ is the
factorisation homology of $\cA$, viewed as an
$\Eone$-topological algebra, over the circle $S^1 = \partial D$:
\begin{equation}\label{eq:top-hh}
  \HH^{\mathrm{top}}_*(\cA)
  := \int_{S^1} \cA.
\end{equation}
This is computed by the cyclic bar complex
\begin{equation}\label{eq:cyc-bar}
  B^{\mathrm{cyc}}_n(\cA)
  = \cA^{\otimes n}
  \Big/ \bigl(a_1 \otimes \cdots \otimes a_n
  \sim (-1)^{|a_n|(|a_1| + \cdots + |a_{n-1}|)}
  a_n \otimes a_1 \otimes \cdots \otimes a_{n-1}\bigr)
\end{equation}
with the Hochschild differential
$d_{\HH}(a_1, \ldots, a_n)
= \sum_{i=1}^{n-1} \pm (a_1, \ldots,
m_2(a_i, a_{i+1}), \ldots, a_n)
+ \pm (m_2(a_n, a_1), a_2, \ldots, a_{n-1})$
and, for algebras with non-trivial $\Ainf$-structure,
the higher terms involving $m_k$ for $k \geq 3$.

The input data: the associative product
$m_2 \colon \cA \otimes \cA \to \cA$ and
the higher $\Ainf$ operations $m_k$.
No dependence on $z$; no meromorphic forms; no Arnold
relations; no configuration-space geometry.
\end{definition}

\begin{definition}[Ordered chiral homology of $D^\times$]
\label{def:ch-hom-Dx}
The \emph{ordered chiral chain complex on $D^\times$ at
arity $n$} is:
\begin{equation}\label{eq:ch-chain-Dx}
  \cC_n^{\mathrm{ord}}(D^\times, \cA)
  := \Bigl(
  \Omega^*_{\log}\bigl(
  \overline{\FM}_n^{\mathrm{ord}}(D^\times)\bigr)
  \otimes (s^{-1}\bar{\cA})^{\otimes n},
  \;\; d_{\mathrm{dR}} + d_{\mathrm{bar}}
  + \nabla_{\mathrm{KZ}}
  \Bigr).
\end{equation}
The ingredients:
\begin{enumerate}[label=\textup{(\alph*)}]
\item $\overline{\FM}_n^{\mathrm{ord}}(D^\times)$:
  the ordered FM compactification of configurations
  of $n$ distinct points in $D^\times$. This is a smooth
  manifold with corners.
\item $\Omega^*_{\log}$:
  logarithmic differential forms, with logarithmic
  singularities along the boundary divisors of the FM
  compactification. The Arnold forms
  $\omega_{ij} = d\log(z_i - z_j)$ are among the
  generators; these are coefficients in $d_{\mathrm{bar}}$,
  not in $\nabla_{\mathrm{KZ}}$.
\item $d_{\mathrm{bar}}$:
  the bar differential encoding the full OPE. At the
  codimension-$1$ boundary stratum where $z_i \to z_j$
  ($i < j$), the bar differential extracts the residue:
  \begin{equation}\label{eq:bar-diff}
    d_{\mathrm{bar}}
    \bigl(a_1 \otimes \cdots \otimes a_n\bigr)
    \big|_{\Delta_{ij}}
    = \sum_k \pm\,
    a_1 \otimes \cdots \otimes a_{i\,(k)}\, a_j
    \otimes \cdots \otimes a_n
    \;\otimes\; \omega_{ij}^{(k)}
  \end{equation}
  where $a_{(k)} b$ are the OPE modes and
  $\omega_{ij}^{(k)}$ are the corresponding forms on the
  boundary stratum.
\item $\nabla_{\mathrm{KZ}}$:
  the KZ connection
  \begin{equation}\label{eq:kz-Dx}
    \nabla_{\mathrm{KZ}} = \sum_{i < j}
    r_{ij}(z_i - z_j)\, d(z_i - z_j),
  \end{equation}
  with connection $1$-form $r(z)\,dz$, not
  $r(z)\,d\log z$.
\end{enumerate}
The \emph{ordered chiral homology of $D^\times$} is the
derived pushforward of $\cF^{\mathrm{ord}}(\cA)$ over
$\Ran^{\mathrm{ord}}(D^\times)$; the arity-by-arity
computation gives:
\begin{equation}\label{eq:ch-hom-Dx}
  \int_{D^\times}^{\mathrm{ord}} \cA
  = R\Gamma_{\mathrm{dR}}\bigl(
  \Ran^{\mathrm{ord}}(D^\times),\;
  \cF^{\mathrm{ord}}(\cA)\bigr)
  = \bigoplus_{n \geq 0}
  H^*\bigl(\cC_n^{\mathrm{ord}}(D^\times, \cA)\bigr).
\end{equation}
The input data: the full OPE
$Y(a,z)b = \sum_k a_{(k)}b \, z^{-k-1}$, the Arnold
forms on $\Conf_n^{\mathrm{ord}}(D^\times)$, the FM
compactification, and the $r$-matrix. The output depends
on the complex structure of $D^\times$ \textup{(}the
coordinate $z$\textup{)}.
\end{definition}

\begin{remark}[Stratified perspective]
\label{rem:stratified-disk}
The disk $D$ is a stratified space with closed stratum
$\partial D = S^1$ and open stratum $D^\times$.
Topological Hochschild
$\HH^{\mathrm{top}}_*(\cA) = \int_{S^1}\cA$ is
factorisation homology over the closed stratum;
ordered chiral homology $\int_{D^\times}^{\mathrm{ord}}\cA$
is factorisation homology over the open stratum with its
complex structure.
The comparison map
$\int_{D^\times}^{\mathrm{ord}} \cA \to
\HH^{\mathrm{top}}_*(\cA)$ is induced by the boundary
retraction $D^\times \to S^1$, which \emph{forgets the
complex structure}. This map is a quasi-isomorphism for
$\Einf$-chiral algebras
\textup{(}Theorem~\textup{\ref{thm:formality-bridge}}
below\textup{)} and fails for genuinely $\Eone$-chiral
algebras
\textup{(}Theorem~\textup{\ref{thm:formality-failure}}%
\textup{)}:
the holomorphic data encoded by the ordered theory
\textup{(}the $r$-matrix, the spectral parameter, the
Arnold relations\textup{)} is invisible to the topological
theory on~$S^1$.
\end{remark}

% ----------------------------------------------------------------
\subsection{Four levels of formality}
\label{subsec:formality}

Operad formality (Kontsevich--Tamarkin) guarantees that
ordered and symmetric give the same answer for
$\Einf$-chiral algebras. But this does not guarantee
simplicity: the transferred $\Ainf$-structure may carry
non-trivial higher operations. Four formality conditions,
each strictly stronger than the last, control the
relationship between ordered chiral homology and
topological Hochschild.

\begin{definition}[Formality hierarchy]
\label{def:formality-hierarchy}
\begin{enumerate}[label=\textup{(\alph*)}]
\item \textup{(Operad formality.)} The chain operad
  $C_*(\overline{\FM}_n(\CC))$ is formal as an operad:
  it is quasi-isomorphic to its cohomology
  $H^*(\overline{\FM}_n(\CC))$ via a quasi-isomorphism
  depending on a Drinfeld associator $\Phi$
  \textup{(}Kontsevich--Tamarkin\textup{)}.
  This is \emph{unconditional}: it depends only on the
  operad, not on any algebra.
\item \textup{(Algebra formality.)} The transferred
  $\Ainf$-structure on $H^*(\cA)$ has
  $m_k^{\mathrm{tr}} = 0$ for $k \geq 3$.
  This depends on the \emph{specific algebra}:
  it holds for class $G$ \textup{(}Heisenberg\textup{)},
  and fails for class $L/C/M$.
\item \textup{(Convolution algebra formality.)}
  The modular $L_\infty$ brackets $\ell_k = 0$ for
  $k \geq 3$ in $\gmod$.
  This holds \emph{only} for class $G$
  \textup{(}\cite[prop:sc-formal-iff-class-g]{Lorgat26I}\textup{)}.
\item \textup{(SC-formality.)} The Swiss-cheese type
  operations $m_k^{\mathrm{SC}} = 0$ for $k \geq 3$.
  This is equivalent to class $G$.
\end{enumerate}
Operad formality \textup{(a)} is the weakest condition and
the only one that holds unconditionally; it guarantees
that ordered and symmetric give the \emph{same answer as
chain complexes} for $\Einf$-chiral algebras. It does
\emph{not} guarantee that the answer is simple: the
transferred $\Ainf$-structure may carry non-trivial higher
operations \textup{(}levels (b)--(d) impose progressively
stronger conditions on the specific algebra,
culminating in SC-formality, which holds only for
class~$G$\textup{)}.
\end{definition}

\begin{theorem}[Formality bridge for $\Einf$-chiral algebras]
\label{thm:formality-bridge}
Let $\cA$ be an $\Einf$-chiral algebra and $\Phi$ a
Drinfeld associator. The Kontsevich--Tamarkin operad
formality quasi-isomorphism induces a quasi-isomorphism
of chain complexes:
\begin{equation}\label{eq:formality-bridge}
  \cC_n^{\mathrm{ord}}(D^\times, \cA)
  \xrightarrow[\Phi]{\;\sim\;}
  B^{\mathrm{cyc}}_n(\cA).
\end{equation}
In particular,
$\int_{D^\times}^{\mathrm{ord}} \cA
\simeq \HH^{\mathrm{top}}_*(\cA)$
as graded vector spaces.

The transferred $\Ainf$-structure on the cohomological
model $H^*(\cA)$ may still carry non-trivial higher
operations $m_k^{\mathrm{tr}}$ for $k \geq 3$. For
class $L$ \textup{(}affine KM\textup{)}, class $C$
\textup{(}$\beta\gamma$\textup{)}, and class $M$
\textup{(}Virasoro\textup{)}: these transferred
operations are non-trivial. What operad formality
guarantees is that the ordered and symmetric complexes
compute the same cohomology.
\end{theorem}

\begin{proof}
The proof proceeds in six steps: the KT formality
theorem, the identification of the two chain complexes,
the transfer of the bar differential, the vanishing of
higher OPE modes, the transferred $\Ainf$-structure, and
the failure for genuinely $\Eone$-chiral algebras.

\medskip
\noindent\textit{Step~1: Kontsevich--Tamarkin operad formality.}
The chain operad $C_*(\overline{\FM}_n(\CC))$, with
structure maps given by the operad composition on
Fulton--MacPherson compactifications, is formal: there
exists a quasi-isomorphism of operads
\begin{equation}\label{eq:kt-formality}
  \varphi_\Phi \colon
  C_*(\overline{\FM}_n(\CC))
  \xrightarrow{\;\sim\;}
  H^*(\overline{\FM}_n(\CC))
\end{equation}
depending on the choice of a Drinfeld associator
$\Phi \in \exp(\hat{\mathfrak{t}}_3)$
\textup{(}Kontsevich~\cite{Kontsevich99},
Tamarkin~\cite{Tamarkin03}\textup{)}.
The target $H^*(\overline{\FM}_n(\CC))$ is the
Arnold algebra: it is generated by classes
$\omega_{ij} = [d\log(z_i - z_j)]$ of degree~$1$
\textup{(}$1 \leq i < j \leq n$\textup{)},
subject to the Arnold relation
$\omega_{ij}\omega_{jk}
+ \omega_{jk}\omega_{ki}
+ \omega_{ki}\omega_{ij} = 0$
for all triples $i < j < k$
\textup{(}Arnold~\cite{Arnold69}\textup{)}.
The quasi-isomorphism $\varphi_\Phi$ is
$\Sigma_n$-equivariant and respects the operad
composition maps; these properties are inherited from the
$\Sigma_n$-equivariance of the FM operad structure.

\medskip
\noindent\textit{Step~2: Identification of the two
chain complexes.}
The ordered chiral chain complex at arity~$n$
\textup{(}Definition~\textup{\ref{def:ch-hom-Dx}}\textup{)}
is the tensor product
\[
  \cC_n^{\mathrm{ord}}(D^\times, \cA)
  = \Omega^*_{\log}\bigl(
  \overline{\FM}_n^{\mathrm{ord}}(D^\times)\bigr)
  \otimes (s^{-1}\bar{\cA})^{\otimes n}
\]
with the total differential
$d_{\mathrm{dR}} + d_{\mathrm{bar}} + \nabla_{\mathrm{KZ}}$.
The de~Rham factor carries the full algebra of
logarithmic differential forms on the FM compactification.
The cyclic bar complex
$B^{\mathrm{cyc}}_n(\cA)$ uses only the cohomology
$H^*(\overline{\FM}_n^{\mathrm{ord}})$, which is the
Arnold algebra.
The quasi-isomorphism~\eqref{eq:kt-formality} induces a
chain map from the de~Rham model to the Arnold model:
\begin{equation}\label{eq:transfer-cx}
  \varphi_\Phi \otimes \id
  \colon
  \Omega^*_{\log}(\overline{\FM}_n^{\mathrm{ord}})
  \otimes (s^{-1}\bar{\cA})^{\otimes n}
  \longrightarrow
  H^*(\overline{\FM}_n^{\mathrm{ord}})
  \otimes (s^{-1}\bar{\cA})^{\otimes n}.
\end{equation}
Since $\varphi_\Phi$ is a quasi-isomorphism of
differential graded vector spaces, so is
$\varphi_\Phi \otimes \id$
\textup{(}by the K\"unneth theorem; the algebra factor
$(s^{-1}\bar{\cA})^{\otimes n}$ is degreewise
finite-dimensional\textup{)}.

\medskip
\noindent\textit{Step~3: Transfer of the bar differential.}
The bar differential $d_{\mathrm{bar}}$
\textup{(}equation~\eqref{eq:bar-diff}\textup{)}
acts at the codimension-$1$ boundary stratum
$\Delta_{ij}$ where $z_i \to z_j$, extracting the
full OPE:
\[
  d_{\mathrm{bar}}
  \bigl(a_1 \otimes \cdots \otimes a_n\bigr)
  \big|_{\Delta_{ij}}
  = \sum_{k \geq 0} \pm\,
  a_1 \otimes \cdots \otimes a_{i\,(k)}\, a_j
  \otimes \cdots \otimes a_n
  \;\otimes\; \omega_{ij}^{(k)}
\]
where $\omega_{ij}^{(k)}$ is the differential form on
the boundary stratum that pairs with OPE mode
$a_{(k)}b$.
For $\cA$ an $\Einf$-chiral algebra, the $\cD$-module
underlying the OPE is $\Sigma_n$-equivariant: the
braiding is trivial \textup{(}$S(z) = \id$\textup{)},
so $d_{\mathrm{bar}}$ commutes with the
$\Sigma_n$-action on $\overline{\FM}_n^{\mathrm{ord}}$.
Since $\varphi_\Phi$ is a $\Sigma_n$-equivariant operad
map, the transferred differential
$d_{\mathrm{bar}}^{\mathrm{tr}}
= \pi \circ d_{\mathrm{bar}} \circ \iota$,
where $\iota \colon H^* \to \Omega^*_{\log}$ is
the inclusion of cohomology representatives and
$\pi \colon \Omega^*_{\log} \to H^*$ is the projection
from the homotopy transfer data of $\varphi_\Phi$,
is well defined on the target complex.

\medskip
\noindent\textit{Step~4: Vanishing of higher OPE modes.}
The OPE $Y(a,z)b = \sum_{k \geq 0} a_{(k)}b\, z^{-k-1}$
pairs each mode with a differential form on the FM
compactification. The leading mode
$a_{(0)}b = m_2(a,b)$ pairs with the $1$-form
$\omega_{ij} = d\log(z_i - z_j)$: this is the Arnold
generator in degree~$1$, and it represents a non-trivial
cohomology class in
$H^1(\overline{\FM}_n^{\mathrm{ord}})$.

The higher modes $a_{(k)}b$ for $k \geq 1$ pair with
forms $\omega_{ij}^{(k)}$ supported on the boundary
strata of the FM compactification. These forms arise
from the Laurent expansion of the OPE along the
exceptional divisor of the blowup
$\overline{\FM}_n \to \CC^n$: the mode $a_{(k)}b$
multiplies the residue form
$z_{ij}^{-k-1}\,dz_{ij}$ by a regular form on the
remaining coordinates, where $z_{ij} = z_i - z_j$ is
the local coordinate on the boundary stratum. For
$k \geq 1$, the pole order exceeds~$1$, so the
corresponding form $\omega_{ij}^{(k)}$ involves
derivatives of the Green kernel along the exceptional
fiber.

Under the projection
$\pi \colon \Omega^*_{\log}(\overline{\FM}_n^{\mathrm{ord}})
\to H^*(\overline{\FM}_n^{\mathrm{ord}})$,
these boundary forms vanish:
\begin{equation}\label{eq:higher-mode-vanishing}
  \pi(\omega_{ij}^{(k)}) = 0
  \in H^*(\overline{\FM}_n^{\mathrm{ord}})
  \quad \text{for all } k \geq 1.
\end{equation}
The argument proceeds by degree analysis in the Arnold
algebra. The cohomology $H^*(\Conf_n(\CC))$ is generated
in degree~$1$ by the classes $[\omega_{ij}]$; every class
of degree $\geq 2$ is decomposable, i.e.\ a product of
degree-$1$ generators subject to the Arnold relation.
The forms $\omega_{ij}^{(k)}$ for $k \geq 1$ are
localised on the codimension-$1$ boundary stratum
$\Delta_{ij}$ and carry fiber degree $\geq 1$ along the
exceptional divisor, in addition to the base degree from
$\omega_{ij}$ itself. They are indecomposable in the
fiber direction: each involves the local coordinate of a
single boundary stratum and cannot be expressed as a
product of Arnold generators from distinct boundary
strata. Since the Arnold algebra has no indecomposable
classes in degree $\geq 2$, these forms must be exact
in $\Omega^*_{\log}(\overline{\FM}_n^{\mathrm{ord}})$.
Explicitly, $\omega_{ij}^{(k)}$ for $k \geq 1$ is
$d_{\mathrm{dR}}$ of a form of lower degree supported
on $\Delta_{ij}$, so it lies in the kernel of $\pi$.

The transferred bar differential on the Arnold model
therefore reduces to the contribution from $k = 0$
alone:
\[
  d_{\mathrm{bar}}^{\mathrm{tr}}
  \bigl(a_1 \otimes \cdots \otimes a_n\bigr)
  = \sum_{i < j} \pm\,
  a_1 \otimes \cdots \otimes m_2(a_i, a_j)
  \otimes \cdots \otimes a_n
  \;\otimes\; [\omega_{ij}].
\]
Here the elements $a_i \in s^{-1}\bar{\cA}$ are desuspended
augmentation-ideal elements, as in
Definition~\ref{def:ch-hom-Dx}. The resulting differential
on $H^*(\overline{\FM}_n^{\mathrm{ord}}) \otimes
(s^{-1}\bar{\cA})^{\otimes n}$ is the bar-construction
model of the Hochschild differential $d_{\HH}$: it
coincides with the differential on $B^{\mathrm{cyc}}_n(\cA)$
\textup{(}Definition~\textup{\ref{def:top-hh}}\textup{)}
under the standard identification between the bar complex
$T^c(s^{-1}\bar{\cA})$ and the Hochschild chain complex
$\cA^{\otimes n}/\mathrm{cyclic}$.

\medskip
\noindent\textit{Step~5: Transferred $\Ainf$-structure
and independence of the chain-level comparison.}
The homotopy transfer theorem
\textup{(}Kadeishvili~\cite{Kadeishvili82}\textup{)}
applied to the quasi-isomorphism $\varphi_\Phi$
produces a transferred $\Ainf$-structure
$\{m_k^{\mathrm{tr}}\}_{k \geq 2}$ on $H^*(\cA)$:
$m_2^{\mathrm{tr}} = m_2$, and for $k \geq 3$ the
operations $m_k^{\mathrm{tr}}$ are given by sums over
planar rooted trees with $k$ leaves, each internal edge
decorated by the contracting homotopy
$h \colon \Omega^*_{\log} \to \Omega^{*-1}_{\log}$
\textup{(}satisfying
$\id - \iota\pi = [d_{\mathrm{dR}}, h]$\textup{)}
and each vertex by the original multiplication.

These transferred operations encode the interaction
between the OPE pole structure and the
configuration-space geometry: they are non-trivial
precisely when algebra-level formality
\textup{(}Definition~\ref{def:formality-hierarchy}(b)\textup{)}
fails. This occurs for class $L$
\textup{(}affine Kac--Moody\textup{)}, class $C$
\textup{(}$\beta\gamma$\textup{)}, and class $M$
\textup{(}Virasoro, $\cW_N$\textup{)}.

The chain-level quasi-isomorphism
$\cC_n^{\mathrm{ord}}(D^\times, \cA)
\xrightarrow{\sim} B^{\mathrm{cyc}}_n(\cA)$
does \emph{not} depend on these higher operations.
The map $\varphi_\Phi \otimes \id$ is a
quasi-isomorphism of chain complexes by Step~2,
and the transferred bar differential equals $d_{\HH}$
by Step~4. The two complexes compute the same
cohomology, establishing~\eqref{eq:formality-bridge},
regardless of whether the $\Ainf$-structure is formal.
What $m_k^{\mathrm{tr}}$ for $k \geq 3$ controls is the
\emph{multiplicative} structure on the cohomology:
the $\Ainf$-algebra refinement of $H^*(\cA)$ as an
algebra over the operad $H^*(\overline{\FM})$.

\medskip
\noindent\textit{Step~6: Why the argument fails for
genuinely $\Eone$-chiral algebras.}
For a genuinely $\Eone$-chiral algebra, the $R$-matrix
braiding $S(z) \neq \id$ breaks the
$\Sigma_n$-equivariance of the bar differential:
$d_{\mathrm{bar}}$ no longer commutes with the
$\Sigma_n$-action, because the OPE $Y(a,z)b$ and
$Y(b,-z)a$ differ by $S(z) \neq \id$. The $\cD$-module
underlying the OPE does not descend to the symmetric
quotient $\Conf_n(\CC)/\Sigma_n$, so the transfer
through $\varphi_\Phi$ cannot reduce the ordered complex
to the cyclic bar complex. The ordered chiral complex
retains holomorphic data from the spectral-parameter
dependence of $S(z)$ that has no counterpart in
$B^{\mathrm{cyc}}_n$
\textup{(}Theorem~\textup{\ref{thm:formality-failure}}\textup{)}.
\end{proof}

\begin{theorem}[Formality failure for genuinely
$\Eone$-chiral algebras]
\label{thm:formality-failure}
Let $\cA$ be a genuinely $\Eone$-chiral algebra
\textup{(}tier~(c): the $R$-matrix $S(z)$ is not derived
from any $\Einf$-chiral OPE; $S(z) \neq \id$\textup{)}.
The chain-level formality
quasi-isomorphism~\eqref{eq:formality-bridge}
does \emph{not} exist.
\begin{enumerate}[label=\textup{(\roman*)}]
\item \textup{(Operad formality holds.)}
  The Kontsevich--Tamarkin formality of
  $C_*(\overline{\FM}_n(\CC))$ as an operad is
  unconditional: it depends only on the operad structure,
  not on the specific algebra.
\item \textup{(Equivariance fails.)}
  The $\cD$-module $\cF_n^{\mathrm{ord}}$ is not
  $\Sigma_n$-equivariant \textup{(}because
  $S(z) \neq \id$\textup{)}, so the ordered complex
  does \emph{not} descend to the symmetric quotient.
  The ordered chiral homology is strictly richer than
  topological Hochschild: the ordered complex carries
  holomorphic data from the spectral parameter
  dependence of $S(z)$ that has no counterpart in the
  cyclic bar complex.
\item \textup{(Dimension count.)}
  By Arnold's theorem, the cohomology algebra
  $H^*(\Conf_n(\CC))$ has Poincar\'e polynomial
  $\prod_{j=0}^{n-1}(1 + jt)$, with total dimension $n!$.
  The ordered chiral chain complex at arity $n$ therefore has
  \begin{equation}\label{eq:dim-count}
    \dim \cC_n^{\mathrm{ord}}(D^\times, \cA)
    = (\dim \bar{\cA})^n \cdot n!,
  \end{equation}
  where $\bar{\cA} = \ker(\varepsilon)$ is the augmentation
  ideal \textup{(}as in
  Definition~\textup{\ref{def:ch-hom-Dx}}\textup{)},
  while the cyclic bar complex has
  $\dim B^{\mathrm{cyc}}_n = (\dim \cA)^n$.
  The factor $n!$ from the Arnold algebra measures the
  additional holomorphic data carried by the ordered chiral
  complex beyond the topological Hochschild complex.
\end{enumerate}
\end{theorem}

\begin{proof}[Proof of the dimension count]
The Poincar\'e polynomial of $H^*(\Conf_n(\CC))$ is
$P_n(t) = \prod_{j=0}^{n-1}(1 + jt)$
(Arnold~\cite{Arnold69}). Setting $t = 1$ gives
$P_n(1) = \prod_{j=0}^{n-1}(1 + j) = n!$.
$P_2(1) = 1 \cdot 2 = 2$; at $n = 3$,
$P_3(1) = 1 \cdot 2 \cdot 3 = 6$; at $n = 4$,
$P_4(1) = 1 \cdot 2 \cdot 3 \cdot 4 = 24$.
The ordered chiral complex is
$\Omega^*_{\log}(\overline{\FM}_n^{\mathrm{ord}})
\otimes (s^{-1}\bar{\cA})^{\otimes n}$;
at the cohomological level, the
$\Omega^*_{\log}$ factor contributes the Arnold algebra,
whose total dimension is $n!$.
\end{proof}


% ================================================================
\section{The $\Ethree$ structure and Chern--Simons theory}
\label{sec:e3}

The symmetric bar $\BarSig(\cA)$ of a chiral algebra on a
curve $X$ is an $\Etwo$-coalgebra (from the coshuffle
coproduct on $\Conf_k(X) \cong \Conf_k(\RR^2)$). The Higher
Deligne Conjecture promotes the Hochschild cohomology of an
$\Etwo$-algebra to $\Ethree$. Applied to $\BarSig(\cA)$,
this gives the derived chiral centre
$Z^{\mathrm{der}}_{\mathrm{ch}}(\cA)$ a natural
$\Ethree$-algebra structure: $\Etwo$ from the curve,
$+1$ from the centre construction, $= \Ethree$. The
comparison with perturbative Chern--Simons theory
(Costello--Francis--Gwilliam) identifies this $\Ethree$
structure with the bulk observables of the $3$d gauge theory.

% ----------------------------------------------------------------
\subsection{The derived chiral centre}
\label{subsec:derived-centre}

The ordered and symmetric chiral Hochschild cochain complexes
differ by the presence or absence of
$\Sigma_k$-coinvariants. The derived chiral centre is the
cohomology of the symmetric version; its ordered refinement
carries the full $\Eone$ data.

\begin{definition}[Symmetric and ordered chiral
Hochschild complexes]
\label{def:ch-hochschild}
Let $\cA$ be a chiral algebra on a curve $X$.
\begin{enumerate}[label=\textup{(\roman*)}]
\item The \emph{symmetric chiral Hochschild cochain complex}
  is:
  \begin{equation}\label{eq:sym-hh}
    \HH^{*,\Sigma}_{\mathrm{ch}}(\cA, \cA)
    = \bigoplus_{k \geq 0}
    \Omega^*(\overline{\FM}_k(\CC))
    \otimes_{\Sigma_k}
    \Hom(\cA^{\otimes k}, \cA),
  \end{equation}
  with the chiral Hochschild differential.
  The derived chiral centre is
  $Z^{\mathrm{der}}_{\mathrm{ch}}(\cA)
  := \HH^{*,\Sigma}_{\mathrm{ch}}(\cA, \cA)$.
\item The \emph{ordered chiral Hochschild cochain complex}
  is:
  \begin{equation}\label{eq:ord-hh}
    \HH^{*,\mathrm{ord}}_{\mathrm{ch}}(\cA, \cA)
    = \bigoplus_{k \geq 0}
    \Omega^*(\overline{\FM}_k^{\mathrm{ord}}(\CC))
    \otimes
    \Hom(\cA^{\otimes k}, \cA),
  \end{equation}
  without $\Sigma_k$-coinvariants, with the ordered chiral
  Hochschild differential.
\end{enumerate}
\end{definition}

\begin{proposition}[The $\Ethree$ structure on the
derived centre]
\label{prop:e3-structure}
Let $\cA$ be a chiral algebra on a smooth algebraic
curve~$X$.
\begin{enumerate}[label=\textup{(\roman*)}]
\item \textup{(Source of the $\Etwo$ structure.)}
  The curve geometry provides the $\Etwo$-coalgebra
  structure on $\BarSig(\cA)$
  \textup{(}Warning~\textup{\ref{warn:e1-vs-e2-source}}%
  \textup{)}.
\item \textup{(Higher Deligne gives $\Ethree$.)}
  By the Higher Deligne Conjecture
  \textup{(}Lurie~\cite{HA}, Francis~\cite{Francis2013}%
  \textup{)}, the Hochschild cohomology of an
  $\Etwo$-algebra carries a natural $\Ethree$-algebra
  structure. Dualising the $\Etwo$-coalgebra
  $\BarSig(\cA)$ to an algebra
  \textup{(}which requires the finiteness conditions
  satisfied by standard chiral algebras,
  see~\cite[Thm~A]{Lorgat26I}\textup{)}:
  \begin{equation}\label{eq:e3-hdc}
    Z^{\mathrm{der}}_{\mathrm{ch}}(\cA)
    = \HH^*(\BarSig(\cA), \BarSig(\cA))
    \quad \text{is an $\Ethree$-algebra.}
  \end{equation}
  The dimension count: $\Etwo$ \textup{(}curve,
  $\dim_{\RR} X = 2$\textup{)} $+ 1$
  \textup{(}centre construction\textup{)} $= \Ethree$.
\item \textup{(Scope.)} For $\Einf$-chiral algebras, the
  symmetric bar $\BarSig(\cA)$ exists as a well-defined
  $\Etwo$-coalgebra and the $\Ethree$ structure on
  the derived centre follows from \textup{(i)--(ii)}.
  No Virasoro element or conformal symmetry is required.
  For genuinely $\Eone$-chiral algebras
  \textup{(}Yangians, EK quantum vertex algebras\textup{)},
  the factorisation $\cD$-module does not descend to
  $X^{(n)}$ \textup{(}Theorem~%
  \textup{\ref{thm:formality-failure}}\textup{)},
  so $\BarSig(\cA)$ does not exist. The ordered bar
  $\Barord(\cA)$ carries an $\Eone$-coalgebra structure,
  giving only $\Etwo$ on
  $\HH^*(\Barord(\cA), \Barord(\cA))$
  via the classical Deligne conjecture
  \textup{(}$\Eone + 1 = \Etwo$\textup{)}.
  The passage from $\Etwo$ to $\Ethree$ for $\Eone$-chiral
  algebras requires the Drinfeld centre construction
  \textup{(}%
  see~\cite[conj:drinfeld-center-equals-bulk]{Lorgat26III}%
  \textup{)}.
\end{enumerate}
\end{proposition}

\begin{remark}[Two distinct $\Ethree$ structures]
\label{rem:two-e3}
The paper produces two $\Ethree$-algebra structures on
the derived chiral centre of $V_k(\fg)$ at non-critical
level.
\begin{enumerate}[label=\textup{(\alph*)}]
\item \emph{Algebraic $\Ethree$}
  \textup{(}Proposition~\textup{\ref{prop:e3-structure}}%
  \textup{)}: the Hochschild cohomology of the
  $\Etwo$-coalgebra $\BarSig(V_k(\fg))$, via the Higher
  Deligne Conjecture. This requires dualisation of the
  coalgebra to an algebra, which uses the finiteness of
  $V_k(\fg)$ \textup{(}finite-type as a graded vector space
  in each conformal-weight stratum\textup{)}. No conformal
  vector is needed.
\item \emph{Topological $\Ethree$}
  \textup{(}Proposition~%
  \textup{\ref{prop:khan-zeng-topological}}\textup{)}:
  the Sugawara Virasoro element at non-critical level
  topologises the holomorphic direction of the
  $\Etwo$ from the curve geometry, producing a topological
  $\Ethree$-algebra on $\RR^3 \cong X \times \RR$.
  This requires the conformal vector and fails at
  $k = -h^\vee$.
\end{enumerate}
Structure~\textup{(a)} is formal-algebraic;
structure~\textup{(b)} is geometric-topological.
Conjecture~\textup{\ref{conj:e3-identification}} identifies
the two as families over $\hbar H^3(\fg)[[\hbar]]$.
\end{remark}

\begin{warning}[The bar complex itself is not $\Ethree$]
\label{warn:bar-not-e3}
The symmetric bar $\BarSig(\cA)$ does \emph{not} carry
an $\Ethree$-algebra structure. The obstruction is
topological: $\Ethree$ requires $\Conf_k(\RR^3)$, but a
curve provides only $\Conf_k(\RR^2)$. For
$\widehat{\mathfrak{sl}}_2$, the obstruction space
$\HH^3_{\Etwo}(\BarSig, \BarSig)$ is one-dimensional
and the obstruction class is non-zero. The
$\Ethree$-algebra is the \emph{derived centre}, not the
bar complex.
\end{warning}


% ----------------------------------------------------------------
\subsection{The CFG comparison}
\label{subsec:cfg}

Costello, Francis, and Gwilliam construct a filtered
$\Ethree$-algebra $\cA^\lambda$ from BV quantisation of
Chern--Simons theory on $\RR^3$. The derived chiral centre
of $V_k(\fg)$ and the CFG algebra both deform $C^*(\fg)$
over $\hbar H^3(\fg)[[\hbar]]$. The question is whether
these deformation families are isomorphic as
$\Ethree$-algebras.

\begin{theorem}[Costello--Francis--Gwilliam~\cite{CFG26}]
\label{thm:cfg}
Let $\fg$ be a simple finite-dimensional Lie algebra.
BV quantisation of Chern--Simons theory on $\RR^3$ with
gauge algebra $\fg$ and invariant pairing $\lambda$ yields
a filtered $\Ethree$-algebra $\cA^\lambda$ on the
Chevalley--Eilenberg cochains $C^*(\fg)$.
\begin{enumerate}[label=\textup{(\roman*)}]
\item The classical observables are
  $C^*(\fg) = \Sym(\fg^\vee[-1])$
  with the CE differential: a commutative
  \textup{(}$\Einf$\textup{)} algebra.
\item The quantum observables
  $\cA^\lambda$ are a filtered $\Ethree$-algebra deforming
  $C^*(\fg)$, with deformation space
  $\hbar H^3(\fg)[[\hbar]]$.
\item For any $\Ethree$-algebra $\cA$,
  $\mathrm{Mod}_\cA$ is $\Etwo$-monoidal \textup{(}braided
  monoidal\textup{)}. Filtered Koszul duality gives
  $\mathrm{Perf}_{C^*(\fg)}
  \simeq \mathrm{Rep}_{\mathrm{fin}}(\fg)^{\mathrm{dg}}$,
  so deforming $C^*(\fg)$ as $\Ethree$ deforms
  $\mathrm{Rep}_{\mathrm{fin}}(\fg)$ as braided monoidal:
  this recovers the Drinfeld--Jimbo quantum group.
\item \textup{(Theorem~1.4 of \cite{CFG26}.)} Canonical
  bijection between: perturbative CS quantisations, filtered
  $\Ethree$-algebras deforming $C^*(\fg)$, braided monoidal
  deformations of $\mathrm{Rep}_{\mathrm{fin}}(\fg)$, and
  quasi-triangular quasi-Hopf algebras deforming $U(\fg)$.
  All non-canonically isomorphic to
  $\hbar H^3(\fg)[[\hbar]]$.
\end{enumerate}
\end{theorem}

\begin{conjecture}[Deformation families of $\Ethree$-algebras]
\label{conj:e3-identification}
For the affine Kac--Moody vertex algebra
$V_k(\fg)$ at generic level $k$, the $\Ethree$-algebra
$Z^{\mathrm{der}}_{\mathrm{ch}}(V_k(\fg))$ of
Proposition~\textup{\ref{prop:e3-structure}} and the
CFG $\Ethree$-algebra $\cA^\lambda$ are related as
follows.
\begin{enumerate}[label=\textup{(\roman*)}]
\item \textup{(Matching deformation spaces.)}
  Both $Z^{\mathrm{der}}_{\mathrm{ch}}(V_k(\fg))$
  and $\cA^\lambda$ are parametrised by
  $\hbar H^3(\fg)[[\hbar]]$:
  for the derived chiral centre, the deformation
  parameter is the departure from critical level
  $k = -h^\vee + \hbar\lambda_1
  + \hbar^2 \lambda_2 + \cdots$;
  for the CFG algebra, it is the Chern--Simons coupling.
  This matching of \emph{deformation spaces} is a
  necessary condition for isomorphism.
\item \textup{(Base point mismatch.)}
  At the critical level $k = -h^\vee$, the derived
  chiral centre $Z^{\mathrm{der}}_{\mathrm{ch}}
  (V_{-h^\vee}(\fg))$ is related to the
  Feigin--Frenkel centre
  $\mathrm{Fun}(\mathrm{Op}_{\fg^\vee}(D))$
  \textup{(}opers on the formal disk\textup{)}, not
  directly to $C^*(\fg)$. The comparison requires a
  localisation or restriction functor, and it is at
  this base point that the identification is most
  delicate.
\item \textup{(Formal deformation families.)}
  We conjecture that the formal deformation
  \emph{families} are isomorphic:
  \begin{equation}\label{eq:e3-conj}
    \{Z^{\mathrm{der}}_{\mathrm{ch}}(V_k(\fg))
    \}_{k \in -h^\vee + \hbar\CC[[\hbar]]}
    \;\simeq\; \{\cA^\lambda\}_{\lambda
    \in \hbar H^3(\fg)[[\hbar]]}
  \end{equation}
  as families of $\Ethree$-algebras over
  $\hbar H^3(\fg)[[\hbar]]$, rather than as individual
  fibers.
\item \textup{(Sugawara constraint.)}
  The topological enhancement
  \textup{(}Theorem~\ref{thm:e3-cs}(iv)\textup{)} holds
  at generic $k$ but fails at $k = -h^\vee$, where
  the Sugawara element is undefined. The perturbative
  identification at
  $k = -h^\vee + \hbar\lambda_1 + \cdots$ holds
  order-by-order in $\hbar$, since the Sugawara element
  exists at any nonzero $\hbar$.
\item \textup{(What is proved.)} The deformation
  \emph{spaces} agree
  \textup{(}both $\hbar H^3(\fg)[[\hbar]]$\textup{)},
  which is a necessary condition but not sufficient
  for isomorphism. The comparison at generic $k$ is
  accessible through the algebraic string dictionary
  of~\cite{Lorgat26I}.
\end{enumerate}
\end{conjecture}

% ----------------------------------------------------------------
\subsection{Explicit $\Ethree$ operations for
$V_k(\mathfrak{sl}_2)$}
\label{subsec:e3-explicit}

The abstract $\Ethree$ structure of
Proposition~\ref{prop:e3-structure} acquires concrete content
for $\fg = \mathfrak{sl}_2$: the five-dimensional derived
chiral centre $\CC \oplus \mathfrak{sl}_2[-1] \oplus \CC[-2]$
carries an explicit cup product, Gerstenhaber bracket,
$\Pthree$ bracket, and BV operator, all computable from the
OPE.

\medskip
\noindent
\textbf{The underlying graded vector space.}
By Proposition~\ref{prop:e3-structure} and the chiral
Hochschild computation of
Example~\ref{ex:km}, the derived chiral centre of
$V_k(\mathfrak{sl}_2)$ at generic level
$k \neq -2$ is the graded vector space
\begin{equation}\label{eq:e3-graded-space}
  Z^{\mathrm{der}}_{\mathrm{ch}}(V_k(\mathfrak{sl}_2))
  \;=\;
  \underbrace{\CC}_{=\, \HH^0}
  \;\oplus\;
  \underbrace{\mathfrak{sl}_2[-1]}_{=\, \HH^1}
  \;\oplus\;
  \underbrace{\CC[-2]}_{=\, \HH^2},
\end{equation}
with total dimension $1 + 3 + 1 = 5$. Write $\mathbf{1}$
for the generator of $\HH^0 = \CC$ (the vacuum),
$\{e, f, h\}$ for the standard basis of
$\HH^1 = \mathfrak{sl}_2$
(outer derivations: infinitesimal level-preserving
deformations parametrised by $\fg$ itself), and $\eta$
for the generator of $\HH^2 = \CC$
(the central charge deformation cocycle). The
$\mathfrak{sl}_2$-module structure is: $\HH^0$ trivial,
$\HH^1$ adjoint, $\HH^2$ trivial.

\medskip
\noindent
\textbf{Conventions.}
The Killing form normalisation is
$(e, f) = (f, e) = 1$, $(h, h) = 2$, all other pairings zero.
Structure constants: $[e, f] = h$, $[h, e] = 2e$,
$[h, f] = -2f$. The inverse Casimir is
$\Omega = e \otimes f + f \otimes e
+ \tfrac{1}{2}\,h \otimes h$, and $h^\vee = 2$.
The KZ coupling constant is
$h_{\mathrm{KZ}} = 1/(k + h^\vee) = 1/(k + 2)$.

\begin{remark}[Notation: $h_{\mathrm{KZ}}$ vs $\hbar$]
\label{rem:hbar-convention}
Two closely related parameters appear throughout this paper.
The \emph{KZ coupling} $h_{\mathrm{KZ}} = 1/(k + h^\vee)$
is the normalisation constant in the KZ connection
$\nabla = d + h_{\mathrm{KZ}} \sum \Omega_{ij}\,dz_{ij}$;
it vanishes as $k \to \infty$ (weak-coupling limit) and
diverges at the critical level $k = -h^\vee$.
The \emph{departure from critical level}
$\hbar = k + h^\vee$ is the complementary parameter:
$\hbar = 0$ at critical level, $\hbar \to \infty$ at
weak coupling, and $h_{\mathrm{KZ}} = 1/\hbar$.
In the $\mathfrak{sl}_2$ computation below, $h^\vee = 2$,
so $h_{\mathrm{KZ}} = 1/(k+2)$ and $\hbar = k + 2$.
We write $h_{\mathrm{KZ}}$ for the $\Pthree$ bracket
normalisation (where the KZ coupling is the natural
parameter) and $\hbar$ for deformation-theoretic contexts
(where the departure from critical level controls the
formal neighbourhood).
\end{remark}

An $\Ethree$-algebra carries three levels of
operations:
\begin{itemize}
\item Level $0$ ($\Einf$): a graded-commutative product
  $\mu\colon \HH^p \otimes \HH^q \to \HH^{p+q}$.
\item Level $1$ ($\Etwo$ beyond $\Eone$): the
  Gerstenhaber bracket
  $[-,-]\colon \HH^p \otimes \HH^q \to \HH^{p+q-1}$
  of degree~$-1$, derived from the fundamental class of
  $S^1 = \Conf_2(\RR^2)/\RR^+$.
\item Level $2$ ($\Ethree$ beyond $\Etwo$): the $\Pthree$
  bracket
  $\{-,-\}\colon \HH^p \otimes \HH^q \to \HH^{p+q-2}$
  of degree~$-2$, derived from the fundamental class of
  $S^2 = \Conf_2(\RR^3)/\RR^+$.
\end{itemize}
We compute each in turn.

\begin{proposition}[Explicit $\Ethree$ operations on
$Z^{\mathrm{der}}_{\mathrm{ch}}(V_k(\mathfrak{sl}_2))$]
\label{prop:e3-explicit-sl2}
Let $k \neq -2$.
The $\Ethree$-algebra structure on
$Z^{\mathrm{der}}_{\mathrm{ch}}(V_k(\mathfrak{sl}_2))$
is determined by the following data.

\smallskip
\noindent
\textup{(I) Cup product} ($\Einf$ level).
\begin{enumerate}[label=\textup{(\alph*)}]
\item $\mathbf{1}$ is the unit:
  $\mathbf{1} \cdot f = f \cdot \mathbf{1} = f$ for all $f$.
\item
  $\mu(X, Y) = 0$ for all $X, Y \in \HH^1 = \mathfrak{sl}_2$.
\item
  $\mu(X, \eta) = 0$ for all $X \in \HH^1$, and
  $\eta \cdot \eta = 0$.
\end{enumerate}

\smallskip
\noindent
\textup{(II) Gerstenhaber bracket} ($\Etwo$ level,
degree~$-1$).
\begin{gather*}
  [\mathbf{1}, -] = 0, \qquad
  [X, Y] = 0 \quad \text{for all }
  X, Y \in \HH^1 = \mathfrak{sl}_2, \\
  [X, \eta] = 0 \quad \text{for all } X \in \HH^1,
  \qquad [\eta, \eta] = 0.
\end{gather*}

\smallskip
\noindent
\textup{(III) The $\Pthree$ bracket} ($\Ethree$ level,
degree~$-2$).
\begin{align}
  \label{eq:e3-p3-killing}
  \{X, Y\}
  &= h_{\mathrm{KZ}} \cdot (X, Y)
  = \frac{1}{k + h^\vee}\,(X, Y)
  \quad \text{for } X, Y \in \HH^1,
  \\
  \label{eq:e3-p3-derivation}
  \{X, \eta\}
  &= 0
  \quad \text{for } X \in \HH^1,
  \\
  \label{eq:e3-p3-unit}
  \{\mathbf{1}, -\}
  &= 0,
  \\
  \label{eq:e3-p3-top}
  \{\eta, \eta\}
  &= 0,
\end{align}
where $(X, Y)$ denotes the normalised Killing form
and $h_{\mathrm{KZ}} = 1/(k + h^\vee) = 1/(k + 2)$ is
the KZ coupling constant
\textup{(}Remark~\textup{\ref{rem:hbar-convention}}\textup{)}.
\end{proposition}

\begin{proof}
We establish each level by $\mathfrak{sl}_2$-equivariance,
degree counting, and the BV relation.

\textit{Cup product.}
The cup product $\mu\colon \HH^p \otimes \HH^q
\to \HH^{p+q}$ is graded-commutative:
$\mu(a, b) = (-1)^{pq}\,\mu(b, a)$.
For $X, Y \in \HH^1$, this gives
$\mu(X, Y) = -\mu(Y, X)$: the restriction to
$\HH^1 \otimes \HH^1 \to \HH^2 = \CC$ is an
$\mathfrak{sl}_2$-invariant antisymmetric bilinear form
on the adjoint representation valued in the trivial
representation.
Since
$\Lambda^2(\mathfrak{sl}_2)^{\mathfrak{sl}_2}
\cong \Hom_{\mathfrak{sl}_2}(\Lambda^2(\mathrm{ad}), \CC)
= 0$
(the exterior square of the adjoint decomposes as
$\Lambda^2(\mathrm{ad}) \cong \mathrm{ad}$,
which has no trivial summand), the cup product
$\mu(X, Y) = 0$ for all $X, Y \in \HH^1$.

The products $\mu(\HH^1, \HH^2) \subset \HH^3 = 0$ and
$\mu(\HH^2, \HH^2) \subset \HH^4 = 0$ vanish for degree
reasons.

\textit{BV operator.}
The framed $\Etwo$ structure on $\BarSig(V_k(\mathfrak{sl}_2))$
(from genus-$1$ sewing) equips the derived centre with a BV
operator $\Delta\colon \HH^n \to \HH^{n-1}$ satisfying
$\Delta^2 = 0$ and the BV relation
\begin{equation}\label{eq:bv-rel-sl2}
  [f, g]
  = \Delta(f \cdot g) - (\Delta f) \cdot g
  - (-1)^{|f|}\,f \cdot (\Delta g).
\end{equation}

The component $\Delta\colon \HH^1 \to \HH^0$ is an
$\mathfrak{sl}_2$-equivariant linear map from the adjoint
representation to the trivial representation. Since the
adjoint is irreducible and non-trivial,
$\Delta|_{\HH^1} = 0$.

The component $\Delta\colon \HH^2 \to \HH^1$ is an
$\mathfrak{sl}_2$-equivariant linear map from the trivial
to the adjoint. Again, $\Delta|_{\HH^2} = 0$.

All other components vanish trivially
($\Delta|_{\HH^0}$ lands in $\HH^{-1} = 0$;
$\Delta|_{\HH^n} = 0$ for $n \geq 3$). Therefore
$\Delta = 0$ on the entire derived centre.

\textit{Gerstenhaber bracket.}
With $\Delta = 0$ and $\mu|_{\HH^1 \otimes \HH^1} = 0$,
the BV relation~\eqref{eq:bv-rel-sl2} gives:
for $X, Y \in \HH^1$,
\[
  [X, Y]
  = \Delta(\mu(X, Y)) - (\Delta X) \cdot Y
  + X \cdot (\Delta Y)
  = \Delta(0) - 0 + 0 = 0.
\]
For $X \in \HH^1$, $\eta \in \HH^2$:
$[X, \eta] = \Delta(\mu(X, \eta)) - (\Delta X) \cdot \eta
- (-1)^1\,X \cdot (\Delta \eta)
= 0 - 0 + 0 = 0$
(since $\mu(X, \eta) \in \HH^3 = 0$,
$\Delta X = 0$, $\Delta\eta = 0$).
Similarly $[\eta, \eta] \in \HH^3 = 0$.

The vanishing of the Gerstenhaber bracket on the derived
centre is a structural consequence of the
$\mathfrak{sl}_2$-equivariance constraints: the adjoint
representation provides no equivariant maps to or from the
trivial representation. This is specific to the
\emph{derived centre} $\HH^*(\BarSig, \BarSig)$; the
Gerstenhaber bracket on the full Hochschild cochain complex
$C^*(\BarSig, \BarSig)$ is non-trivial (it carries the
commutator of derivations at the cochain level).

\textit{$\Pthree$ bracket.}
The degree-$(-2)$ bracket
$\{-,-\}\colon \HH^p \otimes \HH^q \to \HH^{p+q-2}$
satisfies the graded symmetry
\begin{equation}\label{eq:p3-symmetry}
  \{a, b\}
  = -(-1)^{(|a|-2)(|b|-2)}\,\{b, a\}
\end{equation}
of a degree-$(-2)$ Lie bracket, together with the Leibniz rule
$\{a, b \cdot c\}
= \{a, b\} \cdot c + (-1)^{(|a|-2)|b|}\,b \cdot \{a, c\}$.

\textit{Step 1: $\{\mathbf{1}, -\} = 0$.}
The Leibniz rule applied to
$\mathbf{1} = \mathbf{1} \cdot \mathbf{1}$ gives
$\{\mathbf{1}, f\} = 2\,\{\mathbf{1}, f\}$
for all $f$; hence $\{\mathbf{1}, f\} = 0$.

\textit{Step 2: $\{X, Y\} = h_{\mathrm{KZ}}\,(X, Y)$.}
For $X, Y \in \HH^1$,
$\{X, Y\} \in \HH^0 = \CC$.
The symmetry~\eqref{eq:p3-symmetry} at
$|X| = |Y| = 1$ gives
$\{X, Y\} = -(-1)^{(-1)(-1)}\,\{Y, X\}
= -(-1)^1\,\{Y, X\}
= \{Y, X\}$:
the bracket is \emph{symmetric} on $\HH^1 \otimes \HH^1$.
By $\mathfrak{sl}_2$-equivariance, $\{X, Y\}$ is an
invariant symmetric bilinear form on the adjoint
representation valued in $\CC$. The space of such forms is
one-dimensional, spanned by the Killing form $(X, Y)$.
Hence $\{X, Y\} = \alpha \cdot (X, Y)$ for a scalar
$\alpha$ depending on $k$.

The scalar $\alpha$ is determined by the $r$-matrix
residue, independently of any comparison with the CFG
construction. The $\Pthree$ bracket on the derived
chiral centre $Z^{\mathrm{der}}_{\mathrm{ch}}(V_k(\fg))
= \HH^*(\BarSig(V_k(\fg)), \BarSig(V_k(\fg)))$ is the
$\Etwo$-Hochschild cohomology operation derived from
the fundamental class of $S^2 \subset \Conf_2(\RR^3)$.
At arity~$2$, this operation is the residue of the
classical $r$-matrix paired with the fundamental class
of $S^1 \subset \Conf_2(\CC)$. In the KZ normalisation
(see~\eqref{eq:kz-sl2-arity2} and the conventions
of~\S\ref{subsec:sl2-chiral-e3}):
% AP126/AP148: KZ convention r(z) = Omega/((k+h^v)z).
% At k=0, r = Omega/(h^v z) != 0 (non-abelian; correct for KZ).
% At k=-h^v, r diverges (Sugawara singularity).
\begin{equation}\label{eq:p3-from-r-matrix-residue}
  r(z) = \frac{\Omega}{(k + h^\vee)\,z},
  \qquad
  \operatorname{Res}_{z=0}\bigl[r(z)\bigr]
  = \frac{\Omega}{k + h^\vee}.
\end{equation}
The residue acts on $X \otimes Y \in \fg \otimes \fg$
by contraction with the inverse Casimir:
$\Omega(X \otimes Y) = (X, Y)$ (the normalised Killing
form). The $\Pthree$ bracket on $\HH^1 \cong \fg$ is
therefore
\begin{equation}\label{eq:p3-bracket-normalisation}
  \{X, Y\}
  = \operatorname{Res}_{z=0}
  \bigl[r(z)(X \otimes Y)\bigr]
  = \frac{(X, Y)}{k + h^\vee}
  = h_{\mathrm{KZ}} \cdot (X, Y),
\end{equation}
giving $\alpha = h_{\mathrm{KZ}} = 1/(k + h^\vee)$.
For $\mathfrak{sl}_2$, $h^\vee = 2$ and
$\alpha = 1/(k + 2)$.

Verification: at $k \to \infty$,
$h_{\mathrm{KZ}} \to 0$ and the bracket vanishes (the
weak-coupling limit is the abelian $\Ethree$). At the
critical level $k = -h^\vee$, $h_{\mathrm{KZ}}$ diverges
(the bracket degenerates, reflecting the Feigin--Frenkel
singularity). Both limits are consistent with
the behaviour of the KZ connection.

\textit{Step 3: $\{X, \eta\} = 0$.}
The bracket $\{X, \eta\} \in \HH^1 = \mathfrak{sl}_2$
for $X \in \HH^1$, $\eta \in \HH^2$ is an
$\mathfrak{sl}_2$-equivariant map
$\mathrm{ad} \otimes \CC \to \mathrm{ad}$,
hence $\{X, \eta\} = \beta\,X$ for some scalar $\beta$.
The $\Pthree$ Jacobi identity applied to
$X, Y \in \HH^1$ and $\eta \in \HH^2$,
combined with
$\{X, Y\} = h_{\mathrm{KZ}}\,(X, Y) \in \HH^0$, gives
$\{\{X, \eta\}, \eta\} = 0$ (since
$\{\{X, \eta\}, \eta\} \in \HH^{-1} = 0$) and
$\{\eta, \{X, \eta\}\} \in \HH^{-1} = 0$,
while $\{X, \{\eta, \eta\}\} = \{X, 0\} = 0$.
The Jacobi identity is thus satisfied for all $\beta$;
the constraint comes instead from the self-consistency
of the $\Ethree$ deformation.
At the cohomological level, one applies the iterated Jacobi
identity to $X, Y \in \HH^1$ and $\eta$:
\begin{align*}
  \{\{X, Y\}, \eta\}
  &= \{X, \{Y, \eta\}\}
  - (-1)^{(|X|-2)(|Y|-2)}\,\{Y, \{X, \eta\}\} \\
  &= \{X, \beta Y\} - (-1)^1\,\{Y, \beta X\} \\
  &= \beta\,h_{\mathrm{KZ}}\,(X, Y)
  + \beta\,h_{\mathrm{KZ}}\,(Y, X) \\
  &= 2\beta\,h_{\mathrm{KZ}}\,(X, Y).
\end{align*}
The left side:
$\{\{X, Y\}, \eta\}
= \{h_{\mathrm{KZ}}(X, Y)\,\mathbf{1}, \eta\}
= h_{\mathrm{KZ}}(X, Y)\,\{\mathbf{1}, \eta\} = 0$.
Hence $2\beta\,h_{\mathrm{KZ}}\,(X, Y) = 0$ for all
$X, Y$. Since
$(e, f) = 1 \neq 0$ and $h_{\mathrm{KZ}} \neq 0$ at
generic $k$, we conclude $\beta = 0$.

\textit{Step 4: $\{\eta, \eta\} = 0$.}
The bracket $\{\eta, \eta\} \in \HH^2 = \CC$ and
the symmetry~\eqref{eq:p3-symmetry} at
$|\eta| = 2$ gives
$\{\eta, \eta\} = -(-1)^{0 \cdot 0}\,\{\eta, \eta\}
= -\{\eta, \eta\}$, hence $\{\eta, \eta\} = 0$.
\end{proof}

\begin{remark}[The Killing form as the $\Ethree$ shadow
of the Cartan $3$-form]
\label{rem:cartan-3-form}
The bracket~\eqref{eq:e3-p3-killing} has a transparent
topological origin. The $\Ethree$-deformation space is
$\hbar\,H^3(\mathfrak{sl}_2)[[\hbar]] \cong \hbar\,\CC[[\hbar]]$,
one-dimensional at each $\hbar$-order
(Theorem~\ref{thm:cfg}(iv)), where $\hbar = k + h^\vee$
is the departure from critical level. The generator of
$H^3(\mathfrak{sl}_2)$ is the Cartan $3$-form
\begin{equation}\label{eq:cartan-3form}
  \omega_3(X, Y, Z)
  = (X, [Y, Z]),
\end{equation}
the unique (up to scalar) $\mathfrak{sl}_2$-invariant element
of $\Lambda^3(\mathfrak{sl}_2^*)$. The degree-$(-2)$ bracket
$\{X, Y\}$ is the $\Ethree$-operation induced by this
deformation class scaled by
$h_{\mathrm{KZ}} = 1/\hbar$:
\[
  \{X, Y\}
  = h_{\mathrm{KZ}} \cdot (X, Y).
\]
The mechanism is \emph{not} literal contraction of the
antisymmetric $3$-form $\omega_3$: the $\Pthree$ bracket
is symmetric on $\HH^1$ (by the graded
symmetry~\eqref{eq:p3-symmetry} at $|X| = |Y| = 1$),
while any contraction of $\omega_3$ in two slots is
antisymmetric. The connection is indirect:
the one-dimensional deformation space
$H^3(\mathfrak{sl}_2)$, generated by $\omega_3$,
parametrises $\Ethree$-deformations of $C^*(\fg)$
(Theorem~\ref{thm:cfg}(iv));
on cohomology, $\mathfrak{sl}_2$-equivariance forces
the induced $\Pthree$ bracket on
$\HH^1 \cong \mathfrak{sl}_2$ to be proportional to
the Killing form $(X, Y)$, the unique invariant
symmetric bilinear form. The coefficient
$h_{\mathrm{KZ}} = 1/(k + h^\vee)$ is fixed by the
$r$-matrix residue
(equation~\eqref{eq:p3-from-r-matrix-residue}), not by
comparison with CFG.
The $\Pthree$ bracket on $Z^{\mathrm{der}}_{\mathrm{ch}}$
is the binary shadow of the ternary Cartan deformation class,
with the Killing form appearing through equivariance rather
than through direct contraction.
\end{remark}

% ----------------------------------------------------------------
\subsection{Quantum corrections from the vertex $R$-matrix:
$\Ethree$ operations for the EK quantum VOA}
\label{subsec:e3-ek-quantum}

The classical $V_k(\mathfrak{sl}_2)$ has trivial braiding
$S(z) = \id$; the Etingof--Kazhdan quantum VOA
$V_{\mathrm{EK}}$ of Example~\ref{ex:ek-qvoa} has
$S(z) = \id + h_{\mathrm{KZ}}\,\Omega/z + O(z^{-2})$. Does
the non-trivial $R$-matrix change the $\Ethree$-algebra
structure on the derived chiral centre? On cohomology, no:
$\mathfrak{sl}_2$-equivariance forces the $\Pthree$ bracket
to be proportional to the Killing form with the same
coefficient $h_{\mathrm{KZ}} = 1/(k+2)$. On cochains, yes:
the quantum corrections encode the braided monoidal structure
of $\mathrm{Rep}(Y_\hbar(\mathfrak{sl}_2))$.

\medskip
\noindent
\textbf{The underlying graded vector space.}
The EK quantum VOA is a flat formal deformation of
$V_k(\mathfrak{sl}_2)$ over $\CC[[h_{\mathrm{KZ}}]]$, with
$h_{\mathrm{KZ}} = 1/(k+2)$.
Since the chiral Hochschild cohomology is computed by
a bounded complex of finite-dimensional
$\mathfrak{sl}_2$-modules, the deformation does not
change the underlying graded vector space:
\begin{equation}\label{eq:ek-e3-graded}
  Z^{\mathrm{der}}_{\mathrm{ch}}(V_{\mathrm{EK}})
  \;=\;
  \CC[[h_{\mathrm{KZ}}]]
  \;\oplus\;
  \mathfrak{sl}_2[[h_{\mathrm{KZ}}]][-1]
  \;\oplus\;
  \CC[[h_{\mathrm{KZ}}]][-2],
\end{equation}
with total rank $5$ over $\CC[[h_{\mathrm{KZ}}]]$.
The generators
$\mathbf{1} \in \HH^0$, $\{e, f, h\} \in \HH^1$, and
$\eta \in \HH^2$ are defined as in
Proposition~\ref{prop:e3-explicit-sl2}, now taken
over the base ring $\CC[[h_{\mathrm{KZ}}]]$.

\medskip
\noindent
\textbf{Equivariance under the quantum $R$-matrix.}
The Yang $R$-matrix $R(u) = u\,I + \hbar\,P$ on
$V \otimes V$ is $\mathfrak{sl}_2$-equivariant:
$[R(u),\, \rho(X) \otimes I + I \otimes \rho(X)] = 0$
for all $X \in \mathfrak{sl}_2$.
The full vertex $R$-matrix
$S(z)$ on $V_{\mathrm{EK}}^{\otimes 2}$
(equation~\eqref{eq:ek-vertex-rmatrix}) is likewise
$\mathfrak{sl}_2$-equivariant, since it is constructed from
$\mathcal{R}(z)$, the spectral universal $R$-matrix of
the Yangian $Y_\hbar(\mathfrak{sl}_2)$, which
commutes with the $\mathfrak{sl}_2$-action by the
defining property of the Drinfeld coproduct.

This equivariance is the key structural constraint: all
three levels of the $\Ethree$ operations on
$Z^{\mathrm{der}}_{\mathrm{ch}}(V_{\mathrm{EK}})$ are
$\mathfrak{sl}_2$-equivariant maps between representations.
The equivariance arguments of
Proposition~\ref{prop:e3-explicit-sl2} therefore apply
verbatim to the quantum case.

\begin{proposition}[{$\Ethree$ operations on
$Z^{\mathrm{der}}_{\mathrm{ch}}(V_{\mathrm{EK}})$}]
\label{prop:e3-ek-quantum}
Let $V_{\mathrm{EK}}$ be the Etingof--Kazhdan quantum
vertex algebra for $\mathfrak{sl}_2$ at KZ coupling
$h_{\mathrm{KZ}} = 1/(k+2) \neq 0$
\textup{(}Example~\textup{\ref{ex:ek-qvoa}}\textup{)}.
The $\Ethree$-algebra structure on the derived chiral
centre $Z^{\mathrm{der}}_{\mathrm{ch}}(V_{\mathrm{EK}})$
coincides with the classical $\Ethree$ structure of
Proposition~\textup{\ref{prop:e3-explicit-sl2}}: the
cup product and Gerstenhaber bracket vanish, and the
$\Pthree$ bracket is
\begin{equation}\label{eq:ek-p3-bracket}
  \{X, Y\}_q \;=\; h_{\mathrm{KZ}} \cdot (X, Y)
  \qquad
  \textup{for } X, Y \in \HH^1 = \mathfrak{sl}_2,
\end{equation}
with all other brackets zero. In particular, the
vertex $R$-matrix $S(z)$ does not renormalise the
$\Pthree$ bracket on cohomology.
\end{proposition}

\begin{proof}
The argument has four ingredients: equivariance,
one-dimensionality of the obstruction space,
exactness of the rational classical $r$-matrix, and a
$\GRT_1$ consistency check.

\textit{Step 1: cup product and BV operator.}
The cup product $\mu \colon \HH^p \otimes \HH^q
\to \HH^{p+q}$ is a graded-commutative product on
$\mathfrak{sl}_2$-modules over $\CC[[h_{\mathrm{KZ}}]]$.
For $X, Y \in \HH^1
= \mathfrak{sl}_2[[h_{\mathrm{KZ}}]]$,
the product $\mu(X, Y) \in \HH^2
= \CC[[h_{\mathrm{KZ}}]]$
is an $\mathfrak{sl}_2$-equivariant antisymmetric
bilinear form on the adjoint representation:
$\Lambda^2(\mathrm{ad})^{\mathfrak{sl}_2} = 0$,
so $\mu(X, Y) = 0$ over $\CC[[h_{\mathrm{KZ}}]]$.

The BV operator $\Delta \colon \HH^n \to \HH^{n-1}$
is $\mathfrak{sl}_2$-equivariant. The components
$\Delta|_{\HH^1} \colon \mathrm{ad} \to \CC$ and
$\Delta|_{\HH^2} \colon \CC \to \mathrm{ad}$ vanish
by Schur's lemma. Hence $\Delta = 0$ on the derived
centre, and the BV relation gives
$[X, Y] = 0$ for all $X, Y \in \HH^1$, just as in the
classical case.

\textit{Step 2: the $\Pthree$ bracket is proportional
to the Killing form.}
For $X, Y \in \HH^1$, the bracket
$\{X, Y\}_q \in \HH^0 = \CC[[h_{\mathrm{KZ}}]]$ is an
$\mathfrak{sl}_2$-equivariant symmetric bilinear
form on the adjoint representation:
$\Sym^2(\mathrm{ad})^{\mathfrak{sl}_2} \cong \CC$.
Hence $\{X, Y\}_q = f(h_{\mathrm{KZ}}) \cdot (X, Y)$
for some
$f(h_{\mathrm{KZ}}) \in \CC[[h_{\mathrm{KZ}}]]$.

The remaining brackets ($\{X, \eta\}_q$,
$\{\eta, \eta\}_q$, $\{\mathbf{1}, -\}_q$)
vanish by the same equivariance and Jacobi arguments
as in Proposition~\ref{prop:e3-explicit-sl2}
(the argument for $\beta = 0$ uses
$h_{\mathrm{KZ}} \neq 0$, $(e,f) = 1 \neq 0$, and
$\{\mathbf{1}, -\} = 0$; all remain valid over
$\CC[[h_{\mathrm{KZ}}]]$).

\textit{Step 3: the scalar
$f(h_{\mathrm{KZ}}) = h_{\mathrm{KZ}}$
\textup{(}from exactness of the rational
$r$-matrix\textup{)}.}
The $\Pthree$ bracket on cohomology is computed by
the $d$-log residue of the classical $r$-matrix
(Construction~\ref{constr:chiral-p3-bracket},
equation~\eqref{eq:chiral-p3-total}): for
$X, Y \in \HH^1 \cong \fg$, the bracket
$\{X, Y\}_q$ is determined by the coefficient of
the simple pole in $r(z)$.

For the EK quantum vertex algebra, the spectral
universal $R$-matrix $\mathcal{R}(z)$ of
$Y_{h_{\mathrm{KZ}}}(\mathfrak{sl}_2)$
evaluates on the fundamental representation to the
Yang $R$-matrix
% AP126: level prefix h_{KZ} present; at h_{KZ}=0, R(z)=zI, r=0. Verified.
\begin{equation}\label{eq:yang-rmatrix-exact}
  R(z) = z\,I + h_{\mathrm{KZ}}\,P
  \;\in\; \End(V \otimes V)[z],
\end{equation}
which is \emph{polynomial of degree one} in the
spectral parameter $z$
(equation~\eqref{eq:ek-rmatrix-fund}).
The classical $r$-matrix extracted from $R(z)$ is
\begin{equation}\label{eq:rational-rmatrix-exact}
  % AP126: r(z) = h_{KZ}*Omega/z; at h_{KZ}=0, r=0. Verified.
  r(z) = \frac{h_{\mathrm{KZ}}\,\Omega}{z},
\end{equation}
with $\Omega = \sum_a t^a \otimes t_a$ the Casimir
element dual to the Killing form. This is
\emph{exact}: the rational (Yang) $r$-matrix has a
simple pole and no higher-order terms in $z^{-1}$.
The higher-order corrections appearing in the vertex
$R$-matrix expansion~\eqref{eq:ek-s-expansion}
(the $z^{-2}$, $z^{-3}$, \ldots\ terms) contribute
to cochain-level operations but do not modify the
classical $r$-matrix, which is determined by the
$O(h_{\mathrm{KZ}})$ coefficient of $S(z)$ at the
$z^{-1}$ term.

The $\Pthree$ bracket on cohomology is the residue
\begin{equation}\label{eq:p3-residue-exact}
  \{X, Y\}_q
  = \Res_{z=0}\,r(z)(X \otimes Y)\,dz
  = \Res_{z=0}\,
  \frac{h_{\mathrm{KZ}}\,\Omega(X \otimes Y)}{z}\,dz
  = h_{\mathrm{KZ}}\,(X, Y),
\end{equation}
where the last equality uses
$\Omega(X \otimes Y) = (X, Y)$ (the Killing form
pairing through the Casimir). Since $r(z)$ has
\emph{only} a simple pole (no $z^{-k}$ terms for
$k \geq 2$), this residue receives no corrections
at any order in $h_{\mathrm{KZ}}$. Hence
$f(h_{\mathrm{KZ}}) = h_{\mathrm{KZ}}$.

The essential point is the distinction between
rational, trigonometric, and elliptic $R$-matrices.
The Yang $R$-matrix
$R(z) = zI + h_{\mathrm{KZ}} P$
is the unique rational solution of the Yang--Baxter
equation (up to gauge equivalence); it is polynomial
of degree one in $z$. A trigonometric $R$-matrix
(e.g., the XXZ solution
$R(z) \sim \sinh(z + h_{\mathrm{KZ}} P)/\sinh(z)$,
associated to $U_q(\widehat{\mathfrak{sl}}_2)$
rather than $Y_{h_{\mathrm{KZ}}}(\mathfrak{sl}_2)$)
would have poles at $z \in 2\pi i\ZZ$, producing a
Laurent expansion with infinitely many negative
powers of $z$ and potentially renormalising the
$\Pthree$ bracket. For the EK quantum vertex
algebra, which quantises affine Kac--Moody on the
formal disk, the rational solution governs, and the
$\Pthree$ bracket is protected by the polynomial
structure of $R(z)$.

\textit{Step 4: $\GRT_1$ consistency check.}
The preceding argument determines
$f(h_{\mathrm{KZ}}) = h_{\mathrm{KZ}}$
from the explicit form of the Yang $R$-matrix; the
$\GRT_1$-action provides an independent consistency
check. The group $\GRT_1$ acts on the
$\Ethree$-operad maps by conjugation via the
Drinfeld associator $\Phi_{\mathrm{KZ}}$. On the
one-dimensional deformation space
$h_{\mathrm{KZ}}\,\CC[[h_{\mathrm{KZ}}]]$ (since
$H^3(\mathfrak{sl}_2) \cong \CC$), this action is
by rescaling. At the cohomological level, the
conjugation shifts the $\Pthree$ bracket by
Hochschild coboundaries, which vanish on
$\HH^* = H^*(\BarSig, \BarSig)$. The $\Pthree$
bracket is therefore $\GRT_1$-invariant on
cohomology, consistent with the scalar
$f(h_{\mathrm{KZ}}) = h_{\mathrm{KZ}}$ being
independent of the choice of associator.
\end{proof}

\begin{remark}[Cochain-level corrections and bulk rigidity]
\label{rem:ek-cochain-corrections}
\label{rem:ek-bulk-rigidity}
At the \emph{cochain} level, the $\Pthree$ bracket on
$C^1(\BarSig, \BarSig)$ acquires quantum corrections
$\{D_X, D_Y\}^{\mathrm{ch}}
= h_{\mathrm{KZ}}(X,Y)\,\id
+ h_{\mathrm{KZ}}^2\,\Gamma_2 + \cdots$,
where each $\Gamma_n$ is $d_{\mathrm{Hoch}}$-exact. These
corrections encode the braided monoidal structure on
$\mathrm{Rep}(Y_\hbar(\mathfrak{sl}_2))$
(here $\hbar = k + h^\vee$ is the departure from critical
level); they are invisible
on cohomology but recover the Drinfeld--Jimbo quantum group
at the cochain level. The cohomological $\Ethree$ operations
(cup product, Gerstenhaber bracket, $\Pthree$ bracket, BV
operator) are identical for the classical $V_k(\mathfrak{sl}_2)$
and the EK quantum $V_{\mathrm{EK}}$. The sole non-trivial
cohomological $\Ethree$ invariant is
$h_{\mathrm{KZ}} = 1/(k+2)$: the perturbative bulk
observables of Chern--Simons are insensitive to the choice
of associator.
\end{remark}

\begin{remark}[Why the quantum $R$-matrix does not
renormalise the $\Pthree$ bracket]
\label{rem:ek-why-no-renorm}
Four mechanisms, in decreasing order of directness:
\begin{enumerate}[label=\textup{(\roman*)}]
\item
  \textup{(Rational $r$-matrix exactness.)}
  The Yang $R$-matrix $R(z) = zI + h_{\mathrm{KZ}} P$ is polynomial
  degree one in $z$, so the classical $r$-matrix
  $r(z) = h_{\mathrm{KZ}}\,\Omega/z$ has a simple pole and no
  higher-order $z$-corrections. The $\Pthree$ bracket is the
  residue of $r(z)$, hence $h_{\mathrm{KZ}}\,(X,Y)$ exactly.
\item
  \textup{(Equivariance.)}
  $\mathfrak{sl}_2$-equivariance constrains
  $\{X,Y\}_q$ to be proportional to the Killing form $(X,Y)$,
  reducing the problem to a single scalar $f(h_{\mathrm{KZ}})$.
\item
  \textup{($\GRT_1$-invariance.)}
  $\dim H^3(\mathfrak{sl}_2) = 1$ makes the deformation
  space a formal line; the Drinfeld associator shifts the
  $\Pthree$ bracket by Hochschild coboundaries, which
  vanish on cohomology.
\item
  \textup{(Chern--Simons perturbative finiteness.)}
  The $\Pthree$ bracket is the tree-level cubic vertex;
  pure Chern--Simons has no higher-loop corrections.
\end{enumerate}
For simple $\fg$ with $\dim H^3 = 1$, mechanisms
(i)--(iii) all apply and the rigidity persists.
For non-simple $\fg$ with
$\dim H^3 > 1$, mechanism (ii) still constrains the bracket
to the space of invariant symmetric forms, but the
higher-dimensional deformation space could allow
parameter mixing under the $\GRT_1$-action; mechanism (i)
remains valid whenever the governing $R$-matrix is rational.
\end{remark}


% ----------------------------------------------------------------
% The chiral E3-algebra: holomorphic deformation
\label{sec:chiral-e3-algebra}

Both the derived chiral centre
$Z^{\mathrm{der}}_{\mathrm{ch}}(V_k(\fg))$ and the CFG
$\Ethree$-algebra $\cA^\lambda$ deform $C^*(\fg)$, but through
different mechanisms. The chiral bar complex produces a
filtered $\Ethree$-chiral algebra
$\CE^{\mathrm{ch}}_\hbar(\fg)$: a factorisation $\cD$-module on
$\Ran(X)$ whose global sections on the formal disk recover
the CFG $\Ethree$-algebra $\cA^\lambda$.


% ----------------------------------------------------------------
\subsection{The chiral $\Pthree$ bracket from the PVA
$\lambda$-bracket}
\label{subsec:chiral-p3-bracket}

The $\Ethree$ structure on the derived chiral centre is
abstract; it comes from Higher Deligne applied to the
$\Etwo$-coalgebra $\BarSig(\cA)$. To make it explicit, one
needs a concrete bracket. The Poisson vertex algebra (PVA)
$\lambda$-bracket on $V_k(\fg)$ provides exactly this: a
degree-$(-2)$ Lie bracket on the chiral Chevalley--Eilenberg
complex, yielding a chiral $\Pthree$ structure in the
category of $\cD$-modules on $X$.

\begin{definition}[The chiral Chevalley--Eilenberg complex]
\label{def:chiral-ce-complex}
Let $\fg$ be a simple Lie algebra with Killing form
$(\cdot, \cdot)$ and dual Coxeter number~$h^\vee$.
Write $\hbar = k + h^\vee$ for the departure from critical
level. The \emph{chiral Chevalley--Eilenberg complex} is
the curved chiral coalgebra
\begin{equation}\label{eq:chiral-ce-complex}
  \CE^{\mathrm{ch}}(\fg_\hbar)
  \;:=\;
  \bigl(\Sym^c(\fg^*[-1]) \otimes \omega_X,\;
  d_{\CE},\; m_0 = \tfrac{\hbar}{2h^\vee}\,\kappa\bigr),
\end{equation}
where $\omega_X$ denotes the dualising sheaf of~$X$,
$d_{\CE}$ is the chiral Chevalley--Eilenberg differential
encoding the Lie bracket of~$\fg$, and the curvature
$m_0$ is the image of the level under the bar construction
(the notation $\kappa$ here denotes the Koszul invariant,
not the Killing form).
At $\hbar = 0$ (critical level $k = -h^\vee$), the
curvature vanishes and
$\CE^{\mathrm{ch}}(\fg_0) = (\Sym^c(\fg^*[-1]) \otimes
\omega_X,\, d_{\CE})$
is an uncurved commutative ($\Einf$) factorisation
coalgebra on~$X$.
\end{definition}

\begin{construction}[The chiral $\Pthree$ bracket]
\label{constr:chiral-p3-bracket}
The PVA $\lambda$-bracket on $V_k(\fg)$ is
(see \cite{Lorgat26I}, \S5.11):
\begin{equation}\label{eq:pva-bracket-recall}
  \{J^a \mathbin{{}_\lambda} J^b\}
  \;=\;
  f^{ab}{}_c\, J^c \;+\; k\,(a,b)\,\lambda,
\end{equation}
where $f^{ab}{}_c$ are the structure constants of $\fg$ and
$(a,b)$ is the Killing form.

At the level of the chiral Chevalley--Eilenberg complex, the
$\lambda$-bracket induces a degree-$(-2)$ bracket
\begin{equation}\label{eq:chiral-p3-def}
  \{-,-\}^{\mathrm{ch}}
  \colon
  \CE^{\mathrm{ch}}(\fg_\hbar)
  \otimes
  \CE^{\mathrm{ch}}(\fg_\hbar)
  \;\longrightarrow\;
  \CE^{\mathrm{ch}}(\fg_\hbar),
\end{equation}
defined on generators $\phi_a = (J^a)^* \in \fg^*[-1]$ by
the two components of the PVA bracket:
\begin{enumerate}[label=\textup{(\roman*)}]
\item \textup{(Lie bracket term.)}
  The structure constant term $f^{ab}{}_c\,J^c$ contributes
  the Chevalley--Eilenberg differential $d_{\CE}$ at the
  algebra level. On the dual (coalgebra) side, it induces a
  coderivation.
  This is the classical $\Pthree$ bracket of the
  uncurved complex $\CE^{\mathrm{ch}}(\fg_0)$: the degree-$(-2)$
  Lie bracket
  \begin{equation}\label{eq:p3-lie-component}
    \{\phi_a, \phi_b\}^{\mathrm{ch}}_{\mathrm{Lie}}
    \;=\;
    f^{ab}{}_c\,\phi_c.
  \end{equation}
  This bracket is the coinvariant shadow of the ordered bar
  $R$-matrix at the critical level.
\item \textup{(Cocycle term.)}
  The level-dependent term $k\,(a,b)\,\lambda$ contributes a
  degree-$(-2)$ symmetric pairing valued in $\cO_X$:
  \begin{equation}\label{eq:p3-cocycle-component}
    \{\phi_a, \phi_b\}^{\mathrm{ch}}_{\mathrm{cocycle}}
    \;=\;
    k\,(a,b) \cdot \partial,
  \end{equation}
  where $\partial$ denotes the $\cD$-module action
  (the spectral parameter~$\lambda$ of the PVA bracket
  becomes the derivation $\partial_z$ in the $\cD$-module
  structure). This term is the cocycle that controls
  the deformation away from the commutative
  ($\Einf$) factorisation structure.
\end{enumerate}

The total chiral $\Pthree$ bracket is
\begin{equation}\label{eq:chiral-p3-total}
  \{\phi_a, \phi_b\}^{\mathrm{ch}}
  \;=\;
  f^{ab}{}_c\,\phi_c
  \;+\;
  k\,(a,b)\,\partial.
\end{equation}
\end{construction}

\begin{proposition}[The chiral $\Pthree$ structure]
\label{prop:chiral-p3-structure}
The bracket~\eqref{eq:chiral-p3-total} equips the chiral
Chevalley--Eilenberg complex $\CE^{\mathrm{ch}}(\fg_\hbar)$
with the structure of a $\Pthree$-algebra object in
$\cD\textrm{-}\mathrm{mod}(X)$. Concretely:
\begin{enumerate}[label=\textup{(\roman*)}]
\item The bracket is a degree-$(-2)$ Lie bracket satisfying
  graded antisymmetry and the Jacobi identity.
  The Jacobi identity follows from the Jacobi identity of the
  PVA $\lambda$-bracket, which in turn encodes the Jacobi
  identity of the Lie algebra $\fg$ together with the
  $\fg$-invariance of the Killing form.
\item The bracket is compatible with the commutative product
  on $\Sym^c(\fg^*[-1])$ via the Leibniz rule:
  \[
    \{a, b \cdot c\}^{\mathrm{ch}}
    = \{a, b\}^{\mathrm{ch}} \cdot c
    + (-1)^{(|a|-2)|b|}\, b \cdot \{a, c\}^{\mathrm{ch}}.
  \]
\item The bracket is $\cD$-linear in the sense that it
  respects the flat connection on
  $\CE^{\mathrm{ch}}(\fg_\hbar)$ inherited from the
  $\cD$-module structure: the chiral $\Pthree$ bracket
  is a morphism of $\cD$-modules, not merely of
  $\cO_X$-modules.
\end{enumerate}
\end{proposition}

\begin{proof}
The PVA $\lambda$-bracket on $V_k(\fg)$ satisfies the PVA
axioms (sesquilinearity, Leibniz, skew-symmetry, Jacobi);
these are the $\lambda$-bracket translations of the classical
master equation for the Poisson structure.
The identification of the $\lambda$-bracket with a
degree-$(-2)$ Lie bracket on the shifted dual follows the
general correspondence between PVA brackets and shifted
Poisson structures (Beilinson--Drinfeld~\cite{BD04},
\S3.8): a PVA $\lambda$-bracket of conformal weight $1$
(where $\lambda$ has weight $1$) corresponds to a
$\Pthree$ bracket (degree $-2$ = $-\dim_\CC X - 1$, where
$\dim_\CC X = 1$) on the $\cD$-module of chiral cochains.

The $\cD$-linearity follows because the PVA
$\lambda$-bracket is a $\cD$-module morphism by construction:
the derivation $\partial$ acts via $\lambda \mapsto \lambda$
in the $\lambda$-bracket formalism, and this action is
compatible with the Leibniz rule by the sesquilinearity
axiom of the PVA.
\end{proof}

% ----------------------------------------------------------------
\subsection{The filtered $\Ethree$-chiral algebra}
\label{subsec:filtered-e3-chiral}

The chiral $\Pthree$ bracket and the chiral
Chevalley--Eilenberg complex assemble into a single object:
a filtered $\Ethree$-chiral algebra
$\CE^{\mathrm{ch}}_\hbar(\fg)$, a factorisation $\cD$-module
on $\Ran(X)$ whose global sections on the formal disk
recover the CFG $\Ethree$-algebra.

\begin{definition}[The filtered $\Ethree$-chiral algebra]
\label{def:filtered-e3-chiral}
Let $\hbar = k + h^\vee$. The \emph{chiral $\Ethree$-algebra}
\begin{equation}\label{eq:chiral-e3-algebra}
  \CE^{\mathrm{ch}}_\hbar(\fg)
  \;:=\;
  \bigl(\CE^{\mathrm{ch}}(\fg_\hbar),\;
  \{-,-\}^{\mathrm{ch}},\;
  \cD\textrm{-module structure}\bigr)
\end{equation}
is the chiral Chevalley--Eilenberg complex
(Definition~\ref{def:chiral-ce-complex}) equipped with
the chiral $\Pthree$ bracket
(Construction~\ref{constr:chiral-p3-bracket}) and the
$\cD$-module structure from the factorisation algebra
on~$X$. It carries an $\hbar$-adic filtration
$F^\bullet \CE^{\mathrm{ch}}_\hbar(\fg)$ defined by:
\begin{equation}\label{eq:hbar-filtration}
  F^p \CE^{\mathrm{ch}}_\hbar(\fg)
  \;=\;
  \hbar^p \cdot \CE^{\mathrm{ch}}_\hbar(\fg).
\end{equation}
\end{definition}

\begin{theorem}[Structure of the chiral $\Ethree$-algebra]
\label{thm:chiral-e3-structure}
Let $\fg$ be a simple Lie algebra and $\hbar = k + h^\vee$.
The chiral $\Ethree$-algebra
$\CE^{\mathrm{ch}}_\hbar(\fg)$ satisfies:
\begin{enumerate}[label=\textup{(\roman*)}]
\item \textup{(Associated graded.)}
  The associated graded with respect to the $\hbar$-adic
  filtration is the uncurved chiral Chevalley--Eilenberg
  complex at the critical level:
  \begin{equation}\label{eq:assoc-graded}
    \mathrm{gr}_F\,\CE^{\mathrm{ch}}_\hbar(\fg)
    \;\simeq\;
    \CE^{\mathrm{ch}}(\fg_0)
    \;=\;
    \bigl(\Sym^c(\fg^*[-1]) \otimes \omega_X,\; d_{\CE}\bigr),
  \end{equation}
  which is an $\Einf$-chiral algebra \textup{(}commutative
  factorisation $\cD$-module on $\Ran(X)$\textup{)}.
\item \textup{(The $\Ethree$ structure.)}
  The chiral $\Pthree$ bracket
  \textup{(}Proposition~\textup{\ref{prop:chiral-p3-structure}}%
  \textup{)} deforms the $\Einf$ structure to
  a filtered $\Ethree$-chiral algebra:
  $\Etwo$ from the curve geometry
  \textup{(}Warning~\textup{\ref{warn:e1-vs-e2-source}}%
  \textup{)} $+ 1$ from the chiral $\Pthree$ bracket
  $= \Ethree$.
\item \textup{(Curvature and the deformation parameter.)}
  The curvature
  $m_0 = \tfrac{\hbar}{2h^\vee}\,\kappa$ of the chiral
  Chevalley--Eilenberg complex is the obstruction to the
  $\Einf$ structure: at $\hbar = 0$, the curvature vanishes
  and the algebra is commutative; at $\hbar \neq 0$, the
  $\Pthree$ bracket and the curvature together deform the
  commutative structure to a filtered $\Ethree$-chiral
  algebra.
  The deformation parameter is $\hbar \in H^3(\fg)[[\hbar]]$,
  matching the deformation space of
  Theorem~\textup{\ref{thm:e3-cs}(ii)}.
\item \textup{(Factorisation property.)}
  The $\Ethree$-chiral algebra
  $\CE^{\mathrm{ch}}_\hbar(\fg)$ is a factorisation algebra on
  $X$ in the sense that the restriction to disjoint
  opens $U_1 \sqcup U_2 \hookrightarrow X$ induces
  \begin{equation}\label{eq:factorisation}
    \CE^{\mathrm{ch}}_\hbar(\fg)\big|_{U_1 \sqcup U_2}
    \;\simeq\;
    \CE^{\mathrm{ch}}_\hbar(\fg)\big|_{U_1}
    \otimes
    \CE^{\mathrm{ch}}_\hbar(\fg)\big|_{U_2},
  \end{equation}
  where the tensor product is in the derived category of
  $\cD$-modules. The $\Ethree$ structure is compatible with
  the factorisation isomorphism.
\end{enumerate}
\end{theorem}

\begin{proof}
Part (i) is immediate from the filtration: the $\hbar$-adic
filtration is exhaustive, and at $\hbar = 0$ the curvature
and the cocycle term of the $\Pthree$ bracket both vanish,
leaving the uncurved commutative chiral coalgebra
$\CE^{\mathrm{ch}}(\fg_0)$.

Part (ii) requires three things: that the chiral
$\Pthree$ bracket is a degree-$(-2)$ Lie bracket
(Jacobi), that it satisfies the Leibniz rule with respect
to the commutative product (Poisson compatibility), and
that it is compatible with the $\Etwo$ structure from
the curve geometry, so that the combined data assembles
into an $\Ethree$-algebra. We prove each in turn.

\textit{Step 1: the $\Pthree$ Jacobi identity from the
PVA Jacobi.}
The PVA $\lambda$-bracket on $V_k(\fg)$ satisfies the
$(-1)$-shifted Jacobi identity
(see~\cite{Lorgat26II}, Appendix~A, Proposition~A.5):
\begin{equation}\label{eq:pva-jacobi-recall}
  \{a \mathbin{{}_\lambda} \{b \mathbin{{}_\mu} c\}\}
  - (-1)^{(|a|+1)(|b|+1)}\,
  \{b \mathbin{{}_\mu} \{a \mathbin{{}_\lambda} c\}\}
  = \{\{a \mathbin{{}_\lambda} b\}
  \mathbin{{}_{{\lambda+\mu}}} c\}.
\end{equation}
The chiral $\Pthree$ bracket
$\{-,-\}^{\mathrm{ch}}$
on $\CE^{\mathrm{ch}}(\fg_\hbar)$ is obtained from the
$\lambda$-bracket by specialising $\lambda \mapsto \partial$
(the $\cD$-module derivation) and passing to the chiral
Chevalley--Eilenberg complex. On generators
$\phi_a, \phi_b, \phi_c \in \fg^*[-1]$
(cohomological degree $1$, shifted bracket-degree
$|\phi_a| - 2 = -1$), the degree-$(-2)$ Jacobi identity
is:
\begin{equation}\label{eq:p3-jacobi}
  \{\phi_a, \{\phi_b, \phi_c\}^{\mathrm{ch}}\}^{\mathrm{ch}}
  + \text{graded cyclic}
  = 0.
\end{equation}
To derive this from~\eqref{eq:pva-jacobi-recall}, set
$\lambda = \mu = 0$ (the constant-coefficient specialisation
that extracts the $\Pthree$ bracket from the
$\lambda$-bracket). The left side
of~\eqref{eq:pva-jacobi-recall} becomes
$\{\phi_a, \{\phi_b, \phi_c\}^{\mathrm{ch}}\}^{\mathrm{ch}}
- (-1)^{(|\phi_a|+1)(|\phi_b|+1)}\,
\{\phi_b, \{\phi_a, \phi_c\}^{\mathrm{ch}}\}^{\mathrm{ch}}$,
and the right side becomes
$\{\{\phi_a, \phi_b\}^{\mathrm{ch}},
\phi_c\}^{\mathrm{ch}}$ (at $\lambda + \mu = 0$).
The PVA $\lambda$-bracket has degree $-1$ (it is
$(-1)$-shifted); the chiral $\Pthree$ bracket has
degree~$-2$, accounting for the additional degree shift
from the $\cD$-module derivation $\partial$ (conformal
weight~$1$). The Koszul sign
$(-1)^{(|\phi_a|+1)(|\phi_b|+1)}$ in the PVA identity
translates to $(-1)^{(|\phi_a|-2)(|\phi_b|-2)}$ in the
$\Pthree$ identity after incorporating the degree shift
$|\,{-}\,| \mapsto |\,{-}\,| - 2$ appropriate to a
degree-$(-2)$ bracket. On $\fg^*[-1]$ where
$|\phi_a| = 1$, both signs evaluate to
$(-1)^{(-1)(-1)} = -1$. Rearranging
gives~\eqref{eq:p3-jacobi}.

For general elements of $\Sym^c(\fg^*[-1])$ (not just
generators), the Jacobi identity extends by the Leibniz
rule (Step~2 below): the bracket is a biderivation of the
free commutative algebra, so a trilinear identity that
holds on generators holds on all elements.

\textit{Step 2: the Leibniz rule from PVA sesquilinearity.}
The PVA right Leibniz rule
(see~\cite{Lorgat26II}, Appendix~A, Proposition~A.6)
states:
\[
  \{a \mathbin{{}_\lambda} (b \cdot c)\}
  = \{a \mathbin{{}_\lambda} b\} \cdot c
  + (-1)^{(|a|+1)|b|}\,
  b \cdot \{a \mathbin{{}_\lambda} c\}
  + \int_0^\lambda
  \{\{a \mathbin{{}_\mu} b\}
  \mathbin{{}_\lambda} c\}\,d\mu.
\]
At $\lambda = 0$, the integral term vanishes
(the integration range collapses to a point), leaving:
\[
  \{a, b \cdot c\}^{\mathrm{ch}}
  = \{a, b\}^{\mathrm{ch}} \cdot c
  + (-1)^{(|a|+1)|b|}\,
  b \cdot \{a, c\}^{\mathrm{ch}}.
\]
Adjusting the Koszul sign from the $(-1)$-shifted
convention to the degree-$(-2)$ convention:
$(|a|+1)|b|$ becomes $(|a|-2)|b|$ after the degree shift,
recovering the $\Pthree$ Leibniz rule
\[
  \{a, b \cdot c\}^{\mathrm{ch}}
  = \{a, b\}^{\mathrm{ch}} \cdot c
  + (-1)^{(|a|-2)|b|}\,
  b \cdot \{a, c\}^{\mathrm{ch}},
\]
as stated in
Proposition~\ref{prop:chiral-p3-structure}(ii). This
identifies the chiral $\Pthree$ bracket as a derivation of
the commutative product
$\mu$ on $\Sym^c(\fg^*[-1]) \otimes \omega_X$: a
commutative algebra equipped with a compatible
degree-$(-2)$ Lie bracket is a $\Pthree$-algebra.

\textit{Step 3: compatibility with $\Etwo$ and the
$\Ethree$ assembly.}
The $\Etwo$-coalgebra structure on $\BarSig$ comes from
the curve geometry
(Proposition~\ref{prop:e3-structure}(i),
Warning~\ref{warn:e1-vs-e2-source}): the configuration
spaces $\Conf_n(X)$ provide operations parametrised by
$\Conf_k(\RR^2)$, the topological content of
$\Etwo$. The framed $\Etwo$ structure (from genus-$1$
sewing) equips the derived centre with a BV operator
$\Delta$ satisfying
$[f, g] = \Delta(f \cdot g) - (\Delta f) \cdot g
- (-1)^{|f|}\,f \cdot (\Delta g)$
(the BV relation~\eqref{eq:bv-rel-sl2}).
The $\Pthree$ bracket of degree~$-2$ supplements
this $\Etwo$ structure with one additional degree of
operadic freedom. The dimension count:
$\Etwo$ operations are parametrised by
$\Conf_k(\RR^2)$, and the $\Pthree$ bracket
(degree $-2 = -(3-1)$) corresponds to the fundamental
class of $S^2 = \Conf_2(\RR^3)/\RR^+$, the extra
generator needed for $\Ethree$ beyond~$\Etwo$.

The compatibility condition is that the $\Pthree$ bracket
must be a derivation of the Gerstenhaber bracket
$[-,-]$ of degree~$-1$ from the $\Etwo$ structure.
Since the Gerstenhaber bracket is determined by $\Delta$
and $\mu$ via the BV relation, it suffices to show that
$\{-,-\}^{\mathrm{ch}}$ commutes with $\Delta$ and is a
derivation of $\mu$ (already proved in Step~2). The BV
operator $\Delta$ arises from genus-$1$ sewing (the
$S^1$-action on the framed $\Etwo$ structure), while the
$\Pthree$ bracket arises from the arity-$2$ $r$-matrix
residue at genus~$0$. These operations act on independent
geometric data (genus-$1$ vs.\ genus-$0$, sewing vs.\
collision), so $\Delta$ commutes with
$\{-,-\}^{\mathrm{ch}}$ as $\cD$-module endomorphisms.
Given this, for any $a, b, c \in
\CE^{\mathrm{ch}}(\fg_\hbar)$:
\begin{align*}
  \{a, [b, c]\}^{\mathrm{ch}}
  &= \{a, \Delta(b \cdot c)
  - (\Delta b) \cdot c
  - (-1)^{|b|}\,b \cdot (\Delta c)\}^{\mathrm{ch}} \\
  &= \Delta\{a, b \cdot c\}^{\mathrm{ch}}
  - \{a, (\Delta b) \cdot c\}^{\mathrm{ch}}
  - (-1)^{|b|}\,\{a, b \cdot (\Delta c)\}^{\mathrm{ch}},
\end{align*}
where the first term uses
$[\Delta, \{a,-\}^{\mathrm{ch}}] = 0$. Expanding each
term by the Leibniz rule of Step~2 and collecting via the
BV relation yields
\[
  \{a, [b, c]\}^{\mathrm{ch}}
  = [\{a, b\}^{\mathrm{ch}}, c]
  + (-1)^{(|a|-2)|b|}\,[b, \{a, c\}^{\mathrm{ch}}]:
\]
the $\Pthree$ bracket is a derivation of the
Gerstenhaber bracket.

This is the mixed Leibniz rule that characterises
$\Ethree = \Etwo + \Pthree$: the $\Pthree$ bracket acts
as a derivation of the $\Etwo$ Gerstenhaber bracket. By
the recognition principle for $\mathsf{E}_n$-algebras
(Lurie~\cite{HA}, Theorem~5.1.2.2;
Francis~\cite{Francis2013}, Theorem~4.1): a commutative
algebra with a compatible $\Etwo$ structure and a
degree-$(-2)$ Lie bracket satisfying the mutual Leibniz
rules (Jacobi, Poisson compatibility with the product, and
derivation of the Gerstenhaber bracket) is an
$\Ethree$-algebra. The three verifications above establish
all the required relations.

Part (iii) combines two observations. The curvature
$m_0 = \tfrac{\hbar}{2h^\vee}\,\kappa$ is the bar-complex
encoding of the level departure (the identification
$\CE^{\mathrm{ch}}(\fg_\hbar) \simeq B^{\mathrm{ch}}(V_k(\fg))$
at $\hbar = k + h^\vee$). The cocycle term
$k\,(a,b)\,\partial$ of the $\Pthree$ bracket is linear
in~$k$ and hence linear in $\hbar$ (up to the shift
$k = \hbar - h^\vee$). Both vanish at $\hbar = 0$, so the
deformation is controlled by $\hbar$. The identification
of the deformation space as $\hbar H^3(\fg)[[\hbar]]$
then follows from the computation of
Theorem~\ref{thm:e3-cs}(ii): $H^3(\fg)$ is one-dimensional
for simple $\fg$, and the Cartan $3$-form
$\omega_3(X, Y, Z) = (X, [Y,Z])$ is the generator.

Part (iv) is the factorisation property of the chiral
bar complex. The chiral Chevalley--Eilenberg complex, being
a chiral coalgebra, satisfies the factorisation axiom
by construction (Beilinson--Drinfeld~\cite{BD04}, Chapter 3).
The $\Pthree$ bracket, being a $\cD$-module morphism
(Proposition~\ref{prop:chiral-p3-structure}(iii)), is
compatible with the factorisation isomorphisms.
\end{proof}


% ----------------------------------------------------------------
\subsection{The $\cD$-module structure and the KZ connection}
\label{subsec:chiral-e3-dmod}

The chiral $\Ethree$-algebra
$\CE^{\mathrm{ch}}_\hbar(\fg)$ is not an abstract
$\Ethree$-algebra: it is an $\Ethree$-algebra object
in the category of $\cD$-modules on $X$, or equivalently
on $\Ran(X)$. This $\cD$-module structure is the content
that distinguishes the chiral construction from the
topological CFG construction.

\begin{proposition}[The $\cD$-module structure]
\label{prop:chiral-e3-dmod}
The chiral $\Ethree$-algebra
$\CE^{\mathrm{ch}}_\hbar(\fg)$ carries a flat connection
\begin{equation}\label{eq:chiral-e3-connection}
  \nabla^{\mathrm{ch}}
  \;=\;
  d \;+\; \sum_{i < j}
  \frac{k\,\Omega_{ij}}{z_i - z_j}\,dz_{ij}
\end{equation}
on $\Conf_n(X)$, where $\Omega_{ij}$ is the Casimir element
acting on the $i$-th and $j$-th tensor factors. This is the
Knizhnik--Zamolodchikov connection.
\begin{enumerate}[label=\textup{(\roman*)}]
\item The flatness
  $(\nabla^{\mathrm{ch}})^2 = 0$ follows from the
  classical Yang--Baxter equation for the $r$-matrix
  $r(z) = k\,\Omega/z$
  \textup{(}trace-form convention; at $k = 0$, $r = 0$%
  \textup{)}: the curvature of $\nabla^{\mathrm{ch}}$
  is proportional to the LHS of the CYBE, which vanishes.
\item The $\Ethree$-algebra operations are horizontal
  with respect to $\nabla^{\mathrm{ch}}$: the cup product,
  the Gerstenhaber bracket, and the $\Pthree$ bracket are
  all flat sections of the corresponding $\Hom$-bundles.
  This is the $\cD$-linearity of
  Proposition~\textup{\ref{prop:chiral-p3-structure}(iii)}
  in the multi-point form.
\item The monodromy representation of~$\nabla^{\mathrm{ch}}$
  on $\pi_1(\Conf_n(X))$ gives the braid group action.
  For $X = \CC$, the monodromy of the KZ connection
  is computed by the Drinfeld associator
  $\Phi_{\mathrm{KZ}}$
  and recovers the braided monoidal structure on
  $\mathrm{Rep}(U_\hbar(\fg))$.
\end{enumerate}
\end{proposition}

\begin{proof}
The KZ connection on the $n$-point conformal blocks of
$V_k(\fg)$ is classical (Knizhnik--Zamolodchikov, 1984).
The presentation here identifies this connection as the
$\cD$-module structure on the chiral Chevalley--Eilenberg
complex, viewed as a factorisation algebra on $\Ran(X)$.

For (i): the curvature of $\nabla^{\mathrm{ch}}$ on
$\Conf_n(X)$ is
$[\nabla_i, \nabla_j]
= k^2\,[\Omega_{ij}/(z_i-z_j),\,
\Omega_{ik}/(z_i-z_k) + \Omega_{jk}/(z_j-z_k)]
\,dz_{ij} \wedge dz_{ik}$
plus cyclic permutations. The vanishing of this expression
is equivalent to the classical Yang--Baxter equation
$[r_{12}, r_{13}] + [r_{12}, r_{23}] + [r_{13}, r_{23}] = 0$
for $r_{ij} = k\,\Omega_{ij}/(z_i - z_j)$.

For (ii): the $\Ethree$ operations are defined in terms of
the PVA $\lambda$-bracket, which is itself a $\cD$-module
morphism. The cup product is a $\cD$-module map by the
commutative algebra structure on the factorisation algebra.
The $\Pthree$ bracket involves the derivation $\partial$
(from the $\lambda$ in the PVA bracket), which is the
$\cD$-module action; its compatibility with the connection
follows from the PVA sesquilinearity axiom.

For (iii): this is the Drinfeld--Kohno theorem
(Drinfeld~\cite{EK96}, Kohno, Kassel): the monodromy of the
KZ connection on $\Conf_n(\CC)$ factors through the pure
braid group $P_n$ and gives, via the Drinfeld associator
$\Phi_{\mathrm{KZ}} \in \exp(\hat{\mathfrak{t}}_3)$,
a braided monoidal structure on $\mathrm{Rep}(U_\hbar(\fg))$.
\end{proof}

% ----------------------------------------------------------------
\subsection{Comparison with CFG: formal disk global sections}
\label{subsec:chiral-e3-cfg-comparison}

The chiral $\Ethree$-algebra
$\CE^{\mathrm{ch}}_\hbar(\fg)$ lives on $\Ran(X)$ and
carries the KZ connection; the CFG algebra $\cA^\lambda$
lives on $\RR^3$ and is purely topological. The formal
disk $D = \Spec \CC[[z]]$ is the bridge: restricting to
$D$ trivializes the $\cD$-module structure and recovers the
CFG algebra as the fiber.

\begin{theorem}[Formal disk restriction recovers CFG]
\label{thm:chiral-e3-cfg}
Let $D = \Spec \CC[[z]]$ be the formal disk in~$X$.
There is a natural map of filtered $\Ethree$-algebras
\begin{equation}\label{eq:cfg-comparison}
  \Gamma(D,\, \CE^{\mathrm{ch}}_\hbar(\fg))
  \;\longrightarrow\;
  \cA^\lambda
\end{equation}
from the global sections of the chiral $\Ethree$-algebra
on the formal disk to the CFG $\Ethree$-algebra of
Theorem~\textup{\ref{thm:cfg}}, where $\lambda = \hbar$
\textup{(}the Chern--Simons coupling equals the departure
from critical level\textup{)}.
\begin{enumerate}[label=\textup{(\roman*)}]
\item \textup{(Associated graded.)}
  At the level of associated graded algebras,
  \begin{equation}\label{eq:cfg-assoc-graded}
    \mathrm{gr}\,\Gamma(D,\, \CE^{\mathrm{ch}}_\hbar(\fg))
    \;\simeq\;
    C^*(\fg)
    \;=\;
    \Sym(\fg^\vee[-1]),
  \end{equation}
  the Chevalley--Eilenberg cochains of $\fg$: both the
  chiral construction restricted to the formal disk and the
  CFG construction have the same classical limit.
\item \textup{(Matching of $\Pthree$ brackets.)}
  The restriction of the chiral $\Pthree$
  bracket~\eqref{eq:chiral-p3-total} to the formal disk
  recovers the CFG $\Pthree$ bracket.
  On the formal disk, the derivation $\partial = \partial_z$
  acts on $\CC[[z]]$; the global sections
  $\Gamma(D, \fg^*[-1] \otimes \omega_D)
  \cong \fg^*[-1][[z]]\,dz$ produce, after extracting the
  residue at $z = 0$, the finite-dimensional $\Pthree$
  bracket
  \[
    \{\phi_a, \phi_b\}_{\mathrm{CFG}}
    \;=\;
    f^{ab}{}_c\,\phi_c
    \;+\;
    k\,(a,b) \cdot \delta_0,
  \]
  where $\delta_0$ is the delta-function at the origin,
  the formal-disk avatar of $\partial$. Under the
  identification $\Gamma(D, \omega_D) \cong \CC\,dz$
  \textup{(}taking the constant mode\textup{)}, this
  recovers the CFG bracket on $C^*(\fg)$.
\item \textup{(Matching of deformation spaces.)}
  Both $\CE^{\mathrm{ch}}_\hbar(\fg)\big|_D$ and
  $\cA^\lambda$ are deformations of $C^*(\fg)$ over the
  base $\hbar\,H^3(\fg)[[\hbar]]$, with the same
  one-dimensional fibre at each $\hbar$-order for
  simple~$\fg$.
  The map~\eqref{eq:cfg-comparison} intertwines
  the $\hbar$-adic filtrations.
\end{enumerate}
\end{theorem}

\begin{proof}
The argument has two steps: construct the $\Ethree$-algebra
structure on $\Gamma(D, \CE^{\mathrm{ch}}_\hbar(\fg))$
by symmetric monoidal transport, then identify it with the
CFG algebra $\cA^\lambda$.

\smallskip
\noindent\textbf{Step 1: symmetric monoidal trivialization
on~$D$.}
The formal disk $D = \Spec \CC[[z]]$ is contractible:
every $\cD$-module on~$D$ (that is, a $\CC[[z]]$-module
with flat connection) is isomorphic to its fibre at the
origin, since every flat connection on~$D$ is
trivializable ($\pi_1(D^n) = 0$ for all~$n$, so there is
no monodromy obstruction). The global sections functor
\begin{equation}\label{eq:disk-equivalence}
  \Gamma(D, -) \colon
  \cD\textrm{-}\mathrm{mod}(D) \;\xrightarrow{\;\sim\;}\;
  \mathrm{Vect}
\end{equation}
is an equivalence of categories. This equivalence is
symmetric monoidal: the tensor product of $\cD$-modules
on~$D$ (the $!$-tensor product, which on the contractible
formal disk coincides with the $*$-tensor product and with
the ordinary tensor product of the underlying vector
spaces) corresponds under~\eqref{eq:disk-equivalence} to
the tensor product in $\mathrm{Vect}$.

An equivalence of symmetric monoidal categories preserves
all algebraic structures definable in terms of the monoidal
structure. In particular, it preserves $\Ethree$-algebra
structures: if $\cB$ is an $\Ethree$-algebra in
$\cD\textrm{-}\mathrm{mod}(D)$, then
$\Gamma(D, \cB)$ is an $\Ethree$-algebra in
$\mathrm{Vect}$, with structure maps obtained by applying
$\Gamma(D, -)$ to those of~$\cB$.

Applying this to
$\cB = \CE^{\mathrm{ch}}_\hbar(\fg)\big|_D$: the
restriction of the chiral $\Ethree$-algebra to~$D$ is an
$\Ethree$-algebra in $\cD\textrm{-}\mathrm{mod}(D)$ (by
restriction of the $\Ethree$-structure on~$X$), and the
equivalence~\eqref{eq:disk-equivalence} endows
$\Gamma(D, \CE^{\mathrm{ch}}_\hbar(\fg))$ with an
$\Ethree$-algebra structure in $\mathrm{Vect}$. This
$\Ethree$-structure is functorial and carries the
$\hbar$-adic filtration inherited from the chiral side.

\smallskip
\noindent\textbf{Step 2: identification with $\cA^\lambda$.}
It remains to identify
$\Gamma(D, \CE^{\mathrm{ch}}_\hbar(\fg))$ with the CFG
algebra $\cA^\lambda$ as filtered $\Ethree$-algebras,
where $\lambda = \hbar$. Both are filtered
$\Ethree$-algebras deforming the same classical limit over
the same base, and the identification proceeds through
three structural invariants.

Part (i): on~$D$, the $\cD$-module structure trivializes,
so the chiral Chevalley--Eilenberg complex reduces to the
ordinary one. The associated graded of the chiral
$\Ethree$-algebra on~$D$ is
\[
  \mathrm{gr}\,\Gamma(D,\, \CE^{\mathrm{ch}}_\hbar(\fg))
  \;=\;
  \Gamma(D,\, \mathrm{gr}\,\CE^{\mathrm{ch}}_\hbar(\fg))
  \;=\;
  \Gamma(D,\, \Sym^c(\fg^*[-1]) \otimes \omega_D).
\]
On the formal disk, $\omega_D \cong \cO_D\,dz$; the
equivalence~\eqref{eq:disk-equivalence} extracts the
constant mode in the $z$-expansion, giving
$\Sym(\fg^\vee[-1]) = C^*(\fg)$, the
Chevalley--Eilenberg cochains of~$\fg$ as a commutative
algebra. This is the associated graded of $\cA^\lambda$
(Theorem~\ref{thm:cfg}(i)).

Part (ii): the chiral $\Pthree$ bracket on~$X$
(equation~\eqref{eq:chiral-p3-total}) is defined via the
PVA $\lambda$-bracket. On the formal disk, the spectral
parameter $\lambda$ becomes the formal variable of the
completed enveloping algebra: the translation generator
$\partial_z$ acts on $\CC[[z]]$, and taking global sections
via~\eqref{eq:disk-equivalence} extracts the residue at
$z = 0$. The resulting bracket on
$\Gamma(D, \fg^*[-1] \otimes \omega_D)
\cong \fg^*[-1][[z]]\,dz$ is
\[
  \{\phi_a, \phi_b\}_{\mathrm{CFG}}
  \;=\;
  f^{ab}{}_{c}\,\phi_c
  \;+\;
  k\,(a,b) \cdot \delta_0,
\]
where $f^{ab}{}_{c}$ are the structure constants of~$\fg$,
$(a,b)$ is the Killing form, and $\delta_0$ is the
delta-function at the origin (the formal-disk avatar of the
derivation~$\partial$). Under the identification
$\Gamma(D, \omega_D) \cong \CC\,dz$, this is the CFG
$\Pthree$ bracket on $C^*(\fg)$
(Theorem~\ref{thm:cfg}(ii)).

Part (iii): the deformation theory of $\Ethree$-algebras
with fixed associated graded $C^*(\fg)$ is controlled by
the $\Ethree$-Hochschild cohomology, which for
$C^*(\fg)$ with~$\fg$ simple is
$\hbar\,H^3(\fg)[[\hbar]]$, one-dimensional at each
$\hbar$-order. Both
$\Gamma(D, \CE^{\mathrm{ch}}_\hbar(\fg))$ and
$\cA^\lambda$ are deformations of $C^*(\fg)$ over this
base (Theorem~\ref{thm:e3-cs}(ii) and
Theorem~\ref{thm:cfg}(iv) respectively). Since the
deformation space is one-dimensional at each order, and
both deformations have matching associated graded
(Part~(i)) and matching $\Pthree$ bracket (Part~(ii)),
they are isomorphic as filtered $\Ethree$-algebras.
The $\hbar$-adic filtration on the chiral side maps to
the $\lambda$-adic filtration on the CFG side under
$\lambda = \hbar = k + h^\vee$.
\end{proof}

\begin{remark}[The chain of spaces]
\label{rem:disk-vs-point}
The restriction passes through
$\mathrm{pt} \leftarrow D \to \mathbb{A}^1 \to \PP^1
\to X$, with each step adding geometric content: the formal
disk provides the derivation $\partial$ (cocycle term
$k(a,b)\partial$); $\mathbb{A}^1$ adds Arnold relations
and braid monodromy; $\PP^1$ and higher-genus $X$ add
compactness and modular curvature. None of these
restrictions is an identity; each step loses geometric
information.
\end{remark}


% ----------------------------------------------------------------
\subsection{The Khan--Zeng topological enhancement}
\label{subsec:khan-zeng-enhancement}

The chiral $\Ethree$-algebra is holomorphic: it depends on the
complex structure of $X$ through the KZ connection. At
non-critical level $k \neq -h^\vee$, the Sugawara Virasoro
element provides a homotopy between the holomorphic and
topological directions, upgrading the $\Ethree$ structure
from holomorphic to topological. At the critical level,
the Sugawara element is undefined and the enhancement fails.

\begin{proposition}[Topological enhancement via Sugawara]
\label{prop:khan-zeng-topological}
At generic non-critical level $k \neq -h^\vee$, the affine
Kac--Moody vertex algebra $V_k(\fg)$ possesses a Sugawara
Virasoro element
\begin{equation}\label{eq:sugawara-element}
  T_{\mathrm{Sug}} \;=\;
  \frac{1}{2(k + h^\vee)}
  \sum_{a=1}^{\dim \fg}
  \mathopen{:} J^a J_a \mathclose{:},
\end{equation}
which generates an inner action of the Virasoro algebra at
central charge $c = k\,\dim(\fg)/(k + h^\vee)$.

By Khan--Zeng~\cite{KhanZeng25}, the Virasoro element
upgrades the holomorphic-topological symmetry of the $3$d
gauge theory on $\Sigma \times \RR_t$ to fully topological
symmetry on $\Sigma \times \RR_t$. At the level of the
chiral $\Ethree$-algebra, this means:
\begin{enumerate}[label=\textup{(\roman*)}]
\item \textup{(Holomorphic-topological.)}
  The chiral $\Ethree$-algebra
  $\CE^{\mathrm{ch}}_\hbar(\fg)$ lives on $X \times \RR$,
  where $X$ carries holomorphic structure and $\RR$
  carries topological structure. The $\Etwo$ from the
  curve geometry is holomorphic; the $+1$ from the $\Pthree$
  bracket is topological \textup{(}from the $\RR$ factor%
  \textup{)}.
\item \textup{(Full topological.)}
  The Sugawara element provides an identification
  of the holomorphic direction in~$X$ with a topological
  direction: the stress-energy tensor $T_{\mathrm{Sug}}$
  generates translations along $X$, making them homotopically
  trivial. After this identification, the
  $\Ethree$-structure becomes \emph{topological} $\Ethree$:
  an $\Ethree$-algebra in the category of locally constant
  sheaves on $\RR^3 \cong X \times \RR$
  \textup{(}identifying $X \cong \CC \cong \RR^2$%
  \textup{)}.
\item \textup{(Failure at critical level.)}
  At $k = -h^\vee$, the Sugawara element~\eqref{eq:sugawara-element}
  is undefined \textup{(}the denominator $2(k+h^\vee)$
  vanishes\textup{)}. The critical-level chiral
  $\Ethree$-algebra $\CE^{\mathrm{ch}}_0(\fg)$ is
  genuinely holomorphic, not topological: the
  $\Etwo$-coalgebra structure on $\BarSig(\cA)$ comes from
  $\Conf_n(\RR^2)$ and is topological, but the
  \emph{chiral algebra} $\cA$ itself retains essential
  dependence on the complex structure of~$X$
  \textup{(}the OPE, the $\cD$-module, the KZ
  connection all see the holomorphic coordinate\textup{)}.
  Without the Sugawara element, there is no homotopy
  between the holomorphic and topological directions,
  so the full structure is not topological $\Ethree$.
  Perturbatively, at
  $k = -h^\vee + \hbar\lambda_1 + \cdots$,
  the topological enhancement holds order-by-order in
  $\hbar$ \textup{(}since $T_{\mathrm{Sug}}$ exists at
  any nonzero $\hbar$\textup{)}.
\end{enumerate}
\end{proposition}

\begin{proof}
Part (i) is the HT structure of the $3$d gauge theory
whose boundary chiral algebra is $V_k(\fg)$
(Costello--Gwilliam~\cite{CG17}). The factorisation
algebra of the HT theory lives on $\Sigma \times \RR$
with holomorphic dependence on $\Sigma$ and topological
dependence on $\RR$; the chiral $\Ethree$-algebra is the
algebraic encoding of this factorisation structure.

Part (ii).
Khan--Zeng~\cite{KhanZeng25} prove the following at the
classical level: an inner Virasoro element in a Poisson
vertex algebra with non-degenerate central charge provides
a homotopy between the holomorphic and topological
directions, upgrading the HT theory to a fully topological
theory. We extend this to the quantum vertex algebra
$V_k(\fg)$. The Sugawara construction
(Frenkel--Ben-Zvi~\cite{FBZ04}, Chapter~15) produces
$T_{\mathrm{Sug}} \in V_k(\fg)$ satisfying the exact
Virasoro OPE
\[
  T_{\mathrm{Sug}}(z)\,T_{\mathrm{Sug}}(w) \;\sim\;
  \frac{c/2}{(z-w)^4}
  + \frac{2\,T_{\mathrm{Sug}}(w)}{(z-w)^2}
  + \frac{\partial T_{\mathrm{Sug}}(w)}{z-w}
\]
at central charge $c = k\,\dim(\fg)/(k+h^\vee)$, and the
Sugawara OPE with the current generators
\[
  T_{\mathrm{Sug}}(z)\,J^a(w) \;\sim\;
  \frac{J^a(w)}{(z-w)^2}
  + \frac{\partial J^a(w)}{z-w}.
\]
The second relation gives $L_{-1} \cdot J^a = \partial J^a$:
the holomorphic translation operator $L_{-1} = (T_{\mathrm{Sug}})_{(1)}$
acts as $\partial$ on all generators. These are identities in
the quantum vertex algebra, not formal-deformation statements;
the normal ordering in~\eqref{eq:sugawara-element} absorbs
all quantum corrections.  The BRST operator $Q$ of the $3$d
gauge theory satisfies $[Q, b] = L_{-1}$ where $b$ is the
antighost zero-mode (this is the standard BRST relation for
the Sugawara Virasoro; see~\cite{FBZ04}, \S15.4), so
$L_{-1}$ is BRST-exact. The Khan--Zeng
argument~\cite{KhanZeng25} then applies verbatim: BRST-exactness
of holomorphic translation is the datum required to upgrade
the HT symmetry to full topological symmetry. At generic
$k$, the central charge
$c = k\,\dim(\fg)/(k+h^\vee) \neq 0$, so the hypothesis is
satisfied.

Part (iii): at $k = -h^\vee$, the Feigin--Frenkel centre
$Z(V_{-h^\vee}(\fg)) = \mathrm{Fun}(\mathrm{Op}_{\fg^\vee}(D))$
(Feigin--Frenkel) is a commutative
algebra with non-trivial dependence on the complex structure
of the formal disk (through the notion of $\fg^\vee$-oper).
The absence of the Sugawara element means no homotopy
between holomorphic and topological directions exists, and
the chiral $\Ethree$-algebra is genuinely holomorphic.
The perturbative statement follows because the Sugawara
element exists as a formal power series in $\hbar$ for
$\hbar \neq 0$.
\end{proof}


% ----------------------------------------------------------------
\subsection{Example: $\mathfrak{sl}_2$ chiral $\Ethree$-algebra}
\label{subsec:sl2-chiral-e3}

The general construction specialises cleanly to
$\fg = \mathfrak{sl}_2$: the chiral $\Pthree$ bracket is
determined by six pairings on the generators
$\phi_e, \phi_f, \phi_h$, each consisting of a Lie bracket
term (from the structure constants) and a cocycle term
(proportional to the level $k$, vanishing at $k = 0$).
Restricting to the formal disk and extracting the zero mode
recovers the CFG bracket of
Proposition~\ref{prop:e3-explicit-sl2}.

\begin{example}[Explicit chiral $\Pthree$ bracket for
$\mathfrak{sl}_2$]
\label{ex:sl2-chiral-e3}
For $\fg = \mathfrak{sl}_2$ with basis $\{e, f, h\}$,
Killing form $(e,f) = (f,e) = 1$, $(h,h) = 2$,
structure constants $[e,f] = h$, $[h,e] = 2e$, $[h,f] = -2f$,
and $h^\vee = 2$, the chiral $\Pthree$ bracket
(Construction~\ref{constr:chiral-p3-bracket}) on the
generators $\phi_e, \phi_f, \phi_h \in \mathfrak{sl}_2^*[-1]$
is:
\begin{align}
  \label{eq:sl2-p3-ef}
  \{\phi_e, \phi_f\}^{\mathrm{ch}}
  &= \phi_h + k\,\partial, \\
  \label{eq:sl2-p3-he}
  \{\phi_h, \phi_e\}^{\mathrm{ch}}
  &= 2\,\phi_e, \\
  \label{eq:sl2-p3-hf}
  \{\phi_h, \phi_f\}^{\mathrm{ch}}
  &= -2\,\phi_f, \\
  \label{eq:sl2-p3-hh}
  \{\phi_h, \phi_h\}^{\mathrm{ch}}
  &= 2k\,\partial, \\
  \label{eq:sl2-p3-ee}
  \{\phi_e, \phi_e\}^{\mathrm{ch}}
  &= 0, \\
  \label{eq:sl2-p3-ff}
  \{\phi_f, \phi_f\}^{\mathrm{ch}}
  &= 0.
\end{align}
The Lie bracket terms ($\phi_h$ in~\eqref{eq:sl2-p3-ef},
$2\phi_e$ in~\eqref{eq:sl2-p3-he},
$-2\phi_f$ in~\eqref{eq:sl2-p3-hf}) persist at $k = 0$;
the cocycle terms ($k\,\partial$ in~\eqref{eq:sl2-p3-ef}
and~\eqref{eq:sl2-p3-hh}) vanish at $k = 0$
(AP126 check: at $k = 0$, the level-dependent terms
vanish, and the bracket reduces to the Chevalley--Eilenberg
coderivation).

At the critical level $k = -2$ ($\hbar = 0$), the cocycle
terms become $-2\,\partial$, and the chiral
$\Ethree$-algebra degenerates: the curvature $m_0 = 0$
(uncurved), and the commutative factorisation structure is
undeformed.

Taking formal disk global sections and extracting the zero
mode recovers the CFG $\Pthree$ bracket of
Proposition~\ref{prop:e3-explicit-sl2}:
$\{\phi_a, \phi_b\}_{\mathrm{CFG}}
= f^{ab}{}_c\,\phi_c + k\,(a,b)\,\delta_0$, which at the
cohomological level gives
$\{X, Y\} = h_{\mathrm{KZ}}\,(X,Y)$ for
$X, Y \in \HH^1 = \mathfrak{sl}_2$
(equation~\eqref{eq:e3-p3-killing}), with
$h_{\mathrm{KZ}} = 1/(k+2)$.
\end{example}

% ================================================================
\section{Examples}
\label{sec:examples}

The four shadow classes $G$, $L$, $C$, $M$ produce
qualitatively different ordered chiral homology theories.
Class $G$ (Heisenberg) has regular singularities, finite
shadow tower, and trivial kernel $\ker(\av)$. Class $L$
(affine Kac--Moody) has regular singularities but a
non-trivial $R$-matrix kernel. Class $C$
($\beta\gamma$ system) has irregular singularities of
bounded Poincar\'e rank. Class $M$ (Virasoro) has irregular
singularities of unbounded rank and an infinite shadow tower.

\begin{example}[Heisenberg: class $G$, $r = 2$;
explicit ordered chiral Hochschild cohomology]
\label{ex:heisenberg}
The Heisenberg algebra $\cH_k$ has a single generating
field $J(z)$ of conformal weight $1$, OPE
$J(z)J(w) \sim k/(z-w)^2$, and
$r$-matrix $r(z) = k/z$ (trace-form convention;
at $k=0$, $r = 0$). The Koszul invariant is $\kappa = k$.
The algebra is $\Einf$-chiral (symmetric OPE; no simple
pole) and class $G$ (shadow depth $r=2$: $S_4 = 0$, so
$\Delta = 8\kappa S_4 = 0$).
By the formality bridge (Theorem~\ref{thm:formality-bridge}),
ordered and symmetric chiral homologies agree.
We compute the ordered
chiral Hochschild cohomology
$\HH^{*,\mathrm{ord}}_{\mathrm{ch}}(\cH_k, \cH_k)$ on the
formal disk $D$ and on $\PP^1$ explicitly, arity by arity.

\medskip
\noindent
\textbf{Setup.}
The ordered chiral Hochschild cochain complex
(Definition~\ref{def:ch-hochschild}) is
\[
\HH^{*,\mathrm{ord}}_{\mathrm{ch}}(\cH_k, \cH_k)
= \bigoplus_{n \geq 0}
\Bigl(\Omega^*\bigl(\overline{\FM}_n^{\mathrm{ord}}(\CC)\bigr)
\otimes \Hom(\cH_k^{\otimes n}, \cH_k),\;
d_{\mathrm{dR}} + d_{\mathrm{Hoch}} + \nabla_{\mathrm{KZ}}
\Bigr),
\]
where $d_{\mathrm{Hoch}}$ is the Hochschild differential
(encoding OPE residues at boundary strata) and
$\nabla_{\mathrm{KZ}}$ is the KZ connection.
The bar complex is $B(\cH_k) = T^c(s^{-1}\bar{\cH}_k)$
(desuspended augmentation ideal; $|s^{-1}v| = |v| - 1$),
with deconcatenation coproduct.

\medskip
\noindent
\textbf{Arity $0$: the vacuum.}
The chiral centre $Z(\cH_k) = \{a : a_{(n)}b = 0$
for all $b$, $n \geq 0\}$ consists of the scalars only
(since $J_{(1)}J = k \neq 0$ but $J_{(0)}J = 0$).
\begin{equation}\label{eq:heis-hh0}
\HH^0 = \CC.
\end{equation}

\medskip
\noindent
\textbf{Arity $1$: outer derivations.}
A chiral derivation $D$ satisfies $D(J) = c \cdot \mathbf{1}$
for $c \in \CC$ (the Leibniz rule and the absence of fields
with regular OPE against $J$ force this). Since
$J_{(0)}J = 0$ (no simple pole), all inner derivations vanish:
$\mathrm{Inn}(\cH_k) = 0$.
\begin{equation}\label{eq:heis-hh1}
\HH^1 = \CC,
\end{equation}
generated by the translation $D(J) = \mathbf{1}$.
The classical Hochschild $\mathrm{HH}^1$ of the Weyl algebra
is $0$ (every derivation is inner); the discrepancy arises
because chiral inner derivations use $a_{(0)}$, not
$[a, -]$.

\medskip
\noindent
\textbf{Arity $2$: level deformation and the KZ connection.}
Setting $w = z_1 - z_2$, the KZ connection is
\begin{equation}\label{eq:kz-heis-2}
\nabla_{\mathrm{KZ}} = d + k\, d\log w,
\end{equation}
with solution $\Phi(w) = w^k$ and scalar monodromy
\begin{equation}\label{eq:heis-monodromy}
M = e^{2\pi i k}.
\end{equation}
For generic $k \notin \ZZ$, the twisted de~Rham cohomology
of $\CC^\times$ vanishes. The level deformation cocycle
$\phi_k(J \otimes J) = k \cdot \mathbf{1}$ represents
\begin{equation}\label{eq:heis-hh2}
[\phi_k] \in \HH^2 = \CC,
\end{equation}
the unique deformation
$J(z)J(w) \sim (k+\epsilon)/(z-w)^2$.

\medskip
\noindent
\textbf{Arity $n \geq 3$: vanishing.}
The Heisenberg is chirally Koszul with resolution of
length $2$, so the cochain complex at arity $n \geq 3$ is
acyclic:
\begin{equation}\label{eq:heis-higher-vanish}
\HH^{n,\mathrm{ord}}_{\mathrm{ch}}(\cH_k, \cH_k) = 0
\quad \text{for } n \geq 3.
\end{equation}

\medskip
\noindent
\textbf{Summary on $D$.}
The ordered chiral Hochschild cohomology on the formal disk is
\begin{equation}\label{eq:heis-hh-summary}
\HH^{n,\mathrm{ord}}_{\mathrm{ch}}(\cH_k, \cH_k) =
\begin{cases}
\CC & n = 0 \;\text{(vacuum: centre $= \CC \cdot \mathbf{1}$)}, \\
\CC & n = 1 \;\text{(outer derivation: $D(J) = \mathbf{1}$)}, \\
\CC & n = 2 \;\text{(level deformation: $k \mapsto k + \epsilon$)}, \\
0 & n \geq 3 \;\text{(Koszul resolution of length $2$)},
\end{cases}
\end{equation}
with Poincar\'e polynomial $P_{\cH_k}(t) = 1 + t + t^2$.
The total dimension is $3$, and the cohomological amplitude
is $[0, 2]$, consistent with Theorem~H.
The derived chiral centre is
$Z^{\mathrm{der}}_{\mathrm{ch}}(\cH_k) \cong
\CC \oplus \CC[-1] \oplus \CC[-2]$
as a graded vector space.

\medskip
\noindent
\textbf{The computation on $\PP^1$.}
The computation on $\PP^1$ gives the same result:
\begin{equation}\label{eq:heis-hh-P1}
\HH^{n,\mathrm{ord}}_{\mathrm{ch,\PP^1}}(\cH_k, \cH_k)
\cong
\HH^{n,\mathrm{ord}}_{\mathrm{ch},D}(\cH_k, \cH_k)
\quad \text{for all } n.
\end{equation}
The factorisation $\cD$-module is insensitive to the
compactification for chirally Koszul algebras: the Koszul
resolution has length $2$, so only arities $0$, $1$, $2$
contribute, and at these arities the global geometry of $\PP^1$
introduces no new cohomology beyond what the formal disk sees.

\medskip
\noindent
\textbf{The $\Ethree$ structure.}
The derived chiral centre
$Z^{\mathrm{der}}_{\mathrm{ch}}(\cH_k)$ is an
$\Ethree$-algebra by Proposition~\ref{prop:e3-structure}.
The obstruction space
$\HH^3_{\Etwo}(\BarSig(\cH_k), \BarSig(\cH_k))
\cong \Lambda^3(\CC^1) = 0$
vanishes, so the $\Ethree$-algebra structure on
$Z^{\mathrm{der}}_{\mathrm{ch}}(\cH_k)
\cong \CC \oplus \CC[-1] \oplus \CC[-2]$
is the unique (trivial) one.
The $3$d theory is abelian Chern--Simons at level $k$;
the triviality reflects the absence of perturbative quantum
corrections in the abelian theory.

\medskip
\noindent
\textbf{Comparison with topological Hochschild.}
By the formality bridge (Theorem~\ref{thm:formality-bridge}),
$\HH^{*,\mathrm{ord}}_{\mathrm{ch}}(\cH_k, \cH_k)
\simeq \HH^{*,\mathrm{top}}(\cH_k, \cH_k)$:
no information is lost in the passage from holomorphic
(ordered, $z$-dependent) to topological
($S^1$, $z$-independent).
\end{example}

\begin{example}[Affine Kac--Moody: class $L$, $r = 3$;
explicit computation for $\widehat{\mathfrak{sl}}_2$]
\label{ex:km}
We carry out the ordered chiral Hochschild computation for
$V_k(\mathfrak{sl}_2)$ on the formal disk, the simplest
non-abelian Kac--Moody algebra.

\smallskip
\noindent\textbf{Data.}
$\fg = \mathfrak{sl}_2$, $h^\vee = 2$, $\dim(\fg) = 3$.
Generators $\{e, f, h\}$ with brackets $[e,f] = h$,
$[h,e] = 2e$, $[h,f] = -2f$.
Normalised invariant form: $(e,f) = (f,e) = 1$,
$(h,h) = 2$. Inverse Casimir:
$\Omega = e \otimes f + f \otimes e
+ \tfrac{1}{2}h \otimes h$.
% AP1: kappa from landscape_census.tex:
% kappa(V_k(sl_2)) = 3(k+2)/4; k=0 -> 3/2; k=-2 -> 0.
The Koszul invariant is
$\kappa = 3(k + 2)/4$.
The $r$-matrix in trace-form convention is
$r(z) = k\Omega/z$ (at $k = 0$, $r = 0$);
in KZ normalisation,
$r(z) = \Omega/((k + 2)z)$.
Class $L$ (shadow depth $r = 3$): $\ell_1, \ell_2, \ell_3$
non-trivial, $\ell_k = 0$ for $k \geq 4$.

\smallskip
\noindent\textbf{Arity 2: KZ conformal blocks.}
The ordered chiral Hochschild at arity $2$ is the de~Rham
cohomology of $\Conf_2(\CC)$ with coefficients in the
KZ local system. The KZ equation at arity $2$, setting
$z = z_1 - z_2$, is
\begin{equation}\label{eq:kz-sl2-arity2}
  \frac{\partial\Phi}{\partial z}
  = \frac{\Omega}{(k + 2)\,z}\,\Phi.
\end{equation}
Since $V_k(\mathfrak{sl}_2)$ is $\Einf$-chiral, the
coefficients live in
$V_k(\mathfrak{sl}_2)^{\otimes 2}$, which decomposes under
$\mathfrak{sl}_2$ by Clebsch--Gordan. For the adjoint
representation acting on itself ($j = 1$):
$V_1 \otimes V_1 = V_0 \oplus V_1 \oplus V_2$.
The Casimir eigenvalue of $\Omega$ on the isospin-$J$
component is
\begin{equation}\label{eq:kz-adj-casimir}
  \lambda_J
  = \frac{C_2(J) - C_2(1) - C_2(1)}{2}
  = \frac{J(J+1) - 4}{2},
\end{equation}
where $C_2(j) = j(j+1)$ in the normalisation
$(\theta|\theta) = 2$. The eigenvalues are:
\[
  \lambda_0 = -2,\quad \lambda_1 = -1,\quad \lambda_2 = 1.
\]
On each isospin component, the KZ
equation~\eqref{eq:kz-sl2-arity2} has the explicit solution
$\Phi_J(z) = z^{\lambda_J/(k+2)}$. The monodromy around
$z = 0$ (the diagonal $\Delta_{12} \subset X^2$) is
\[
  M_J = \exp\!\bigl(2\pi i\,\lambda_J/(k + 2)\bigr).
\]
The de~Rham cohomology $H^*_{\mathrm{dR}}(\Conf_2(\CC),
\nabla_{\mathrm{KZ}})$ decomposes as:
\begin{itemize}
\item $H^0 = \ker(M - \id)$: the subspace on which the
  monodromy is trivial. At generic $k$, this is zero
  (no isospin channel has
  $\lambda_J/(k+2) \in \ZZ$).
\item $H^1 = \mathrm{coker}(M - \id)$: the space of
  KZ conformal blocks at arity $2$. At generic $k$,
  $H^1 \cong V_0 \oplus V_1 \oplus V_2$
  (three-dimensional over the isospin decomposition:
  one block per channel).
\end{itemize}
The averaging map
$\av_2\colon H^1_{\mathrm{ord}} \to H^1_{\Sigma}$
projects each conformal block to its $\Sigma_2$-coinvariant.
Under the flip $z \mapsto -z$, the monodromy on $V_J$ picks
up a sign $(-1)^J$ (from the permutation of the two
adjoint insertions); at generic $k$, $\av_2$ is surjective
with $\ker(\av_2)$ carrying the odd-$J$ blocks.
The scalar $\kappa = \av(r(z)) + \dim(\fg)/2
= 3k/(2 \cdot 2) + 3/2 = 3(k+2)/4$
% AP1: cross-check: kappa(V_k(sl_2)) = dim(sl_2)(k+h^v)/(2h^v) = 3(k+2)/4.
% k=0 -> 3/2. k=-2 -> 0. Verified.
is the complete information surviving the averaging.

\smallskip
\noindent\textbf{Arity 3: the Drinfeld associator.}
At arity $3$, the KZ system on $\Conf_3(\CC)$ is
\begin{equation}\label{eq:kz-3pt}
  d\Phi = \frac{1}{k + 2}
  \biggl(\frac{\Omega_{12}}{z_1 - z_2}\,d(z_1 - z_2)
  + \frac{\Omega_{13}}{z_1 - z_3}\,d(z_1 - z_3)
  + \frac{\Omega_{23}}{z_2 - z_3}\,d(z_2 - z_3)
  \biggr)\Phi.
\end{equation}
The fundamental group $\pi_1(\Conf_3(\CC)) = PB_3$
(the pure braid group on three strands) has generators
$\sigma_{12}^2, \sigma_{23}^2$ (Dehn twists around the
diagonals). The KZ monodromy gives a representation
$PB_3 \to \mathrm{GL}(V_k(\fg)^{\otimes 3})$.
The Drinfeld associator
$\Phi_{\mathrm{KZ}} \in
\exp(\widehat{\Lie}(t_{12}, t_{23}))$ is the regularised
holonomy along the path from the boundary stratum
$|z_1 - z_2| \ll |z_2 - z_3|$ to the stratum
$|z_2 - z_3| \ll |z_1 - z_2|$ in the real locus of
$\overline{\FM}_3^{\mathrm{ord}}(\CC)$.

The kernel of the averaging map $\av_3$ is the space of
ordered data invisible to $\BarSig$. For
$\fg = \mathfrak{sl}_2$, this kernel is one-dimensional:
$\ker(\av_3) \cong \Lambda^3(\CC^3) = \CC$, carrying the
Drinfeld associator. The associator $\Phi_{\mathrm{KZ}}$
is the unique (up to $\GRT_1$-action) non-trivial element
of this kernel.

For general $\fg$, the ordered data at arity $3$ is
controlled by
$\Lambda^3(\fg) = \Lambda^3(\CC^{\dim\fg})$, with
$\dim \Lambda^3(\fg) = \binom{\dim\fg}{3}$.
At $\fg = \mathfrak{sl}_2$: $\binom{3}{3} = 1$.

\smallskip
\noindent\textbf{$\Ethree$ obstruction.}
The symmetric bar $\BarSig(V_k(\mathfrak{sl}_2))$ carries
an $\Etwo$-coalgebra structure from the curve geometry
(Proposition~\ref{prop:e3-structure}(i): the coshuffle
coproduct on configuration spaces of the curve, not
from bar-cobar). The obstruction to
promoting $\BarSig$ to an $\Ethree$-coalgebra lives in
\begin{equation}\label{eq:e3-obstruction-km}
  \HH^3_{\Etwo}(\BarSig, \BarSig)
  \;\cong\; H^3(\mathfrak{sl}_2)
  \;\cong\; \Lambda^3(\mathfrak{sl}_2^*)^{\mathfrak{sl}_2}
  \;=\; \CC,
\end{equation}
which is one-dimensional and carries the non-trivial
obstruction class (Warning~\ref{warn:bar-not-e3}).
The $\Ethree$-algebra is the derived chiral centre
$Z^{\mathrm{der}}_{\mathrm{ch}}(V_k(\mathfrak{sl}_2))$,
not $\BarSig$ itself.

\smallskip
\noindent\textbf{Identification with CFG.}
The $\Ethree$-deformation space is
$\hbar H^3(\mathfrak{sl}_2)[[\hbar]]
\cong \hbar\CC[[\hbar]]$:
one-dimensional at each $\hbar$-order.
For $\mathfrak{sl}_2$,
$H^3(\mathfrak{sl}_2) = \CC$ (generated by the
three-form $\omega_3 = (e^* \wedge f^* \wedge h^*) \circ
[\cdot, \cdot]$, the Cartan form).
The deformation parameter is the departure from critical
level: $k = -2 + \hbar\lambda_1 + \hbar^2\lambda_2 + \cdots$.
By Theorem~\ref{thm:chiral-e3-cfg}, the formal disk
restriction $\Gamma(D, \CE^{\mathrm{ch}}_\hbar(\mathfrak{sl}_2))$
recovers the CFG filtered $\Ethree$-algebra
$C^*_\hbar(\mathfrak{sl}_2)$ of perturbative Chern--Simons
quantisations; the full identification of the deformation
families is Conjecture~\ref{conj:e3-identification}.
At the associative level, the Drinfeld--Kohno theorem
identifies the KZ monodromy with the $R$-matrix of
$U_q(\mathfrak{sl}_2)$ at $q = \exp(\pi i/(k + 2))$:
the chiral quantum group equivalence
(Theorem~\ref{thm:chiral-qg-equiv}) lifts this from
monodromy representations to the algebra.

\smallskip
\noindent\textbf{Summary of ordered chiral Hochschild for
$\widehat{\mathfrak{sl}}_2$.}
\begin{center}
\begin{tabular}{@{} c l l @{}}
\toprule
Arity $n$ & Ordered datum & Symmetric shadow \\
\midrule
$1$ & $V_k(\mathfrak{sl}_2)$ & $V_k(\mathfrak{sl}_2)$ \\
$2$ & KZ local system; conformal blocks
      $V_0 \oplus V_1 \oplus V_2$
    & scalar $\kappa = 3(k+2)/4$ \\
$3$ & KZ associator $\Phi_{\mathrm{KZ}}$;
      $\ker(\av_3) = \CC$
    & trivial (associator lost) \\
$\geq 4$ & higher KZ holonomy
    & trivial \\
\bottomrule
\end{tabular}
\end{center}

The KZ connection~\eqref{eq:kz-km} has simple poles:
regular singularities. The $3$d theory is Chern--Simons
gauge theory with gauge group $G$.
\end{example}

\begin{example}[$\beta\gamma$ system: class $C$, $r = 4$;
explicit ordered chiral Hochschild cohomology]
\label{ex:betagamma}
The $\beta\gamma$ system at conformal weight $\lambda$ has two
generating fields $\beta(z)$ of weight $\lambda$ and $\gamma(z)$
of weight $1 - \lambda$, with OPE
\[
\beta(z)\gamma(w) \sim \frac{1}{z - w}, \qquad
\gamma(z)\beta(w) \sim \frac{1}{z - w}, \qquad
\beta(z)\beta(w) \sim 0, \qquad
\gamma(z)\gamma(w) \sim 0.
\]
The central charge is $c_{\beta\gamma}(\lambda)
= 2(6\lambda^2 - 6\lambda + 1)$ and the Koszul invariant is
$\kappa = c_{\beta\gamma}/2 = 6\lambda^2 - 6\lambda + 1$.
(Verification: at $\lambda = 1/2$, $c_{\beta\gamma} = -1$;
at $\lambda = 2$, $c_{\beta\gamma} = 26$.)
The algebra is $\Einf$-chiral (the OPE is symmetric after
clearing the simple pole by locality) and class $C$ (shadow
depth $r = 4$: the quartic contact shadow is nontrivial,
but the quintic obstruction vanishes by rank-one abelian
rigidity).

By the formality bridge (Theorem~\ref{thm:formality-bridge}),
ordered and symmetric chiral homologies agree.
The distinctive feature of this example is that the KZ connection has
\emph{unipotent} monodromy (rather than scalar or semisimple),
and the nilpotency $\Theta^2 = 0$ makes the parallel transport
computable in closed form at every arity.
The $\beta\gamma$ system is the \emph{hydrogen atom} for
ordered chiral homology with unipotent local systems: all KZ
solutions are explicit polynomials in logarithms, with no
hypergeometric functions and no Drinfeld associator.

\medskip
\noindent
\textbf{The collision residue and $R$-matrix.}
The generating space is $V = \CC\beta \oplus \CC\gamma$
(two-dimensional). The OPE has a simple pole between $\beta$
and $\gamma$. After $d$-log absorption (pole order $1 - 1 = 0$),
the collision residue is a \emph{constant} nilpotent endomorphism
\begin{equation}\label{eq:bg-theta}
  \Theta \in \End(V \otimes V), \qquad
  \Theta(\beta \otimes \gamma) = \mathbf{1}, \quad
  \Theta(\gamma \otimes \beta) = \mathbf{1}, \quad
  \Theta(\beta \otimes \beta) = 0, \quad
  \Theta(\gamma \otimes \gamma) = 0,
\end{equation}
where $\mathbf{1}$ denotes the vacuum in the full algebra $A$
(not in $V \otimes V$). The map $\Theta$ sends
$V \otimes V \to A$, with image in the one-dimensional
vacuum subspace $\CC \cdot \mathbf{1} \subset A$.
The key algebraic fact:
$\Theta^2 = 0$ (nilpotent of order $2$), because the image
of $\Theta$ is the vacuum subspace, and
$\Theta$ annihilates the vacuum. Concretely: $\Theta$ has
rank $1$ on the four-dimensional space $V^{\otimes 2}$
(both $\beta \otimes \gamma$ and $\gamma \otimes \beta$
map to $\mathbf{1}$, hence the image is one-dimensional),
with
$\ker(\Theta) \supset \CC(\beta \otimes \beta)
\oplus \CC(\gamma \otimes \gamma)$
(the ``same-type'' tensors, on which the OPE is regular)
and $\mathrm{im}(\Theta) = \CC \cdot \mathbf{1}$.

The $R$-matrix is the monodromy $R = \exp(2\pi i\,\Theta)$.
Since $\Theta^2 = 0$, the exponential terminates:
\begin{equation}\label{eq:bg-rmatrix}
  R_{\beta\gamma} = I + 2\pi i\,\Theta.
\end{equation}
This termination at the linear term is the defining feature
of class~$C$. For affine Kac--Moody (class~$L$), the
$R$-matrix is an infinite series in the spectral parameter;
for $\beta\gamma$, nilpotency truncates the series.

\medskip
\noindent
\textbf{The KZ connection.}
At arity $n$, the KZ connection on
$\Conf_n^{\mathrm{ord}}(\CC)$ is
\begin{equation}\label{eq:kz-bg}
  \nabla_n^{\beta\gamma}
  = d + \sum_{1 \leq i < j \leq n}
  \Theta_{ij}\, \omega_{ij}, \qquad
  \omega_{ij} = d\log(z_i - z_j),
\end{equation}
where $\Theta_{ij} \in \End(V^{\otimes n})$ acts as $\Theta$
on the $i$-th and $j$-th tensor factors and as the identity
elsewhere. The connection is logarithmic: regular singularities
along every diagonal $\Delta_{ij}$. This is the analogue of
the KZ connection for affine KM, with the simplification
that $\Theta_{ij}^2 = 0$.

Flatness: the condition $(\nabla_n^{\beta\gamma})^2 = 0$
reduces to the classical Yang--Baxter equation for $\Theta$:
\begin{equation}\label{eq:cybe-bg}
  [\Theta_{12}, \Theta_{13}]
  + [\Theta_{12}, \Theta_{23}]
  + [\Theta_{13}, \Theta_{23}] = 0
  \quad \text{in } \End(V^{\otimes 3}).
\end{equation}
This holds because $\Theta$ is the collision residue of a
vertex algebra OPE. Each commutator
$[\Theta_{ij}, \Theta_{jk}]$ produces a ``triple collision''
term (contraction at all three sites via the shared index $j$),
and the three such terms cancel by the cyclic symmetry
of the Arnold relation.

\medskip
\noindent
\textbf{Explicit KZ solutions via nilpotency.}
Since $\Theta^2 = 0$, the parallel transport is polynomial
in logarithms. At arity $2$:
\begin{equation}\label{eq:bg-kz-sol-2}
  \Phi(w) = \exp(\Theta_{12} \log w)
  = I + \Theta_{12}\,\log w,
\end{equation}
with monodromy $R = I + 2\pi i\,\Theta_{12}$.
At arity $3$, the solution terminates at degree two in
logarithms (triple contractions vanish on the
two-dimensional $V$), and the Drinfeld associator is
trivial. Since $\Theta^2 = 0$, the iterated-integral
corrections from non-commutativity of the connection
vanish: the Picard series
$\Phi = I + \sum_{i<j}\Theta_{ij}\log(z_i - z_j)
+ \sum_{(i<j),(k<l)}\Theta_{ij}\Theta_{kl}
\log(z_i - z_j)\log(z_k - z_l) + \cdots$
reduces to ordinary products of logarithms (rather than
Chen iterated integrals), because the ordering ambiguity
in the iterated integral $\int \Theta_{ij}\omega_{ij}
\circ \Theta_{kl}\omega_{kl}$ versus
$\int \Theta_{kl}\omega_{kl} \circ \Theta_{ij}\omega_{ij}$
differs by the commutator $[\Theta_{ij}, \Theta_{kl}]$,
which vanishes when $\Theta^2 = 0$ (the image of each
$\Theta_{ij}$ is the vacuum, annihilated by every
$\Theta_{kl}$).

\medskip
\noindent
\textbf{The chiral Hochschild computation.}
The $\beta\gamma$ system is chirally Koszul with
Koszul resolution of length~$2$ (generating space
$V = \CC\beta \oplus \CC\gamma$, relation space
$R = \CC \cdot (\beta_{(0)}\gamma - \mathbf{1})$), so only
arities $0$, $1$, $2$ contribute.

\medskip
\noindent
\textbf{Arity $0$: the vacuum.}
Since $\beta_{(0)}\gamma = \mathbf{1} \neq 0$, neither
$\beta$ nor $\gamma$ is central; only the scalars survive.
\begin{equation}\label{eq:bg-hh0}
\HH^0 = \CC.
\end{equation}

\medskip
\noindent
\textbf{Arity $1$: outer derivations.}
The simple-pole OPE $\beta_{(0)}\gamma = \mathbf{1}$ makes
$a_{(0)}$ surjective on the generating space, so every
chiral derivation is inner ($\beta \mapsto \beta + \epsilon$
is generated by $(\epsilon\gamma)_{(0)}$).
\begin{equation}\label{eq:bg-hh1}
\HH^1 = 0.
\end{equation}
This inverts the Heisenberg pattern: a double pole makes
inner derivations trivial ($\HH^1 = \CC$); a simple pole
makes all derivations inner ($\HH^1 = 0$). The dichotomy
is a shadow of the Koszul duality
$(\beta\gamma)^! \simeq bc$.

\medskip
\noindent
\textbf{Arity $2$: the unipotent KZ connection.}
The KZ connection~\eqref{eq:kz-bg} at arity $2$ is
\begin{equation}\label{eq:kz-bg-2}
\nabla_{\mathrm{KZ}}^{\beta\gamma}
= d + \Theta_{12}\,\omega_{12},
\qquad \omega_{12} = d\log(z_1 - z_2).
\end{equation}
The solution is $\Phi(w) = I + \Theta_{12}\,\log w$
(equation~\eqref{eq:bg-kz-sol-2}), with
\emph{unipotent} monodromy
\begin{equation}\label{eq:bg-monodromy}
M = R_{\beta\gamma} = I + 2\pi i\,\Theta_{12},
\qquad (M - I)^2 = 0.
\end{equation}
For a unipotent local system ($N = M - I$, $N^2 = 0$),
the de~Rham cohomology is $H^0 = \ker(N)$,
$H^1 = \mathrm{coker}(N)$.
Since $\Theta_{12}$ has rank $2$ on $V^{\otimes 2} \cong \CC^4$
(killing $\beta \otimes \beta$ and $\gamma \otimes \gamma$),
$\dim \ker = \dim \mathrm{coker} = 2$.
The monodromy is scalar for the Heisenberg
($e^{2\pi ik}$, generic), semisimple for affine KM
($e^{2\pi i\lambda_J/(k+2)}$, channel-dependent), and
unipotent for $\beta\gamma$ (always nontrivial flat
sections).

The unique deformation cocycle
$\phi(\beta \otimes \gamma)
= \epsilon \cdot \mathbf{1}$ represents
\begin{equation}\label{eq:bg-hh2}
[\phi] \in \HH^2 = \CC,
\end{equation}
corresponding to $\beta(z)\gamma(w) \sim (1 + \epsilon)/(z - w)$.

\medskip
\noindent
\textbf{Arity $n \geq 3$: vanishing.}
By the Koszul resolution of length~$2$:
\begin{equation}\label{eq:bg-higher-vanish}
\HH^{n,\mathrm{ord}}_{\mathrm{ch}}(\beta\gamma, \beta\gamma) = 0
\quad \text{for } n \geq 3.
\end{equation}

\medskip
\noindent
\textbf{Summary on $D$.}
The ordered chiral Hochschild cohomology on the formal disk is
\begin{equation}\label{eq:bg-hh-summary}
\HH^{n,\mathrm{ord}}_{\mathrm{ch}}(\beta\gamma, \beta\gamma) =
\begin{cases}
\CC & n = 0 \;\text{(vacuum: centre $= \CC \cdot \mathbf{1}$)}, \\
0 & n = 1 \;\text{(all derivations inner: simple-pole OPE)}, \\
\CC & n = 2 \;\text{(OPE residue deformation:
  $1/(z{-}w) \mapsto (1{+}\epsilon)/(z{-}w)$)}, \\
0 & n \geq 3 \;\text{(Koszul resolution of length $2$)},
\end{cases}
\end{equation}
with Poincar\'e polynomial $P_{\beta\gamma}(t) = 1 + t^2$.
The total dimension is $2$, and the cohomological amplitude
is $[0, 2]$, consistent with Theorem~H.
The derived chiral centre is
$Z^{\mathrm{der}}_{\mathrm{ch}}(\beta\gamma) \cong
\CC \oplus \CC[-2]$
as a graded vector space.

\begin{remark}[Poincar\'e polynomials and Koszul duality]
\label{rem:bg-poincare-comparison}
$P_{\cH_k}(t) = 1 + t + t^2$ versus
$P_{\beta\gamma}(t) = 1 + t^2$: the missing $t$-term
reflects $\HH^1 = 0$. By Koszul duality
$(\beta\gamma)^! \simeq bc$, the $bc$ ghost system has
$P_{bc}(t) = 1 + t^2$ as well
($\HH^n(\beta\gamma) \cong \HH^{2-n}(bc)^*$).
\end{remark}

\medskip
\noindent
\textbf{The computation on $\PP^1$.}
The computation on $\PP^1$ gives the same result
(as for the Heisenberg, the factorisation $\cD$-module is
insensitive to the compactification for chirally Koszul
algebras):
\begin{equation}\label{eq:bg-hh-P1}
\HH^{n,\mathrm{ord}}_{\mathrm{ch,\PP^1}}(\beta\gamma,
\beta\gamma)
\cong
\HH^{n,\mathrm{ord}}_{\mathrm{ch},D}(\beta\gamma,
\beta\gamma)
\quad \text{for all } n.
\end{equation}

\medskip
\noindent
\textbf{Koszul duality and complementarity.}
The Koszul dual is the $bc$ ghost system at weight $\lambda$,
with $c_{bc}(\lambda) = 1 - 3(2\lambda - 1)^2$, satisfying
the complementarity
\begin{equation}\label{eq:bg-bc-compl}
  \kappa(\beta\gamma) + \kappa(bc) = 0.
\end{equation}
(Verification: $c_{\beta\gamma}(2) = 26$,
$c_{bc}(2) = -26$, sum $= 0$.)
The complementarity $K = 0$ places $\beta\gamma$ in the same
family as the Heisenberg and lattice algebras: the Koszul
conductor vanishes.

\medskip
\noindent
\textbf{The $\Ethree$ structure.}
As for the Heisenberg, the $\Ethree$ obstruction space
$\HH^3_{\Etwo}(\BarSig(\beta\gamma), \BarSig(\beta\gamma))
\cong \Lambda^3(\CC^2) = 0$
vanishes for dimensional reasons, and the $\Ethree$-algebra
structure on
$Z^{\mathrm{der}}_{\mathrm{ch}}(\beta\gamma)
\cong \CC \oplus \CC[-2]$
is the unique (trivial) one.

\medskip
\noindent
\textbf{Comparison with topological Hochschild.}
By the formality bridge (Theorem~\ref{thm:formality-bridge}),
$\HH^{*,\mathrm{ord}}_{\mathrm{ch}}(\beta\gamma, \beta\gamma)
\simeq \HH^{*,\mathrm{top}}(\beta\gamma, \beta\gamma)$.
The transferred $\Ainf$-structure carries
nontrivial operations $m_3^{\mathrm{tr}} \neq 0$ and
$m_4^{\mathrm{tr}} \neq 0$ (shadow depth $r = 4$): the
formality bridge preserves the chain complex but does
not kill the higher operations.
The information lost in passing from ordered to topological
is the explicit $z$-dependence of the flat sections:
$\Phi(w) = I + \Theta\,\log w$ is replaced by the abstract
homotopy transfer data.
\end{example}

\begin{example}[Yangian $Y_\hbar(\mathfrak{sl}_2)$
as a chiral quantum group: three structures and the
ordered chiral center]
\label{ex:yangian}
The Yangian $Y_\hbar(\mathfrak{sl}_2)$ is the first example
in which the ordered chiral homology differs from the
Beilinson--Drinfeld (symmetric) chiral homology. It is a
genuinely $\Eone$-chiral algebra: the $R$-matrix is
independent input, not derived from any $\Einf$-chiral OPE.
The formality bridge (Theorem~\ref{thm:formality-bridge})
does not apply, because the factorisation $\cD$-module is not
$\Sigma_n$-equivariant. We exhibit all three structures of
the chiral quantum group equivalence
(Theorem~\ref{thm:chiral-qg-equiv}) concretely, then compute
the ordered chiral center at low arities and identify
$\ker(\av) \neq 0$.

\medskip
\noindent
\textbf{Data.} Let $\fg = \mathfrak{sl}_2$ with standard
basis $\{e, f, h\}$, brackets $[e,f] = h$, $[h,e] = 2e$,
$[h,f] = -2f$, and normalised Killing form $(e,f) = 1$,
$(h,h) = 2$. The dual Coxeter number is $h^\vee = 2$.
The deformation parameter is
$\hbar = 1/(k + h^\vee) = 1/(k + 2)$
($= h_{\mathrm{KZ}}$ of
Remark~\ref{rem:hbar-convention}), where $k$ is the
level of the affine input $V_k(\mathfrak{sl}_2)$.
The inverse Casimir is
\begin{equation}\label{eq:yangian-casimir}
  \Omega = e \otimes f + f \otimes e
  + \tfrac{1}{2}\,h \otimes h
  \;\in\; \fg \otimes \fg.
\end{equation}
In the fundamental representation $V = \CC^2$:
$(\rho \otimes \rho)(\Omega) = P - \tfrac{1}{2}\,I$,
where $P$ is the permutation operator on $V \otimes V$.

\medskip
\noindent
\textbf{Structure~(I): the $R$-matrix.}
The Yang $R$-matrix on $V \otimes V$ is
\begin{equation}\label{eq:yangian-rmatrix}
  R(u) = u\,I + \hbar\,P
  \;\in\; \End(V \otimes V)[u],
\end{equation}
where $u$ is the spectral parameter (the difference of
evaluation points on $\CC$). This satisfies the quantum
Yang--Baxter equation
\begin{equation}\label{eq:yangian-qybe}
  R_{12}(u - v)\,R_{13}(u)\,R_{23}(v)
  = R_{23}(v)\,R_{13}(u)\,R_{12}(u - v)
\end{equation}
in $\End(V^{\otimes 3})[u, v]$, and the unitarity condition
$R_{12}(u)\,R_{21}(-u) = (u^2 - \hbar^2)\,I$
(scalar: invertible for $u \neq \pm\hbar$).

At the classical level, $R(u) = u\,I + \hbar\,P
= u(I + \hbar\,\Omega/u + O(\hbar^2))$
after restoring the Casimir from $P = \Omega + \tfrac{1}{2}I$
on $V \otimes V$. The classical $r$-matrix
$r(z) = \hbar\,\Omega/z$ is in KZ normalisation;
% AP126 check: at hbar=0, i.e. k -> infty, r = 0.
% At critical level k = -2, hbar diverges. Verified.
at $\hbar = 0$ (i.e.\ $k \to \infty$), $r = 0$;
at the critical level $k = -2$, $\hbar$ diverges.

The eigenvalues of $R(u)$ on $V \otimes V$ are
$u + \hbar$ on $\operatorname{Sym}^2 V$ (dimension $3$)
and $u - \hbar$ on $\bigwedge^2 V$ (dimension $1$);
the $R$-matrix has a zero on $\bigwedge^2 V$ at $u = \hbar$,
selecting the antisymmetric channel.

\medskip
\noindent
\textbf{Structure~(II): the $\Ainf$ structure.}
The chiral $\Ainf$-structure maps
$m_k^{\mathrm{ch}} \colon Y_\hbar^{\otimes k} \to Y_\hbar$
are the operations extracted from the ordered bar complex via
the bar-cobar correspondence. The transfer matrix
\begin{equation}\label{eq:yangian-transfer}
  T(u) =
  \begin{pmatrix} A(u) & B(u) \\ C(u) & D(u) \end{pmatrix}
  \;\in\; \End(V) \otimes Y_\hbar(\mathfrak{sl}_2)[\![u^{-1}]\!]
\end{equation}
encodes the generating functions of the algebra.

The binary operation $m_2^{\mathrm{ch}}$ is the RTT
relation: for any two entries $T_{ij}(u)$ and $T_{kl}(v)$
of the transfer matrix,
\begin{equation}\label{eq:yangian-rtt}
  R_{12}(u - v)\,T_1(u)\,T_2(v)
  = T_2(v)\,T_1(u)\,R_{12}(u - v),
\end{equation}
where subscripts indicate the auxiliary space factor.
Expanding this $4 \times 4$ matrix equation produces the
Yangian relations in RTT form. The ten independent
component relations include:
\begin{itemize}
\item diagonal commutativity: $[A(u), A(v)] = 0$,
  $[D(u), D(v)] = 0$;
\item off-diagonal commutativity: $[B(u), B(v)] = 0$,
  $[C(u), C(v)] = 0$;
\item the cross-relations, e.g.\
  $(u{-}v)[B(u), C(v)]
  = \hbar\bigl(D(u)A(v) - D(v)A(u)\bigr)$.
\end{itemize}

The higher operations $m_k^{\mathrm{ch}}$ for $k \geq 3$
involve the Drinfeld associator $\Phi_{\mathrm{KZ}}$. The
binary product $m_2^{\mathrm{ch}}$ is strictly associative
only up to the holonomy of the KZ connection on
$\overline{\FM}_3^{\mathrm{ord}}(\CC)$:
the correction is $m_3^{\mathrm{ch}}(a, b, c)
= \Phi_{\mathrm{KZ}} \cdot m_2(m_2(a,b), c)
- m_2(a, m_2(b,c))$, where $\Phi_{\mathrm{KZ}}$ is the
regularised holonomy between the two maximal boundary
strata of $\overline{\FM}_3^{\mathrm{ord}}(\CC)$.

The Gauss decomposition
$T(u) = F(u)\,K(u)\,E(u)$ (lower-triangular, diagonal,
upper-triangular) translates the RTT relations into
the Drinfeld presentation: generators
$\{E_r, F_r, H_r\}_{r \geq 0}$ with relations
\begin{align}
  [H_r, H_s] &= 0, \label{eq:yangian-D1}\\
  [H_0, E_s] &= 2E_s, \quad [H_0, F_s] = -2F_s,
  \label{eq:yangian-D2}\\
  [E_r, F_s] &= H_{r+s}, \label{eq:yangian-D4}\\
  [H_{r+1}, E_s] - [H_r, E_{s+1}]
  &= \hbar(H_r E_s + E_s H_r),
  \label{eq:yangian-D5}
\end{align}
and the Serre relation
$[E_{r+1}, [E_r, E_s]] - [E_r, [E_{r+1}, E_s]] = 0$
(and similarly for $F$). The leading modes $E_0 = e$,
$F_0 = f$, $H_0 = h$ generate a copy of
$U(\mathfrak{sl}_2)$ inside $Y_\hbar(\mathfrak{sl}_2)$;
the PBW filtration has
$\operatorname{gr} Y_\hbar(\mathfrak{sl}_2)
\cong U(\mathfrak{sl}_2[t])$.

\medskip
\noindent
\textbf{Structure~(III): the chiral coproduct.}
The Drinfeld coproduct on $Y_\hbar(\mathfrak{sl}_2)$ is
\begin{equation}\label{eq:yangian-drinfeld-coprod}
  \Delta_z(T(u))
  = T(u) \;\dot{\otimes}\; T(u - z),
\end{equation}
where $\dot{\otimes}$ denotes matrix multiplication in the
auxiliary space $\End(V)$ and ordinary (completed) tensor
product in the $Y_\hbar$ factor. The spectral parameter $z$
is the difference coordinate $z_1 - z_2$ on $\Conf_2(\CC)$.
On Drinfeld generators:
\begin{align}
  \Delta_z(E_0) &= E_0 \otimes 1
  + 1 \otimes E_0
  + \hbar\,H_0 \otimes E_0 \cdot z^{-1}
  + O(z^{-2}), \label{eq:yangian-coprod-E}\\
  \Delta_z(H_0) &= H_0 \otimes 1
  + 1 \otimes H_0
  + 2\hbar\,(E_0 \otimes F_0
  - F_0 \otimes E_0) \cdot z^{-1}
  + O(z^{-2}).
  \label{eq:yangian-coprod-H}
\end{align}

The coproduct is coassociative up to the Drinfeld
associator $\Phi$:
\[
  (\Delta_{z_1} \otimes \id) \circ \Delta_{z_2}
  = \Phi_{123}(z_1, z_2) \cdot
  (\id \otimes \Delta_{z_2}) \circ \Delta_{z_1}
  \cdot \Phi_{123}(z_1, z_2)^{-1},
\]
where $\Phi_{123}(z_1, z_2) \in
\exp(\widehat{\Lie}(\Omega_{12}, \Omega_{23}))$ is
the KZ associator with spectral parameters.

The $R$-matrix is recovered from the coproduct:
$R(u) = \sigma \circ \Delta_z \circ (\Delta_z^{\op})^{-1}$
evaluated on $V \otimes V$, where $\sigma$ is the flip.
This is the content of the equivalence
(I)$\Leftrightarrow$(III) in
Theorem~\ref{thm:chiral-qg-equiv}.

\medskip
\noindent
\textbf{The three structures are equivalent.}
The ordered bar complex $\Barord(Y_\hbar(\mathfrak{sl}_2))
= T^c(s^{-1}\overline{Y}_\hbar)$ (desuspended augmentation
ideal, deconcatenation coproduct) is the universal datum
from which all three structures are recovered:
\begin{enumerate}[label=\textup{(\alph*)}]
\item The bar differential encodes the $\Ainf$ maps
  (coderivations of $T^c$ biject with multilinear
  operations): the RTT relation~\eqref{eq:yangian-rtt} is
  the degree-$2$ bar differential, and the Serre relations
  are the degree-$3$ component.
\item The deconcatenation coproduct on $T^c$ induces the
  Drinfeld coproduct~\eqref{eq:yangian-drinfeld-coprod}
  via cobar dualisation.
\item The $R$-matrix~\eqref{eq:yangian-rmatrix} is the
  arity-$2$ KZ monodromy: the parallel transport of the
  connection $\nabla = d + \hbar\,\Omega/(z_1 - z_2)
  \cdot d(z_1 - z_2)$ around the diagonal.
\end{enumerate}
The passage
$r(z) = \hbar\,\Omega/z
\;\xrightarrow{\text{exponentiation}}\;
R(u) = u\,I + P
\;\xrightarrow{\text{RTT}}\;
\text{Yangian relations}
\;\xrightarrow{\text{Gauss}}\;
\text{Drinfeld presentation}$
is the ordered bar construction applied to
$V_k(\mathfrak{sl}_2)$, producing
$Y_\hbar(\mathfrak{sl}_2)$ as the $\Eone$-chiral Koszul
dual.

\medskip
\noindent
\textbf{Part B: the ordered chiral center.}
We compute $\HH^{*,\mathrm{ord}}_{\mathrm{ch}}
(Y_\hbar(\mathfrak{sl}_2),\,Y_\hbar(\mathfrak{sl}_2))$
at low arities on the formal disk $D$, and identify
the kernel of the averaging map $\av$.

\smallskip
\noindent\emph{Arity~$0$: the center.}
The arity-$0$ cohomology is the chiral center
$Z(Y_\hbar(\mathfrak{sl}_2))$: those elements $a$ with
$a_{(n)}b = 0$ for all $b$ and $n \geq 0$.
The quantum determinant
$\operatorname{qdet}\,T(u) = A(u)\,D(u - \hbar)
- B(u)\,C(u - \hbar)$
is central in $Y_\hbar(\mathfrak{sl}_2)$~\cite{Molev07}.
The center is the polynomial algebra generated by the
coefficients of $\operatorname{qdet}\,T(u)$:
\begin{equation}\label{eq:yangian-hh0}
  \HH^{0,\mathrm{ord}}_{\mathrm{ch}}
  = Z(Y_\hbar(\mathfrak{sl}_2))
  = \CC[\operatorname{qdet}\,T(u)].
\end{equation}
The averaging map $\av_0$ is an isomorphism at arity~$0$
(there is nothing to symmetrise), so
$\ker(\av_0) = 0$.

\smallskip
\noindent\emph{Arity~$2$: the KZ local system and $\ker(\av_2) \neq 0$.}
This is the arity where the ordered theory diverges from
the symmetric. The ordered chiral Hochschild complex at
arity~$2$ involves the de~Rham cohomology of
$\Conf_2^{\mathrm{ord}}(\CC)$ with coefficients in the KZ
local system. At arity~$2$, the leading-order KZ connection
$\nabla = d + \hbar\,\Omega\,d\log(z_1 - z_2)$
captures the classical $r$-matrix data
$r(z) = \hbar\,\Omega/z$. The full quantum $R$-matrix
$R(u) = u\,I + \hbar\,P$ contributes higher-order
corrections visible at arity $\geq 3$ through the Drinfeld
associator; at arity~$2$ the connection is:
\begin{equation}\label{eq:yangian-kz-arity2}
  \nabla_{\mathrm{KZ}}
  = d + \frac{\hbar\,\Omega}{z_1 - z_2}\,d(z_1 - z_2)
  = d + \hbar\,\Omega\,d\log(z_1 - z_2),
\end{equation}
with coefficients in
$\End(Y_\hbar^{\otimes 2})$.
Setting $w = z_1 - z_2$, the flat sections on $\CC^\times$
satisfy $\partial \Phi / \partial w
= (\hbar\,\Omega/w)\,\Phi$.

For the Yangian on the fundamental representation
$V \otimes V$, the Casimir decomposes along the
irreducible $\mathfrak{sl}_2$-components:
\[
  V \otimes V = \operatorname{Sym}^2 V \oplus \bigwedge^2 V,
  \qquad
  \Omega\big|_{\operatorname{Sym}^2 V} = +\tfrac{1}{2},
  \quad
  \Omega\big|_{\bigwedge^2 V} = -\tfrac{3}{2}.
\]
The monodromy of $\nabla_{\mathrm{KZ}}$ around $w = 0$
on each isospin component is
\begin{equation}\label{eq:yangian-monodromy}
  M_{\mathrm{sym}} = e^{2\pi i \hbar/2}
  = e^{\pi i/(k+2)}, \qquad
  M_{\mathrm{alt}} = e^{-3\pi i/(k+2)}.
\end{equation}

The averaging map at arity~$2$ is
$\av_2 \colon H^*_{\mathrm{dR}}(\Conf_2^{\mathrm{ord}},
\nabla_{\mathrm{KZ}}) \to
H^*_{\mathrm{dR}}(\Conf_2,
\nabla_{\mathrm{KZ}})^{\Sigma_2}$,
where $\Sigma_2$ acts by the flip $w \mapsto -w$
composed with the $R$-matrix.

For an $\Einf$-chiral algebra, $R = \id$ and the flip acts
by the Koszul sign; $\av_2$ is an isomorphism and
$\ker(\av_2) = 0$. For the Yangian, $R(u) \neq \id$: the
$R$-matrix mixes the $\operatorname{Sym}^2 V$ and
$\bigwedge^2 V$ channels non-trivially, and the
$\Sigma_2$-action on the ordered de~Rham cohomology is the
composition of the geometric flip with $R$.

The kernel is spanned by those flat sections that are
antisymmetric under this twisted $\Sigma_2$-action:
\begin{equation}\label{eq:yangian-ker-av2}
  \ker(\av_2)
  \cong \bigwedge\nolimits^2 V = \CC.
\end{equation}
The generator is the flat section
$\Phi_{\mathrm{alt}}(w) = w^{-3\hbar/2}$
on the antisymmetric channel of $V \otimes V$. The full
$R$-matrix $R(u) = u\,I + \hbar\,P$ evaluates on the
antisymmetric component as $(u - \hbar)\,\id$; the
$\Sigma_2$-flip sends $\Phi_{\mathrm{alt}}(w)$ to
$R_{21}(-w)\,\Phi_{\mathrm{alt}}(-w)
= -\Phi_{\mathrm{alt}}(w)$
(the $R$-matrix eigenvalue $-1$ on $\bigwedge^2 V$ at the
leading order cancels the geometric sign).
This class is in the ordered cohomology but maps to zero
under $\av_2$: it is invisible to the symmetric chiral
homology.

The scalar shadow surviving $\av_2$ is
% AP1: kappa(V_k(sl_2)) = dim(sl_2)(k+h^v)/(2h^v) = 3(k+2)/4.
% k=0 -> 3/2. k=-2 -> 0. Verified.
$\kappa = \av(r(z)) + \dim(\fg)/2
= 3k/(2 \cdot 2) + 3/2
= 3(k + 2)/4$.

\smallskip
\noindent\emph{Arity~$3$: the Drinfeld associator in $\ker(\av_3)$.}
At arity~$3$, the KZ system on $\Conf_3^{\mathrm{ord}}(\CC)$
is
\[
  d\Phi
  = \hbar\biggl(
  \frac{\Omega_{12}}{z_1 - z_2}\,d(z_1 - z_2)
  + \frac{\Omega_{13}}{z_1 - z_3}\,d(z_1 - z_3)
  + \frac{\Omega_{23}}{z_2 - z_3}\,d(z_2 - z_3)
  \biggr)\Phi.
\]
The fundamental group $\pi_1(\Conf_3^{\mathrm{ord}}(\CC))
= PB_3$ (pure braid group on three strands). The KZ
monodromy gives a representation
$\rho_{\mathrm{KZ}} \colon PB_3
\to \operatorname{GL}(V^{\otimes 3})$.

The Drinfeld associator $\Phi_{\mathrm{KZ}}
\in \exp(\widehat{\Lie}(t_{12}, t_{23}))$
is the regularised holonomy of this connection along the
real path from the boundary stratum
$|z_1 - z_2| \ll |z_2 - z_3|$
to the stratum
$|z_2 - z_3| \ll |z_1 - z_2|$
in $\overline{\FM}_3^{\mathrm{ord}}(\CC)$.

The averaging map $\av_3$ sends ordered data to
$\Sigma_3$-coinvariants. At arity~$3$, the relevant
representation is the adjoint $\fg = \mathfrak{sl}_2$
(dimension $3$), not the fundamental $V = \CC^2$
(dimension $2$) used at arity~$2$: the convolution
algebra acts on $\fg^{\otimes n}$ via the adjoint action
(the Casimir operators $\Omega_{ij}$ act through the
adjoint representation on each tensor factor),
and the Drinfeld associator lives in $\Lambda^3(\fg)$.
The kernel of $\av_3$ contains the associator:
\begin{equation}\label{eq:yangian-ker-av3}
  \ker(\av_3)
  \cong \Lambda^3(\fg) = \CC,
\end{equation}
one-dimensional, carrying the associator class. This
computation follows from the structure of the cohomology
of $\Conf_3^{\mathrm{ord}}(\CC)$: the top-degree Arnold
class $\omega_{12} \wedge \omega_{23}
\in H^2(\Conf_3^{\mathrm{ord}})$ maps to the alternating
representation of $\Sigma_3$, and the KZ-twisted version
(with coefficients in $\fg^{\otimes 3}$) gives
$\ker(\av_3) \cong \Lambda^3(\fg)$.
For $\fg = \mathfrak{sl}_2$:
$\dim \Lambda^3(\fg) = \binom{3}{3} = 1$.

For $\Einf$-chiral algebras, the formality bridge forces
$\ker(\av_3) = 0$: the associator is in the image of a
$\Sigma_3$-equivariant quasi-isomorphism and carries no
independent ordered data. For the Yangian, the
$\cD$-module is not $\Sigma_3$-equivariant (because
$R \neq \id$), and the associator survives as a non-trivial
element of the ordered chiral center invisible to any
symmetric theory.

\smallskip
\noindent\emph{Explicit computation in $\ker(\av_3)$.}
On the fundamental representation $V = \CC^2$ of
$\mathfrak{sl}_2$, the KZ system at three points reduces
to a $2 \times 2$ Fuchsian system on the multiplicity
space $M = \Hom_{\mathfrak{sl}_2}(V_0 \oplus V_2,\,
V^{\otimes 3})$, where $V_j$ is the spin-$j/2$
irreducible. The Casimir $\Omega_{12}$ acts on $M$ as
$A = \operatorname{diag}(-\tfrac{3}{2},\,\tfrac{1}{2})$
(eigenvalue $-\tfrac{3}{2}$ on the $j=0$ channel,
$\tfrac{1}{2}$ on $j=1$).
The representation of $\Sigma_3$ on $M$ has character
$(2,0,0,0,-1,-1)$: it is the standard representation
with no trivial summand. The Reynolds operator
$\tfrac{1}{6}\sum_{\sigma \in \Sigma_3} \sigma$ therefore
vanishes identically on $M$, so the \emph{entire}
$2$-dimensional multiplicity space lies in $\ker(\av_3)$.
The Drinfeld associator acts on $M$ as a pure $\Eone$
invariant with no symmetric shadow.

At level $k = 2$ ($\kappa = 3(k{+}2)/4 = 3$), direct
numerical integration of the Fuchsian system gives
\[
  \Phi_{\mathrm{KZ}}\big|_M
  \;=\;
  \begin{pmatrix}
    0.884 & -0.204 \\
    0.306 & \phantom{-}1.061
  \end{pmatrix},
  \qquad
  \det \Phi_{\mathrm{KZ}}\big|_M = 1,
\]
with eigenvalues on the unit circle (the connection is
unitary on $M$). The off-diagonal entry
$|\Phi_{12}| = 0.204$ measures the channel-mixing
amplitude between the $j = 0$ and $j = 1$ sectors; it
decays as $2.849/\kappa^2$ at large $\kappa$, consistent
with the semiclassical expansion
\[
  \Phi_{\mathrm{KZ}}\big|_M
  \;=\; I
  \;-\; \frac{\zeta(2)}{\kappa^2}\,[t_{12},\, t_{23}]
  \;+\; O(1/\kappa^3),
\]
where $t_{ij} = \Omega_{ij}/(k + h^\vee)$ and
$[t_{12}, t_{23}] \neq 0$ on $M$ (reflecting the
non-abelian Lie bracket). This is an explicit non-trivial
class in $\ker(\av_3)$: an ordered invariant of the
Yangian that is invisible to every symmetric
chiral-homological observable.

\smallskip
\noindent\emph{Arity~$n$: braid group monodromy.}
At arity~$n$, the ordered chiral center carries the full
pure braid group monodromy of the KZ system with the
quantum $R$-matrix:
\[
  \rho_{\mathrm{KZ}}^{(n)} \colon PB_n
  \;\to\; \operatorname{GL}\bigl(V^{\otimes n}\bigr).
\]
The symmetric shadow at arity~$n$ retains only the scalar
invariants: the traces of the monodromy matrices and the
Markov-trace data. The kernel $\ker(\av_n)$ carries the
full matrix-valued monodromy, from which knot invariants
(the Jones polynomial and its coloured generalisations)
are extracted.

The Poincar\'e series of the ordered chiral center has the
form
\[
  P^{\mathrm{ord}}_{Y_\hbar}(t)
  = P^{\Sigma}_{Y_\hbar}(t) + \Delta P(t),
\]
where $\Delta P(t) = t^2 + t^3 + \cdots$ records the extra
data in $\ker(\av)$ at each arity. At arity $2$:
$\Delta P_2 = 1$ (the antisymmetric KZ block). At arity $3$:
$\Delta P_3 = 1$ (the Drinfeld associator class).

\medskip
\noindent
\textbf{Why this is new.}
This is the first instance where the formality bridge
(Theorem~\ref{thm:formality-bridge}) fails: the
$\cD$-module is not $\Sigma_n$-equivariant because
$R(u) \neq \id$, so the ordered chiral homology carries
strictly more data than the symmetric theory.

\begin{center}
\renewcommand{\arraystretch}{1.3}
\begin{tabular}{@{} c l l l @{}}
\toprule
Arity $n$ & Ordered datum & $\ker(\av_n)$
  & Symmetric shadow \\
\midrule
$0$ & $Z(Y_\hbar) = \CC[\operatorname{qdet}]$ & $0$
  & $Z(Y_\hbar)$ \\
$2$ & full $R$-matrix $R(u) = uI + \hbar P$
  & $\bigwedge^2 V = \CC$
  & scalar $\kappa = 3(k{+}2)/4$ \\
$3$ & KZ associator $\Phi_{\mathrm{KZ}}$
  & $\Lambda^3(\fg) = \CC$
  & trivial \\
$n$ & $PB_n$-monodromy of KZ
  & matrix-valued
  & scalar traces only \\
\bottomrule
\end{tabular}
\end{center}

\end{example}

\begin{example}[Virasoro: class $M$, $r = \infty$;
the irregular $\cD$-module and Stokes phenomenon]
\label{ex:virasoro}
The Virasoro algebra $\mathrm{Vir}_c$ has a single generating
field $T(z)$ of conformal weight~$2$, OPE
\[
  T(z)\,T(w) \;\sim\;
  \frac{c/2}{(z-w)^4} + \frac{2\,T(w)}{(z-w)^2}
  + \frac{\partial T(w)}{z-w},
\]
% AP19: OPE pole order 4; d-log absorbs one pole; r-matrix pole order 3.
% C11: r^Vir(z) = (c/2)/z^3 + 2T/z.
and $r$-matrix $r^{\mathrm{Vir}}(z) = (c/2)/z^3 + 2T/z$
(quartic OPE pole, $d$-log absorption gives cubic
$r$-matrix pole).
% C2: kappa(Vir) = c/2. c=0 -> 0; c=13 -> 13/2 self-dual.
The Koszul invariant is $\kappa = c/2$.
The algebra is $\Einf$-chiral (symmetric OPE: the commutator
$[T(z), T(w)]$ is commutative after clearing poles)
and class $M$ (shadow depth $r = \infty$: all higher shadow
operations $\ell_k \neq 0$ in the symmetric convolution
algebra~$\gmod$).

\medskip
\noindent\textbf{Formality bridge and $\Einf$ descent.}
By the formality bridge (Theorem~\ref{thm:formality-bridge}),
ordered and symmetric chiral homologies agree;
$\ker(\av) = 0$ at each arity. The infinite shadow tower of
class $M$ is a \emph{symmetric} phenomenon. The content
genuinely new to class $M$ is the \emph{singularity type}
of the factorisation $\cD$-module: irregular, with Stokes
data that has no analogue in class $G/L/C$.

\medskip
\noindent\textbf{The irregular KZ connection at arity $2$.}
At arity~$2$, setting $z = z_1 - z_2$, the KZ equation
(Proposition~\ref{prop:fact-dmod-props}(ii)) reads
\begin{equation}\label{eq:kz-vir-arity2}
  \frac{\partial\Phi}{\partial z}
  = \biggl(\frac{c/2}{z^3} + \frac{2T}{z}\biggr)\Phi.
\end{equation}
This has a pole of order~$3$ in $z$: the singular point
$z = 0$ is \emph{irregular} of Poincar\'e rank~$2$
(pole order minus one). The classical Riemann--Hilbert
correspondence of Deligne does not apply; one needs the
irregular Riemann--Hilbert correspondence of
D'Agnolo--Kashiwara~\cite{DAgnoloKashiwara16}.

For comparison, the affine Kac--Moody KZ equation at arity~$2$
(Example~\ref{ex:km}) is
$\partial\Phi/\partial z = (\Omega/((k + h^\vee)z))\,\Phi$,
a \emph{simple} pole. Its solution is the power function
$\Phi(z) = z^{\Omega/(k+h^\vee)}$:
purely algebraic monodromy, no exponential corrections,
no Stokes data. The difference between class~$L$ and
class~$M$ is the difference between regular and irregular
singularities.

\medskip
\noindent\textbf{Formal solution and the exponential factor.}
On the primary sector, $T$ acts by the eigenvalue $h$
(conformal weight of the field being transported).
Equation~\eqref{eq:kz-vir-arity2} restricts to the scalar ODE
\begin{equation}\label{eq:kz-vir-scalar}
  \frac{d\varphi}{dz}
  = \biggl(\frac{c/2}{z^3} + \frac{2h}{z}\biggr)\varphi,
\end{equation}
which integrates exactly:
\begin{equation}\label{eq:kz-vir-solution}
  \varphi(z)
  = z^{2h}\,\exp\!\biggl(-\frac{c}{4z^2}\biggr).
\end{equation}
The factor $z^{2h}$ is the monodromy piece: it contributes
$\exp(4\pi i h)$ upon circling $z = 0$, matching the
conformal spin. The factor $\exp(-c/(4z^2))$ is the
\emph{irregular} piece: it has an essential singularity
at $z = 0$, invisible to monodromy, carrying exponentially
small corrections in different angular sectors of the
$z$-plane.

In the Hukuhara--Turrittin--Levelt normal
form~\cite{Wasow65,Sabbah13},
the formal fundamental matrix is
$\Phi^{\mathrm{form}}(z)
= \widehat{F}(z)\,z^{L}\,
  \exp(-c/(4z^2)\,\mathrm{Id})$,
where $\widehat{F}(z) \in \mathrm{GL}(V)[[z^{-1}]]$
is a Gevrey-$1/2$ formal gauge transformation.

\medskip
\noindent\textbf{Stokes phenomenon.}
The exponential $\exp(-c/(4z^2))$ produces four Stokes rays
at $\theta = 0, \pi/2, \pi, 3\pi/2$ (matching Poincar\'e
rank $2$: $2 \times 2 = 4$ rays).
Across each ray, the true solution jumps by a Stokes
multiplier $\mathfrak{S}$:
\begin{equation}\label{eq:stokes-vir}
  \Phi_+^{\mathrm{true}}(z)
  = \Phi_-^{\mathrm{true}}(z)
  \cdot
  (\mathrm{Id} + \mathfrak{S}).
\end{equation}
The flat sections are classified by a
\emph{Stokes-filtered local system}, not a monodromy
representation. This contrasts with class $G/L/C$, where
regular singularities produce no Stokes data.

\medskip
\noindent\textbf{Wall-crossing at $c = -22/5$.}
The Stokes multipliers depend on the central charge~$c$.
At $c = -22/5$ (the $(2,5)$ minimal model, the Yang--Lee
edge singularity), two Stokes rays coalesce in a
wall-crossing: the connection resonates, and the Stokes
multipliers undergo a discontinuous change. In the language
of the shadow obstruction tower~\cite{Lorgat26I},
the shadow invariant $S_4 = 10/(c(5c+22))$ has a pole
at $c = -22/5$, signalling the coalescence.

\medskip
\noindent\textbf{De Rham cohomology of the irregular
$\cD$-module.}
The factorisation $\cD$-module
$\cF_n^{\mathrm{ord}}(\mathrm{Vir}_c)$ is holonomic on
$X^n$ for any smooth proper curve~$X$
(Proposition~\ref{prop:fact-dmod-props}(i)).
By Mebkhout's theorem~\cite{Mebkhout84}, holonomic
$\cD$-modules on proper varieties have finite-dimensional
de~Rham cohomology regardless of regularity.
The finite-dimensionality of
$H^*_{\mathrm{dR}}(X^n, \cF_n^{\mathrm{ord}})$ is
therefore unconditional: the irregular singularities of
class~$M$ do not obstruct finiteness.

The de~Rham cohomology computes the flat sections of the
connection with coefficients in the
\emph{Stokes-filtered local system}, not the naive local
system (which is the wrong coefficient system for irregular
connections). For $X = \PP^1$ at arity $n = 2$: the
irregular de~Rham space
$H^*_{\mathrm{dR}}(\Conf_2(\PP^1), \nabla_{\mathrm{KZ}}^{\mathrm{Vir}})$
is finite-dimensional, with the dimension depending on $c$
through the Stokes multipliers.

\medskip
\noindent\textbf{Shadow tower and the gravitational
partition function.}
The class $M$ shadow obstruction tower is infinite:
\begin{align*}
  S_2 &= c/2 = \kappa,\\
  S_3 &= 2\quad\text{(}c\text{-independent)},\\
  S_4 &= 10/(c(5c+22)),\\
  S_5 &= -48/(c^2(5c+22)),\\
  S_6 &= 80(45c+193)/(3\,c^3(5c+22)^2),
\end{align*}
with all $S_r$ belonging to $\mathbb{Q}(c)$ and having
poles only at $c = 0$ and $c = -22/5$.
The spectral discriminant is $\Delta = 8\kappa S_4
= 40/(5c+22) \neq 0$: the tower does not terminate,
confirming class~$M$.

The scalar gravitational partition function on the
scalar lane is
\begin{equation}\label{eq:vir-Z-grav}
  Z_{\mathrm{grav}}^{\mathrm{scal}}(\mathrm{Vir}_c;\,\hbar)
  = \frac{c}{2}\cdot
  \frac{\sqrt{\hbar}/2}{\sin(\sqrt{\hbar}/2)}
  \qquad\text{(UNIFORM-WEIGHT)},
\end{equation}
converging for $|\hbar| < 4\pi^2$.
This is $\kappa \cdot (\sqrt{\hbar}/2)/\sin(\sqrt{\hbar}/2)$
evaluated at $\kappa = c/2$.
At $c = 26$ (the bosonic string),
$\kappa_{\mathrm{eff}} = \kappa(\text{matter})
+ \kappa(\text{ghost}) = 13 + (-13) = 0$:
the scalar gravitational partition function vanishes, and
the non-scalar corrections from the shadow tower
become the leading contribution.

The full gravitational partition function receives
corrections from all arities of the shadow tower.
For class~$M$, infinitely many arities contribute at
each genus $g \geq 2$. The leading non-scalar correction
involves the quartic contact invariant
$S_4 = 10/(c(5c+22))$, which encodes the quartic
gravitational counterterm in the $3$d theory.

\medskip
\noindent\textbf{Koszul duality and complementarity.}
The Koszul dual is $\mathrm{Vir}_{26 - c}$, with
$\kappa + \kappa' = c/2 + (26 - c)/2 = 13$
(self-dual at $c = 13$; Virasoro complementarity $K = 13$,
distinct from the vanishing $K = 0$ of KM/Heis/lattice).
Under $c \mapsto 26 - c$:
\begin{equation}\label{eq:vir-compl-delta}
  \Delta(c) + \Delta(26 - c)
  = \frac{40}{5c + 22} + \frac{40}{152 - 5c}
  = \frac{6960}{(5c+22)(152 - 5c)}.
\end{equation}
The shadow growth rates are complementary:
$\rho(c) \cdot \rho(26 - c)$ is \emph{not} constant, but
both vanish in the limit $c \to \infty$.

\medskip
\noindent\textbf{The $3$d gravitational theory.}
The $3$d theory dual to $\mathrm{Vir}_c$ is
holomorphic-topological gravity in the Beltrami
parametrisation. The irregular singularity of the
factorisation $\cD$-module is the algebraic counterpart of
the unbounded spin tower in the gravitational theory:
the $T(z)$ OPE generates fields of arbitrarily high spin
through repeated collisions, producing the infinite shadow
tower $S_2, S_3, S_4, \ldots$. In the $3$d language,
each $S_r$ is a perturbative gravitational counterterm
at $r$-point contact. The essential singularity
$\exp(-c/(4z^2))$ in the flat sections of the KZ connection
is the chiral-algebraic incarnation of the non-perturbative
gravitational instanton corrections.

\medskip
\noindent\textbf{Summary.}
\begin{center}
\begin{tabular}{@{} c l l @{}}
\toprule
\textbf{Feature}
  & \textbf{Class $L$ (KM)}
  & \textbf{Class $M$ (Vir)} \\
\midrule
$r$-matrix pole & $1$ (simple) & $3$ (cubic) \\
Singularity type & regular (Fuchsian) & irregular (Poincar\'e rk.\ $2$) \\
Flat sections & $z^{\Omega/(k+h^\vee)}$ &
  $z^{2h}\exp(-c/(4z^2))$ \\
Monodromy data & representation of $\pi_1$
  & Stokes-filtered local system \\
RH correspondence & Deligne (classical)
  & D'Agnolo--Kashiwara \\
Stokes rays & $0$ & $4$ \\
Shadow tower & terminates ($r = 3$)
  & infinite ($r = \infty$) \\
$3$d theory & CS gauge theory
  & HT gravity \\
\bottomrule
\end{tabular}
\end{center}
\end{example}


% ================================================================
\section{Comparison with related work}
\label{sec:comparison}

Two existing programmes address the integration of braided
algebraic structures over surfaces:
Ben-Zvi--Brochier--Jordan integrate braided monoidal
\emph{categories}, while Latyntsev constructs chiral
homology for factorisation \emph{coalgebras}. Ordered chiral
homology occupies a different position: it integrates chiral
\emph{algebras} at the chain level, applies uniformly to all
four shadow classes, and does not require a finite ribbon
category.

\begin{remark}[Ben-Zvi--Brochier--Jordan]
\label{rem:bbj}
Ben-Zvi, Brochier, and Jordan~\cite{BBJ15,BBJ16}
construct a theory of factorisation homology for braided
monoidal categories: for a ribbon category $\cC$ and a
surface $\Sigma$, they define
$\int_\Sigma \cC$ as a categorical invariant.
Their construction and the present one operate at different
levels:
\begin{itemize}
\item BBJ integrate braided monoidal \emph{categories}
  over surfaces (categorical level).
\item Ordered chiral homology integrates chiral
  \emph{algebras} over curves (chain level).
\end{itemize}
For affine Kac--Moody: BBJ's $\int_\Sigma
\mathrm{Rep}_q(\fg)$ is a category, while
$\int_X^{\mathrm{ord}} V_k(\fg)$ is a chain complex.
The expected comparison is:
\begin{equation}\label{eq:bbj-comparison}
  \HH_*\Bigl(\int_\Sigma \mathrm{Rep}_q(\fg)\Bigr)
  \;\simeq\;
  \int_X^{\mathrm{ord}} V_k(\fg),
\end{equation}
where the left side is the Hochschild homology of the
categorical invariant and the right side is the ordered
chiral homology of the algebraic one.

For class $M$ (Virasoro at generic $c$): the BBJ
framework requires a finite ribbon category, which is
not available for generic central charge. Ordered chiral
homology applies without restriction to class $M$,
providing an integration theory where the categorical
approach does not (currently) reach.
\end{remark}

\begin{remark}[Latyntsev]
\label{rem:latyntsev}
Latyntsev~\cite{Latyntsev23} constructs chiral homology
for factorisation coalgebras on curves, providing an
algebraic counterpart to the analytic construction of
this paper. The relationship between the two approaches
at the chain level, particularly for irregular
singularities (class $M$), deserves further investigation.
\end{remark}

\begin{remark}[Wilson lines by ambient dimension]
\label{rem:wilson-line-classification}
The type of Wilson line, the form of the $R$-matrix, and
the resulting Yang--Baxter equation are controlled by the
ambient space dimension. The following table summarizes the
principal cases.

\begin{center}
\renewcommand{\arraystretch}{1.25}
\small
\begin{tabular}{@{} l l l l l @{}}
\toprule
Ambient & $R$-matrix & YBE & Algebra & Integrable system \\
\midrule
$\RR^3$ & constant $R$ & quantum $R_{12}R_{13}R_{23} = R_{23}R_{13}R_{12}$
  & $U_q(\fg)$ & Jones polynomial \\
$\RR \times \CC$ & $r(z)$, $1$ spectral & CYBE $[r_{12},r_{13}] + \mathrm{cyc} = 0$
  & $Y_\hbar(\fg)$ & Gaudin \\
$\RR^2 \times \CC$ & $R(u)$, spectral & quantum spectral YBE
  & $U_q(L\fg)$ & XXZ \\
$\CC^2$ & $r(z,w)$, $2$ params & $2$d CYBE
  & $\fg[u,v]$ & Hitchin on surface \\
$\RR \times \CC^2$ & $R(u,v)$, $2$ spectral & quantum + dynamical
  & $U_{q,t}(\widehat{\widehat{\fg}})$ & eCM + Hitchin \\
$\CC^3$ (CY) & $r + \ZZ \times \ZZ$ grading & $6$d YBE
  & $\Eone$-chiral QG & DT/BPS \\
\bottomrule
\end{tabular}
\end{center}

\noindent
This paper works in the $\RR \times \CC$ setting
(Costello's $4$d Chern--Simons~\cite{CFG26,CG17}),
producing the Yangian $Y_\hbar(\fg)$ as the
line-operator algebra.
The extensions to $\RR^2 \times \CC$ (quantum loop
algebras), $\CC^2$ (double current algebras),
$\RR \times \CC^2$ (quantum toroidal algebras), and
$\CC^3$ ($6$d holomorphic Chern--Simons defects) are
frontier directions; ordered chiral homology provides
the natural integration theory for the
$\RR \times \CC$ column.
\end{remark}


% ================================================================
\section{Chiral quantum groups: $R$-matrices, $\Ainf$-structures,
and chiral coproducts}
\label{sec:chiral-quantum-groups}

The classical theory of quantum groups rests on an equivalence
between three descriptions of the same structure: a
quasi-triangular Hopf algebra $(H, \Delta, R)$ can be
specified by its $R$-matrix, by its coproduct, or by an
operadic $\Ainf$-algebra map. The goal of this section is to
establish the chiral analogue: an equivalence between three
descriptions of an $\Eone$-chiral algebra equipped with
quantum-group-type structure. The primitive datum is the
$\Eone$ ordered bar complex
$\Barord(\cA) = T^c(s^{-1}\bar{\cA})$ with its
deconcatenation coproduct, coassociative over
$(\mathrm{ChirAss})^!$; the $R$-matrix and the chiral
coproduct are both recoverable from this datum.

% ----------------------------------------------------------------
\subsection{Three structures on an $\Eone$-chiral algebra}
\label{subsec:three-structures}

An $\Eone$-chiral algebra admits three equivalent
descriptions: by a vertex $R$-matrix $S(z)$ (braiding data),
by chiral $\Ainf$ structure maps $m_k^{\mathrm{ch}}$
(higher associativity data), and by a chiral coproduct
$\Delta^{\mathrm{ch}}$ (coalgebra data). Each description
captures a different projection of the ordered bar complex
$\Barord(\cA) = T^c(s^{-1}\bar{\cA})$.

\begin{definition}[$\Eone$-chiral algebra with vertex
$R$-matrix]
\label{def:e1-chiral-rmatrix}
An $\Eone$-chiral algebra with vertex $R$-matrix is a pair
$(\cA, S(z))$ where $\cA$ is a graded vertex algebra and
$S(z) \in \End(\cA \otimes \cA)[[z, z^{-1}]]$ is a vertex
$R$-matrix satisfying:
\begin{enumerate}[label=\textup{(\roman*)}]
\item \textup{($S$-locality.)} For all $a, b \in \cA$
  and sufficiently large $N$:
  \begin{equation}\label{eq:s-locality-sec8}
    (z - w)^N\, Y(a,z)\, Y(b,w)
    \;=\;
    (z - w)^N\, Y(b,w)\, Y(a,z) \cdot S(z - w).
  \end{equation}
\item \textup{(Quantum Yang--Baxter equation.)} In
  $\End(\cA^{\otimes 3})$:
  \begin{equation}\label{eq:qybe}
    S_{12}(z_1 - z_2)\, S_{13}(z_1 - z_3)\, S_{23}(z_2 - z_3)
    \;=\;
    S_{23}(z_2 - z_3)\, S_{13}(z_1 - z_3)\, S_{12}(z_1 - z_2).
  \end{equation}
\item \textup{(Unitarity.)}
  $S_{12}(z)\, S_{21}(-z) = \id$.
\item \textup{(Shift condition.)} The $R$-matrix
  intertwines the translation operator $T$:
  $[T \otimes \id + \id \otimes T,\, S(z)]
  = \partial_z S(z)$.
\item \textup{(Hexagon axiom.)} On
  $\overline{\FM}_3^{\mathrm{ord}}(\CC)$, the
  braiding and associativity satisfy the coherence:
  for the two standard paths in
  $\overline{\FM}_3^{\mathrm{ord}}(\CC)$ connecting the
  same boundary strata,
  \begin{equation}\label{eq:hexagon}
    S_{12}(z_{12})\, \Phi_{123}^{-1}\, S_{13}(z_{13})\, \Phi_{132}
    \;=\; \Phi_{213}^{-1}\, S_{23}(z_{23})\, \Phi_{123}
  \end{equation}
  where $\Phi_{ijk}$ is the Drinfeld associator
  evaluated on the three-point KZ connection.
\end{enumerate}
These are the axioms of Etingof--Kazhdan~\cite{EK00} for a
quantum vertex operator algebra.
\end{definition}

\begin{definition}[$\Ainf$-algebra in the chiral
endomorphism operad]
\label{def:ainf-chiral-endo}
An $\Ainf$-algebra structure on $\cA$ in the chiral
endomorphism operad $\End^{\mathrm{ch}}_\cA$ is an operadic
map
\begin{equation}\label{eq:ainf-map}
  \Ainf
  \;\longrightarrow\; \End^{\mathrm{ch}}_\cA
\end{equation}
where $\End^{\mathrm{ch}}_\cA(n)
= \Hom(\cA^{\otimes n},\;
\cA((\lambda_1))\cdots((\lambda_{n-1})))$ is the spectral
substitution operad
\textup{(}\cite[\S3]{Lorgat26I}\textup{)}. Concretely, the
data consist of chiral structure maps
\begin{equation}\label{eq:ch-mk}
  m_k^{\mathrm{ch}} \colon
  \cA^{\otimes k}
  \longrightarrow
  \cA((\lambda_1))\cdots((\lambda_{k-1})),
  \qquad k \geq 1,
\end{equation}
satisfying the Stasheff $\Ainf$-relations in the chiral
endomorphism operad: for each $n \geq 1$,
\begin{equation}\label{eq:stasheff-chiral}
  \sum_{\substack{i + j = n + 1 \\ 1 \leq p \leq i}}
  (-1)^{\epsilon}\,
  m_i^{\mathrm{ch}} \circ_p m_j^{\mathrm{ch}}
  \;=\; 0,
\end{equation}
where $\circ_p$ is the partial composition via spectral
substitution \textup{(}
$\lambda$-variable nesting\textup{)}, and $\epsilon$
carries the standard Koszul sign from the suspended degree
$\|m_k\| = k - 2$.
The structure maps $m_k^{\mathrm{ch}}$ are meromorphic in
the spectral parameters $\lambda_1, \ldots, \lambda_{k-1}$.
At $k = 2$, the map $m_2^{\mathrm{ch}}$ is the OPE vertex
operation $Y(-,z)$. At $k = 1$, $m_1^{\mathrm{ch}} = d$
is the differential. The operations $m_k^{\mathrm{ch}}$ for
$k \geq 3$ record the obstruction to strict associativity
of the chiral product.
\end{definition}

\begin{definition}[$\Eone$-chiral algebra with compatible
chiral coproduct]
\label{def:chiral-coproduct}
An $\Eone$-chiral algebra with compatible chiral coproduct
is a triple $(\cA, Y, \Delta^{\mathrm{ch}})$ where
$(\cA, Y)$ is a chiral algebra and
\begin{equation}\label{eq:chiral-coprod}
  \Delta^{\mathrm{ch}} \colon
  \cA
  \longrightarrow
  (\cA \mathbin{\hat{\otimes}} \cA)((z))
\end{equation}
is a chiral coproduct \textup{(}a map to the completed
tensor product with spectral parameter\textup{)} satisfying:
\begin{enumerate}[label=\textup{(\roman*)}]
\item \textup{(Coassociativity up to associator.)}
  The two iterated coproducts agree up to conjugation by
  the Drinfeld associator $\Phi$:
  \begin{equation}\label{eq:coassoc}
    (\Delta^{\mathrm{ch}} \otimes \id)
    \circ \Delta^{\mathrm{ch}}
    \;=\; \Phi \cdot
    (\id \otimes \Delta^{\mathrm{ch}})
    \circ \Delta^{\mathrm{ch}}
    \cdot \Phi^{-1}.
  \end{equation}
\item \textup{(OPE compatibility.)} The coproduct
  intertwines the vertex operation:
  for $a, b \in \cA$,
  \begin{equation}\label{eq:ope-compat}
    \Delta^{\mathrm{ch}}(Y(a,z)\, b)
    \;=\;
    Y^{(2)}(\Delta^{\mathrm{ch}}(a), z)
    \cdot \Delta^{\mathrm{ch}}(b)
  \end{equation}
  where $Y^{(2)}$ denotes the tensorwise vertex
  operation on $\cA \mathbin{\hat{\otimes}} \cA$.
\item \textup{(Counit.)} There exists a counit
  $\varepsilon \colon \cA \to \CC$ with
  $(\varepsilon \otimes \id) \circ \Delta^{\mathrm{ch}}
  = \id = (\id \otimes \varepsilon) \circ
  \Delta^{\mathrm{ch}}$, recovering the augmentation of
  $\cA$.
\end{enumerate}
The chiral coproduct determines a vertex $R$-matrix by
\begin{equation}\label{eq:r-from-coprod}
  S(z)
  \;:=\; \sigma \circ \Delta^{\mathrm{ch}}
  \circ (\Delta^{\mathrm{ch},\op})^{-1}
\end{equation}
where $\sigma$ is the graded flip and
$\Delta^{\mathrm{ch},\op}
= \sigma \circ \Delta^{\mathrm{ch}}$ is the opposite
coproduct. This is the chiral analogue of the classical
formula $R = \Delta^{\op}(\Delta^{-1}(-))$ for a
quasi-triangular Hopf algebra.
\end{definition}

\begin{remark}[The role of the Drinfeld associator]
\label{rem:drinfeld-assoc-role}
The Drinfeld associator $\Phi \in \exp(\widehat{\Lie}(t_{12},
t_{23}))$ controls the failure of strict coassociativity:
the two parenthesizations traverse different paths in
$\overline{\FM}_3^{\mathrm{ord}}(\CC)$, and $\Phi$ is the
KZ parallel transport between them. It also appears in
$\ker(\av_3)$ of
Proposition~\ref{prop:lossy-descent},
invisible to the symmetric shadow. For $\Einf$-chiral
algebras, $\Phi$ acts trivially on coinvariants; for genuinely
$\Eone$-chiral algebras, it is essential structural data.
\end{remark}


% ----------------------------------------------------------------
\subsection{The equivalence theorem}
\label{subsec:equiv-thm}

The three descriptions of the preceding subsection are not
independent: a vertex $R$-matrix determines the chiral
$\Ainf$ operations via the bar complex, the $\Ainf$
operations determine the chiral coproduct via cobar
dualization, and the coproduct recovers the $R$-matrix.
The equivalences close into a triangle, each arrow depending
on the choice of a Drinfeld associator $\Phi$.

\begin{theorem}[Chiral quantum group equivalence]
\label{thm:chiral-qg-equiv}
The following three structures on an $\Eone$-chiral
algebra $\cA$ on the Koszul locus determine each other,
up to the choice of a Drinfeld associator $\Phi$:
\begin{center}
\renewcommand{\arraystretch}{1.3}
\begin{tabular}{cl}
\textup{(I)} & $\Eone$-chiral algebra with vertex
  $R$-matrix $S(z)$ \\
  & \textup{(}Definition~\textup{\ref{def:e1-chiral-rmatrix}}\textup{)} \\
\textup{(II)} & $\Ainf$-algebra in the chiral endomorphism
  operad $\End^{\mathrm{ch}}_\cA$ \\
  & \textup{(}Definition~\textup{\ref{def:ainf-chiral-endo}}\textup{)} \\
\textup{(III)} & $\Eone$-chiral algebra with compatible
  chiral coproduct $\Delta^{\mathrm{ch}}$ \\
  & \textup{(}Definition~\textup{\ref{def:chiral-coproduct}}\textup{)}
\end{tabular}
\end{center}
The equivalences form a triangle:
\begin{equation}\label{eq:equiv-triangle}
\begin{tikzcd}[column sep=3em, row sep=3em]
  \textup{(I)} \ar[rr, shift left=1, "\alpha"]
  \ar[dr, shift left=1, "\gamma^{-1}"']
  && \textup{(II)} \ar[ll, shift left=1, "\alpha^{-1}"]
  \ar[dl, shift left=1, "\beta"] \\
  & \textup{(III)}
  \ar[ul, shift left=1, "\gamma"]
  \ar[ur, shift left=1, "\beta^{-1}"']
\end{tikzcd}
\end{equation}
where each arrow is an equivalence of categories.
\end{theorem}

\begin{proof}
We prove each direction of the
triangle~\eqref{eq:equiv-triangle}.

\medskip
\noindent
\textbf{(I) $\boldsymbol{\to}$ (II): $R$-matrix gives
$\Ainf$ structure.}
The Stasheff associahedron $K_n$ is isomorphic to the
real Fulton--MacPherson compactification
$\overline{\FM}_n^{\mathrm{ord}}(\RR)$
(Sinha~\cite{Sinha04}); the inclusion
$\RR \hookrightarrow \CC$ gives the embedding
$K_n \hookrightarrow \overline{\FM}_n^{\mathrm{ord}}(\CC)$
of~\eqref{eq:assoc-embed}.
The chiral $\Ainf$-maps $m_k^{\mathrm{ch}}$ are defined by
restricting the OPE to this real locus:
\begin{equation}\label{eq:mk-from-ope}
  m_k^{\mathrm{ch}}(a_1, \ldots, a_k)
  \;=\;
  \int_{K_k}
  Y(a_1, z_1) \cdots Y(a_k, z_k)\, \omega_k,
\end{equation}
where $\omega_k$ is the canonical volume form on $K_k$
induced by the bar propagator.

The Stasheff identity at arity $n$
(equation~\eqref{eq:stasheff-chiral}) follows from Stokes'
theorem on $K_{n+1}
= \overline{\FM}_{n+1}^{\mathrm{ord}}(\RR)$.
The codimension-$1$ boundary $\partial K_{n+1}$ has
$\binom{n+1}{2} - 1$ facets, one for each consecutive
cluster $(z_i, z_{i+1}, \ldots, z_{i+j})$ of length
$j + 1 \geq 2$ (with $i + j \leq n + 1$ and
$j + 1 \leq n$). Each facet is
isomorphic to $K_{r+1} \times K_{s+1}$ with $r + s = n$,
recording the cluster and the residual configuration.
The integral of $d\omega_{n+1}$ over $K_{n+1}$ vanishes,
and the boundary contributions give
$\sum (-1)^{\epsilon}\, m_i^{\mathrm{ch}} \circ_p
m_j^{\mathrm{ch}} = 0$: the ordered AOS reduction
(\cite[\S5]{Lorgat26I}).

At codimension $2$, where two simultaneous consecutive
collisions $(z_i \to z_{i+1}, z_j \to z_{j+1})$ with
$|i - j| \geq 2$ occur, the compatibility condition
is the hexagon axiom~\eqref{eq:hexagon}: the two
codimension-$2$ strata correspond to the two standard
paths in $\overline{\FM}_3^{\mathrm{ord}}(\CC)$, and
the hexagon ensures that the iterated boundary
contributions match. At codimension $3$ (triple
simultaneous collision in
$\overline{\FM}_4^{\mathrm{ord}}(\CC)$), the
compatibility is the pentagon relation for $\Phi$, which
is equivalent to the QYBE~\eqref{eq:qybe} for $S(z)$.

$S$-locality~\eqref{eq:s-locality-sec8} ensures that the
integrals~\eqref{eq:mk-from-ope} converge: the factor
$(z - w)^N$ in $S$-locality bounds the pole orders
of the OPE along each boundary divisor.

This construction defines the functor $\alpha$. Its
inverse $\alpha^{-1}$ recovers the $R$-matrix from the
$\Ainf$-structure as the arity-$2$ holonomy:
$S(z) = \cP \exp\bigl(\int_{\gamma(z)}
m_2^{\mathrm{ch}}\bigr)$,
where $\gamma(z)$ is the KZ parallel-transport path in
$\Conf_2^{\mathrm{ord}}(\CC)$.

\medskip
\noindent
\textbf{(II) $\boldsymbol{\to}$ (III): $\Ainf$ structure
gives chiral coproduct.}
The ordered bar complex
$\Barord(\cA) = T^c(s^{-1}\bar{\cA})$
(where $\bar{\cA} = \ker(\varepsilon)$ is the augmentation
ideal) carries the deconcatenation coproduct
\begin{equation}\label{eq:deconc-coprod}
  \Delta[a_1 | \cdots | a_n]
  = \sum_{i=0}^{n}
  [a_1 | \cdots | a_i]
  \otimes [a_{i+1} | \cdots | a_n].
\end{equation}
The $\Ainf$-maps $m_k^{\mathrm{ch}}$ determine the bar
differential $d_{\Barord}$ by the standard formula
$d_{\Barord} = \sum_{k \geq 1} \hat{m}_k^{\mathrm{ch}}$,
where $\hat{m}_k^{\mathrm{ch}}$ is the coderivation
extension of $m_k^{\mathrm{ch}}$ to
$T^c(s^{-1}\bar{\cA})$. The relation
$d_{\Barord}^2 = 0$ is equivalent to the Stasheff
identities~\eqref{eq:stasheff-chiral}, so the $\Ainf$
data make $(\Barord(\cA), d_{\Barord}, \Delta)$ into a
$E_1$-chiral coassociative dg coalgebra
over $(\mathsf{ChirAss})^!$
(\cite[Theorem~1.2]{Lorgat26I}): the differential
$d_{\Barord}$ comes from the chiral product (OPE collision residues) and the
deconcatenation $\Delta$ comes from the cofree tensor coalgebra structure.

The chiral coproduct on $\cA$ is constructed as follows.
On the Koszul locus (where the bar-cobar adjunction is
an equivalence), the cobar functor $\Omega$ applied to
$\Barord(\cA)$ recovers $\cA$:
$\cA \simeq \Omega(\Barord(\cA))$
(\cite[Theorem~1.1(i)]{Lorgat26I}).
The deconcatenation $\Delta$ on $\Barord(\cA)$ is a
morphism of dg coalgebras
$\Delta \colon \Barord(\cA) \to
\Barord(\cA) \otimes \Barord(\cA)$; applying $\Omega$
produces
\begin{equation}\label{eq:coprod-from-bar}
  \Delta^{\mathrm{ch}} \colon
  \cA \simeq \Omega(\Barord(\cA))
  \xrightarrow{\Omega(\Delta)}
  \Omega(\Barord(\cA) \otimes \Barord(\cA))
  \to (\cA \mathbin{\hat{\otimes}} \cA)((z)),
\end{equation}
where the last map uses the K\"unneth quasi-isomorphism
$\Omega(C \otimes D)
\xrightarrow{\sim}
\Omega(C) \mathbin{\hat{\otimes}} \Omega(D)$ for
conilpotent coalgebras.

The deconcatenation coproduct is strictly coassociative
on $T^c(s^{-1}\bar{\cA})$: the two iterated coproducts
$(\Delta \otimes \id) \circ \Delta
= (\id \otimes \Delta) \circ \Delta$ by direct
computation (the $(n+1)(n+2)/2$ terms on each side match
by the unique decomposition of a word into three
consecutive subwords). The associator $\Phi$ enters when
passing through the cobar functor: the formality
quasi-isomorphism underlying $\Omega$ depends on $\Phi$,
so $\Delta^{\mathrm{ch}}$ is coassociative only up to
conjugation by $\Phi$, as required by
Definition~\ref{def:chiral-coproduct}(i).

The OPE compatibility~\eqref{eq:ope-compat} follows from
the coderivation property of $d_{\Barord}$ with respect
to $\Delta$: the bar differential commutes with the
coproduct, which dualizes under $\Omega$ to the
intertwining of $\Delta^{\mathrm{ch}}$ with the vertex
operation $Y$. The counit
$\varepsilon \colon \cA \to \CC$ is the augmentation, and
the counit axiom
(Definition~\ref{def:chiral-coproduct}(iii)) holds because
the bar complex is built on
$\bar{\cA} = \ker(\varepsilon)$.

This construction defines the functor $\beta$. Its
inverse $\beta^{-1}$ recovers the $\Ainf$-maps from a
chiral coproduct by reading off the bar differential: the
$k$-th structure map $m_k^{\mathrm{ch}}$ is the component
of $d_{\Barord}$ at tensor degree $k$.

\medskip
\noindent
\textbf{(III) $\boldsymbol{\to}$ (I): chiral coproduct
gives $R$-matrix.}
Given a chiral coproduct $\Delta^{\mathrm{ch}}$ satisfying
Definition~\ref{def:chiral-coproduct}, the vertex
$R$-matrix is defined by~\eqref{eq:r-from-coprod}:
\begin{equation}\label{eq:r-from-coprod-proof}
  S(z)
  \;:=\; \sigma \circ \Delta^{\mathrm{ch}}
  \circ (\Delta^{\mathrm{ch},\op})^{-1},
\end{equation}
the chiral analogue of the classical Drinfeld formula for
quasi-triangular quasi-Hopf
algebras~\cite{Drinfeld90}.
We verify the five axioms of
Definition~\ref{def:e1-chiral-rmatrix}.

\emph{QYBE.}
Expand the left side of~\eqref{eq:qybe} using the
definition~\eqref{eq:r-from-coprod-proof}. Each factor
$S_{ij}(z_{ij})$ introduces one coproduct and one
inverse opposite coproduct; composing the three factors
yields an expression involving three applications of the
coassociativity relation~\eqref{eq:coassoc}. In each
application, the two iterated coproducts differ by
conjugation by $\Phi$. The pentagon identity for $\Phi$
(Drinfeld~\cite{Drinfeld90}; see also
Etingof--Kazhdan~\cite{EK96}, \S3) ensures that the
three associator insertions cancel pairwise, and the
remaining coproduct factors rearrange to give the right
side of~\eqref{eq:qybe}.

\emph{Unitarity.}
By definition, $S_{21}(-z) = \sigma \circ S_{12}(-z)
\circ \sigma$, and
$\Delta^{\mathrm{ch},\op,\op} = \Delta^{\mathrm{ch}}$
(applying the flip twice is the identity). Composing:
\[
  S_{12}(z)\, S_{21}(-z)
  = \sigma \circ \Delta^{\mathrm{ch}} \circ
  (\Delta^{\mathrm{ch},\op})^{-1}
  \cdot
  \sigma \circ \Delta^{\mathrm{ch},\op} \circ
  (\Delta^{\mathrm{ch}})^{-1}
  = \id.
\]

\emph{Hexagon.}
The hexagon~\eqref{eq:hexagon} encodes the compatibility
of $S(z)$ with $\Phi$ on three-point configurations.
Substituting~\eqref{eq:r-from-coprod-proof} into the left
side, each $S_{ij}$ contributes a coproduct and an inverse
opposite coproduct, and each $\Phi_{ijk}^{\pm 1}$
reparenthesizes. The OPE
compatibility~\eqref{eq:ope-compat} (the chiral Leibniz
rule: $\Delta^{\mathrm{ch}}$ intertwines $Y$) ensures
that the product-coproduct rearrangement through $\Phi$
matches the right side. This computation is the
vertex-algebraic analogue of Drinfeld's proof that a
quasi-Hopf coproduct produces a quasi-triangular
structure~\cite{Drinfeld90}; the adaptation to the chiral
setting (spectral-parameter dependence, meromorphic
completions) follows
Etingof--Kazhdan~\cite{EK00}, \S4.

\emph{Shift condition.}
The translation operator $T = L_{-1}$ satisfies
$[T, Y(a,z)] = \partial_z Y(a,z)$.
The OPE compatibility~\eqref{eq:ope-compat} applied to
the vacuum gives
$\Delta^{\mathrm{ch}}(T \cdot a)
= (T \otimes \id + \id \otimes T) \cdot
\Delta^{\mathrm{ch}}(a)$ (the counit axiom ensures the
two tensor factors share $T$ additively).
Differentiating~\eqref{eq:r-from-coprod-proof} in $z$
and using this intertwining yields
$[T \otimes \id + \id \otimes T,\, S(z)]
= \partial_z S(z)$.

\emph{$S$-locality.}
The OPE compatibility~\eqref{eq:ope-compat} maps
singular OPE terms to tensor products of singular terms.
For $a, b \in \cA$ and $N$ sufficiently large,
$(z - w)^N Y(a,z) Y(b,w)$ is regular. Applying
$\Delta^{\mathrm{ch}}$ to both sides of the OPE and using
the definition of $S(z-w)$
via~\eqref{eq:r-from-coprod-proof} gives
$(z - w)^N Y(a,z) Y(b,w)
= (z - w)^N Y(b,w) Y(a,z) \cdot S(z - w)$,
which is~\eqref{eq:s-locality-sec8}.

This construction defines $\gamma$. Its inverse
$\gamma^{-1}$ recovers the coproduct from
$(\cA, S(z))$ by composing $\alpha$ and $\beta$:
$\gamma^{-1} = \beta \circ \alpha$.

\medskip
\noindent
\textbf{Triangle coherence.}
The three functors satisfy
$\gamma = \alpha^{-1} \circ \beta^{-1}$. Each functor
is an equivalence with explicit inverse as constructed
above, and the triangle~\eqref{eq:equiv-triangle}
commutes: $\alpha \circ \gamma \circ \beta = \id$.
\end{proof}

The equivalence theorem packages three equivalent descriptions
of quantum-group-type structure on an $\Eone$-chiral algebra.
The following definition records the datum that the equivalence
classifies: the chiral Yangian, assembling all four components
into a single object.

\begin{definition}[Chiral Yangian]
\label{def:chiral-yangian}
For simple $\fg$, smooth curve $X$, and formal parameter
$\hbar$, the \emph{chiral Yangian}
$Y^{\mathrm{ch}}_\hbar(\fg)$ is the datum
$(\cA,\, S(z),\, \Delta^{\mathrm{ch}},\,
\{m_k^{\mathrm{ch}}\})$ where:
\begin{enumerate}[label=\textup{(\roman*)}]
\item $\cA$ is an $\Eone$-chiral algebra on
  $\Ran^{\mathrm{ord}}(X)$ with $R$-twisted $S$-locality
  \textup{(}Definition~\textup{\ref{def:e1-chiral-rmatrix}}\textup{)};
\item $S(z) = R_{21}(z)\, R_{12}(-z)$ is the spectral
  $S$-matrix \textup{(}monodromy on
  $\Conf_2^{\mathrm{ord}}(X)$\textup{)};
\item $\Delta^{\mathrm{ch}} \colon
  \cA \to \cA \mathbin{\hat{\otimes}} \cA$ is a
  quasi-coassociative chiral coproduct satisfying
  quasi-triangularity:
  $R(z)\, \Delta^{\mathrm{ch}}
  = \Delta^{\mathrm{ch},\op}\, R(z)$
  \textup{(}Definition~\textup{\ref{def:chiral-coproduct}}\textup{)};
\item $\{m_k^{\mathrm{ch}}\}_{k \geq 1}$ is the curved
  $\Ainf$-structure from the bar Maurer--Cartan element
  $\Theta_\cA \in \MC(\gEone)$
  \textup{(}Definition~\textup{\ref{def:ainf-chiral-endo}}\textup{)}.
\end{enumerate}
These data satisfy four axioms:
\begin{enumerate}[label=\textup{(A\arabic*)}]
\item \textup{(Koszul self-recovery.)}
  $\Omega^{\mathrm{ch}}(\Barord(\cA))
  \xrightarrow{\;\sim\;} \cA$;
\item \textup{(Verdier by $R$-inversion.)}
  $\cA^{!,\mathrm{ch}}
  \cong Y^{\mathrm{ch}}_{-\hbar}(\fg)$;
\item \textup{(Braided module category.)}
  $\Mod_\cA$ is $\Etwo$ with braiding
  $\sigma = P \circ R(u)$;
\item \textup{(Modular MC tower.)}
  $R^{\mathrm{mod}}(z;\hbar)
  = \sum_{g \geq 0} \hbar^{2g}\, r_g(z)
  \in \MC(\widehat{L}^{\mathrm{mod}}_\cA)$.
\end{enumerate}
\end{definition}

\begin{remark}[Constructive content of the bar complex]
\label{rem:chiral-yangian-bar}
The ordered bar complex
$\Barord(\cA) = T^c(s^{-1}\bar{\cA})$ produces three of the
four components constructively: the $\Eone$-chiral algebra
structure~(i) is the bar-cobar adjunction
\textup{(}Theorem~A of \cite{Lorgat26I}\textup{)}; the
spectral $S$-matrix~(ii) is the arity-$2$ holonomy of the
KZ connection on $\Conf_2^{\mathrm{ord}}(X)$; and the
$\Ainf$-structure~(iv) is the coderivation extension of
$\Theta_\cA$ to $T^c(s^{-1}\bar{\cA})$. The chiral
coproduct~(iii) requires the additional datum of the
$E_1$-chiral coassociative structure: the deconcatenation coproduct on the
bar complex is the cofree tensor coalgebra structure, and its passage
to $\cA$ via the cobar functor $\Omega$ uses the
K\"unneth quasi-isomorphism for conilpotent coalgebras
\textup{(}equation~\textup{\eqref{eq:coprod-from-bar}}\textup{)}.
\end{remark}


% ----------------------------------------------------------------
\subsection{The classical limit}
\label{subsec:classical-limit}

The chiral quantum group equivalence
(Theorem~\ref{thm:chiral-qg-equiv}) is a vertex-algebraic
lift of Drinfeld's classical equivalence between
quasi-triangular Hopf algebras, $R$-matrices, and
coproducts. The lift replaces finite-dimensional vector
spaces with modules over $\CC((z))$, the symmetric group
with the braid group, and strict coassociativity with
coassociativity up to the Drinfeld associator $\Phi$.

\begin{remark}[Comparison with the classical Hopf algebra
story]
\label{rem:classical-hopf}
The chiral quantum group equivalence
(Theorem~\ref{thm:chiral-qg-equiv}) is the
vertex-algebraic analogue of the following classical
equivalence, due to Drinfeld:
\begin{center}
\renewcommand{\arraystretch}{1.3}
\begin{tabular}{cl}
\textup{(I)$_{\mathrm{cl}}$} &
  Quasi-triangular quasi-Hopf algebra $(H, \Delta, R, \Phi)$ \\
\textup{(II)$_{\mathrm{cl}}$} &
  $\Ainf$-algebra in the endomorphism operad $\End_H$ \\
\textup{(III)$_{\mathrm{cl}}$} &
  Bialgebra with $R$-matrix: $R \cdot \Delta(x)
  = \Delta^{\op}(x) \cdot R$ for all $x \in H$
\end{tabular}
\end{center}
The passage from classical to chiral introduces two features:
the coproduct acquires a spectral parameter
($\Delta^{\mathrm{ch}} \colon
\cA \to (\cA \mathbin{\hat{\otimes}} \cA)((z))$),
and strict coassociativity is replaced by coassociativity
up to the Drinfeld associator $\Phi$, which appears
universally in the chiral setting because the KZ connection
on $\overline{\FM}_3^{\mathrm{ord}}(\CC)$ has non-trivial
holonomy.
\end{remark}


% ----------------------------------------------------------------
\subsection{Examples}
\label{subsec:cqg-examples}

The three structures are most transparent for affine
Kac--Moody, where the $R$-matrix is the KZ monodromy, the
$\Ainf$ operations are the Sugawara OPE restricted to the
associahedron, and the chiral coproduct is the Drinfeld
coproduct. For the Yangian $Y_\hbar(\fg)$, the $R$-matrix
is the starting point rather than derived from an
$\Einf$-chiral OPE. For $\Einf$-chiral algebras, all three
structures are trivial.

\begin{example}[Affine Kac--Moody: Sugawara coproduct]
\label{ex:km-coproduct}
For $V_k(\fg)$ at generic level $k \neq -h^\vee$, the three
structures are:
\begin{enumerate}[label=\textup{(\roman*)}]
\item \textup{($R$-matrix.)} The vertex $R$-matrix is the
  KZ $R$-matrix: the monodromy of the KZ connection
  % AP126/AP148: trace-form convention; k=0 -> r=0. Verified.
  $\nabla = d + k\Omega/(z_1 - z_2) \cdot d(z_1 - z_2)$
  around $\Delta_{12}$. The classical $r$-matrix
  (trace-form) is
  $r(z) = k\Omega/z$; at $k = 0$, $r = 0$; at
  $k = -h^\vee$, the critical level, the Sugawara
  construction degenerates.
\item \textup{($\Ainf$ structure.)} The chiral $\Ainf$
  structure maps are the Sugawara OPE operations restricted
  to $K_n \hookrightarrow \overline{\FM}_n^{\mathrm{ord}}(\CC)$.
  At $k = 2$: $m_2^{\mathrm{ch}}(J^a, J^b; \lambda)
  = [J^a, J^b] + k\, \delta^{ab}/\lambda$ is the affine
  OPE. The transferred higher operations
  $m_k^{\mathrm{ch}}$ for $k \geq 3$ encode the
  associator corrections from the KZ holonomy.
\item \textup{(Chiral coproduct.)} The Drinfeld
  coproduct on $Y_\hbar(\fg)$ is the paradigmatic example.
  On generators: $\Delta^{\mathrm{ch}}(J^a) = J^a \otimes 1
  + 1 \otimes J^a + \hbar\, f^a_{bc}\, J^b \otimes J^c / z
  + \cdots$, where the spectral parameter $z$ comes from
  the difference coordinate $z_1 - z_2$ on
  $\Conf_2(\CC)$. The deconcatenation coproduct on the
  ordered bar $\Barord(V_k(\fg))$ induces, via the cobar
  functor, this spectral Drinfeld coproduct.
\end{enumerate}
The Drinfeld--Kohno theorem identifies the KZ monodromy
representation (the $R$-matrix of (i)) with the quantum
group $R$-matrix on $U_q(\fg)$-modules, where
$q = \exp(\pi i / (k + h^\vee))$. The chiral quantum group
equivalence (Theorem~\ref{thm:chiral-qg-equiv}) lifts this
identification from monodromy representations to the level
of the algebra itself.
\end{example}

\begin{example}[Yangian: the Drinfeld coproduct as chiral
primitive]
\label{ex:yangian-coproduct}
For the Yangian $Y_\hbar(\fg)$, structure~(I) is the
starting point: the $R$-matrix
$R(z) = 1 + \hbar\, r/z + O(\hbar^2)$ is independent
input (not derived from any $\Einf$-chiral OPE). The
Drinfeld coproduct
\begin{equation}\label{eq:yangian-coprod}
  \Delta_z(T^{(n)})
  \;=\; \sum_{k=0}^{n}
  T^{(k)} \otimes T^{(n-k)} \cdot z^{-k}
\end{equation}
(on the generating series
$T(u) = 1 + \sum_{n \geq 1} T^{(n)} u^{-n}$) provides
structure~(III). The ordered bar complex
$\Barord(Y_\hbar(\fg))$ is the bridge: its
deconcatenation coproduct is strictly coassociative and
induces, via cobar dualization, the Drinfeld
coproduct~\eqref{eq:yangian-coprod} on $Y_\hbar(\fg)$.
The associator $\Phi$ appears when passing from the strictly
coassociative bar level to the algebra level.

The $\Ainf$ structure maps $m_k^{\mathrm{ch}}$ for
$k \geq 3$ are the RTT relations of the Yangian: they
encode the obstruction to strict associativity of the
spectral OPE, governed by the fusion procedure for the
$R$-matrix. The bar-cobar inversion
$\Omega(\Barord(Y_\hbar(\fg))) \xrightarrow{\sim}
Y_\hbar(\fg)$ identifies these operations with the
standard Yangian relations.
\end{example}

\begin{example}[Etingof--Kazhdan quantum vertex algebra for
$\mathfrak{sl}_2$: full braided VOA structure]
\label{ex:ek-qvoa}
The Yangian $Y_\hbar(\mathfrak{sl}_2)$ of
Example~\ref{ex:yangian} is a quasi-triangular Hopf algebra
in the category of topological algebras with a spectral
parameter. The Etingof--Kazhdan construction~\cite{EK00}
lifts it to a \emph{quantum vertex operator algebra}:
an algebra satisfying the full suite of braided VOA axioms
($S$-locality, vertex $R$-matrix, hexagon) of
Definition~\ref{def:e1-chiral-rmatrix}.
The purpose of this example is to exhibit the construction
for $\fg = \mathfrak{sl}_2$ concretely, verify the axioms,
compute the ordered chiral center with the full quantum
$R$-matrix $S(z)$ (not merely its classical leading order),
and connect $S$-locality to the chiral $\Ainf$ structure.

\medskip
\noindent
\textbf{Part A: the EK construction.}
The input is the Yang $R$-matrix
$R(u) = u\,I + \hbar\,P \in \End(V \otimes V)[u]$
from~\eqref{eq:yangian-rmatrix}, where
$V = \CC^2$ is the fundamental representation,
$P = (\rho \otimes \rho)(\Omega) + \tfrac{1}{2}\,I$ is the
permutation operator, and $\hbar = 1/(k+2)$.
The EK quantum VOA is constructed in four steps:
quantum function algebra, dual QUE algebra (state space),
vertex operator (creation--annihilation decomposition), and
vertex $R$-matrix.

\smallskip
\noindent\emph{Step 1: the quantum function algebra $F_0(R)$.}
The algebra $F_0(R)$ is generated by the matrix entries
of a formal power series
\begin{equation}\label{eq:ek-transfer}
  T(u) = \sum_{n \geq 0} T^{(n)}\, u^n
  \;\in\; \End(V) \otimes F_0(R)[\![u]\!],
\end{equation}
subject to the RTT relation
\begin{equation}\label{eq:ek-rtt}
  R_{12}(u - v)\,T_{13}(u)\,T_{23}(v)
  = T_{23}(v)\,T_{13}(u)\,R_{12}(u - v)
\end{equation}
in $\End(V^{\otimes 2}) \otimes F_0(R)[\![u, v]\!]$.
For $R(u) = u\,I + \hbar\,P$, writing
$T^{(n)} = \bigl(\begin{smallmatrix}
  a_n & b_n \\ c_n & d_n
\end{smallmatrix}\bigr)$, equation~\eqref{eq:ek-rtt}
expands at order $u^m v^n$ to give:
\begin{align}
  [a_m, a_n] &= 0, \quad [d_m, d_n] = 0,
  \quad [b_m, b_n] = 0, \quad [c_m, c_n] = 0,
  \label{eq:ek-diag}\\
  a_m\,b_n - b_n\,a_m
  &= \hbar\,(b_m\,a_n - b_{m-1}\,a_{n+1}),
  \label{eq:ek-cross-ab}\\
  [a_m, d_n] &= \hbar\,(c_m\,b_n - c_{m-1}\,b_{n+1}
  - b_m\,c_n + b_{m-1}\,c_{n+1}).
  \label{eq:ek-cross-ad}
\end{align}
The algebra $F_0(R)$ carries a coalgebra structure
$\Delta_F(T(u)) = T(u)\,\dot{\otimes}\,T(u)$ (matrix
multiplication in $\End(V)$, tensor product in $F_0(R)$),
a counit $\varepsilon_F(T(u)) = I$, and an antipode
$S_F(T(u)) = T(u)^{-1}$.

\smallskip
\noindent\emph{Step 2: the state space $V_{\mathrm{EK}}
= U_0(R)$.}
The quantum universal enveloping algebra $U_0(R)$ is the
continuous dual of $F_0(R)$:
\begin{equation}\label{eq:ek-state-space}
  V_{\mathrm{EK}}
  \;:=\; U_0(R)
  \;=\; F_0(R)^*_{\mathrm{cont}},
\end{equation}
a topological Hopf algebra with multiplication dual to the
comultiplication on $F_0(R)$. The vacuum vector is the
counit:
\begin{equation}\label{eq:ek-vacuum}
  \Omega = \varepsilon_F \in U_0(R),
  \qquad
  \langle \Omega, f \rangle = \varepsilon_F(f)
  \quad \text{for all } f \in F_0(R).
\end{equation}
The translation operator $D \colon U_0(R) \to U_0(R)$ is
defined by the shift of spectral parameter:
\begin{equation}\label{eq:ek-translation}
  e^{z D}\, T(u_1) \cdots T(u_n)\, \Omega
  \;=\;
  T(u_1 + z) \cdots T(u_n + z)\, \Omega.
\end{equation}
On the generating modes: $D\,T^{(n)} = n\,T^{(n-1)}$
(with $T^{(-1)} := 0$), reproducing the conformal translation
of mode operators.

For $\fg = \mathfrak{sl}_2$: the graded components of
$U_0(R)$ are finite-dimensional, with
\begin{equation}\label{eq:ek-hilbert}
  \dim U_0(R)_n
  \;=\; \dim \Sym^n(\mathfrak{sl}_2[t])
  \;=\;
  \begin{cases}
  1 & n = 0, \\
  3 & n = 1, \\
  9 & n = 2, \\
  22 & n = 3,
  \end{cases}
\end{equation}
the Hilbert series of the PBW basis
$\{E_r, F_r, H_r : r \geq 0\}$ of
$Y_\hbar(\mathfrak{sl}_2)$, confirming the identification
$U_0(R) \cong Y_\hbar(\mathfrak{sl}_2)$ as vector spaces.

\smallskip
\noindent\emph{Step 3: the vertex operator $Y$.}
The vertex operator
$Y \colon V_{\mathrm{EK}} \to \End(V_{\mathrm{EK}})[\![z, z^{-1}]\!]$
is constructed from the creation--annihilation decomposition
of the quantum current. Define:
\begin{itemize}
\item the \emph{creation operator}
  $\mathbf{T}^+(u, z) = T(u + z)$, acting on
  $V_{\mathrm{EK}}$ by left multiplication;
\item the \emph{annihilation operator}
  $\mathbf{T}^-(u, z) = (T^*(u + z - K\hbar/2))^{-1}$,
  where $T^*$ is the adjoint transfer matrix
  (the transpose of $T$ with respect to the Killing form)
  and $K$ is the level.
\end{itemize}
The quantum current $\mathbf{T}(u, z)$ is the product
\begin{equation}\label{eq:ek-current}
  \mathbf{T}(u, z)
  \;=\;
  \mathbf{T}^+(u, z) \cdot
  \bigl(\mathbf{T}^-(u, z)\bigr)^{-1}
  \;\in\;
  \End(V) \otimes \End(V_{\mathrm{EK}})[\![u^{-1}, z, z^{-1}]\!].
\end{equation}
The vertex operator on generators is
\begin{equation}\label{eq:ek-vertex-gen}
  Y(T^{(n)}\,\Omega,\, z)
  \;=\;
  \Res_{u=0}\, u^n\, \mathbf{T}(u, z)\, du,
\end{equation}
and extends to products by the normal-ordering prescription:
$Y(a\,b,\,z) = {:}\,Y(a,z)\,Y(b,z)\,{:}$, where the
normal ordering places annihilation modes to the right.

On the leading generators ($n = 0$): identifying
$T^{(0)} = I + \hbar(e \otimes e_{12} + f \otimes e_{21}
+ \tfrac{1}{2}h \otimes (e_{11} - e_{22}))$ at leading
order, the vertex operators reduce to the standard affine
Kac--Moody currents
\begin{equation}\label{eq:ek-vertex-leading}
  Y(J^a\,\Omega,\, z)
  \;=\;
  J^a(z) = \sum_{n \in \ZZ} J^a_n\, z^{-n-1}
\end{equation}
for $a \in \{e, f, h\}$, matching the OPE of
$V_k(\mathfrak{sl}_2)$ at leading order in $\hbar$.

\smallskip
\noindent\emph{Step 4: the vertex $R$-matrix $S(z)$.}
The vertex $R$-matrix is the braiding operator that
governs the exchange of vertex operators. It is constructed
from the $R$-matrix data via the universal formula
of~\cite{EK00}:
\begin{equation}\label{eq:ek-vertex-rmatrix}
  S(z)
  \;=\;
  \mathcal{R}(z)
  \cdot \sigma
  \;\in\; \End(V_{\mathrm{EK}} \otimes V_{\mathrm{EK}})[\![z, z^{-1}]\!],
\end{equation}
where $\sigma$ is the graded flip and
$\mathcal{R}(z) = \sum_{n \geq 0} \mathcal{R}_n\, z^{-n}$
is the \emph{spectral universal $R$-matrix} of
$Y_\hbar(\mathfrak{sl}_2)$, satisfying
\begin{equation}\label{eq:ek-rmatrix-coprod}
  \mathcal{R}(z) \cdot \Delta_z(x)
  \;=\;
  \Delta_z^{\mathrm{op}}(x) \cdot \mathcal{R}(z)
  \qquad \text{for all } x \in Y_\hbar(\mathfrak{sl}_2).
\end{equation}
On the fundamental representation $V \otimes V$, the
spectral universal $R$-matrix evaluates to the Yang
$R$-matrix:
\begin{equation}\label{eq:ek-rmatrix-fund}
  (\rho \otimes \rho)\bigl(\mathcal{R}(z)\bigr)
  \;=\; R(z)
  \;=\; z\,I + \hbar\,P
  \;\in\; \End(V \otimes V)[z].
\end{equation}

The leading-order expansion of $S(z)$ on
$V_{\mathrm{EK}} \otimes V_{\mathrm{EK}}$ is
\begin{equation}\label{eq:ek-s-expansion}
  S(z) = \id + \frac{\hbar\,\Omega}{z}
  + \frac{\hbar^2}{z^2}\Bigl(
  \Omega^{(2)} + \tfrac{1}{2}\,\Omega^2
  \Bigr) + O(z^{-3}),
\end{equation}
where $\Omega^{(2)} \in U(\fg)^{\otimes 2}$ is the
second-order Casimir correction (determined by the
$R$-matrix unitarity condition) and
$\Omega^2 = \Omega \cdot \Omega \in U(\fg)^{\otimes 2}$.
At leading order ($z^{-1}$ term), the classical $r$-matrix
$r(z) = \hbar\,\Omega/z$ is recovered.

\medskip
\noindent
\textbf{Part A (continued): verification of the braided VOA
axioms.}
The five axioms of
Definition~\ref{def:e1-chiral-rmatrix} are verified for
$(V_{\mathrm{EK}},\, Y,\, S(z))$ as follows.

\smallskip
\noindent\emph{$S$-locality.}
For the generating currents, the $S$-locality
\begin{equation}\label{eq:ek-s-locality}
  (z - w)^N\, Y(a, z)\, Y(b, w)
  \;=\;
  (z - w)^N\, Y(b, w)\, Y(a, z) \cdot S(z - w)
\end{equation}
holds at $N = 2$: the $\hbar\,\Omega/(z-w)$ correction from
$S(z-w)$ accounts for the structure-constant term in the OPE,
and the higher-order terms account for the double pole.
For descendants, $S$-locality follows by induction on
conformal weight via the shift condition.

\smallskip
\noindent\emph{QYBE, unitarity, shift, hexagon.}
The quantum Yang--Baxter equation on
$V_{\mathrm{EK}}^{\otimes 3}$ reduces on $V^{\otimes 3}$
to~\eqref{eq:yangian-qybe}, verified in
Example~\ref{ex:yangian}; the extension to
$V_{\mathrm{EK}}$ follows from the universal property of
$\mathcal{R}(z)$.
Unitarity $S_{12}(z) \cdot (\sigma \circ S_{12}(-z)
\circ \sigma) = \id$ follows from Drinfeld--Jimbo unitarity.
The shift condition
\begin{equation}\label{eq:ek-shift-cond}
  [D \otimes \id + \id \otimes D,\; S(z)]
  \;=\; \partial_z\, S(z)
\end{equation}
follows from the translation-covariance of the Drinfeld
coproduct. The hexagon follows from the pentagon relation
for $\Phi$ (from the Stasheff boundary of
$\overline{\FM}_4^{\mathrm{ord}}(\CC) = K_4$) combined
with the QYBE for $S$, by the standard argument of
Drinfeld~\cite{EK96} lifted to the vertex-algebraic setting.

\medskip
\noindent
\textbf{Part B: the ordered chiral center of the EK quantum
VOA.}
We compute the ordered chiral Hochschild cohomology
$\HH^{*,\mathrm{ord}}_{\mathrm{ch}}
(V_{\mathrm{EK}},\,V_{\mathrm{EK}})$
at low arities on the formal disk $D$, paralleling the
computation of Example~\ref{ex:yangian} (Part B) but now
using the \emph{full} vertex $R$-matrix $S(z)$ rather than
the Yang $R$-matrix on $V \otimes V$ alone.

\smallskip
\noindent\emph{Arity~$2$: the KZ connection with the full
vertex $R$-matrix.}
At arity~$2$, the KZ connection on
$\Conf_2^{\mathrm{ord}}(\CC)$ with coefficients in
$\End(V_{\mathrm{EK}}^{\otimes 2})$ is
\begin{equation}\label{eq:ek-kz-full}
  \nabla_{\mathrm{KZ}}^S
  = d + \frac{\hbar\,\Omega}{z}\,dz
  + \frac{\hbar^2}{z^2}\,
  \Bigl(\Omega^{(2)} + \tfrac{1}{2}\,\Omega^2\Bigr)\,dz
  + O(z^{-3}),
\end{equation}
where $z = z_1 - z_2$. The connection~\eqref{eq:ek-kz-full}
is obtained from $S(z)$ via the general formula
$\nabla = d + (\partial_z \log S(z))\,dz$ (logarithmic
derivative of the vertex $R$-matrix).

For the classical $V_k(\mathfrak{sl}_2)$ (trivial braiding
$S(z) = \id$), the flat sections are
$\Phi(z) = z^{\Omega/(k+2)}$ with algebraic monodromy.
For the EK quantum VOA ($S(z) \neq \id$), the flat sections
acquire quantum corrections:
\begin{equation}\label{eq:ek-flat-sections}
  \Phi^S(z)
  = z^{\hbar\,\Omega}\,
  \Bigl(I + \frac{\hbar^2\,\Omega^{(2)}}{z}
  + O(z^{-2})\Bigr).
\end{equation}

\emph{The kernel $\ker(\av)$.}
The ordered chiral centre computation parallels
Example~\ref{ex:yangian}; the result is identical on
cohomology. On $V \otimes V$:
$\ker(\av_2)\big|_{V \otimes V}
\cong \bigwedge^2 V = \CC$,
the same one-dimensional space as
equation~\eqref{eq:yangian-ker-av2}.
The quantum corrections from $S(z)$ shift the subleading
terms of the flat sections but do not alter the
$\ker(\av_2)$ dimension. The scalar shadow surviving
$\av_2$ is $\kappa = 3(k+2)/4$, identical to the classical
computation. At arity~$3$, the kernel
$\ker(\av_3) \cong \Lambda^3(\fg) = \CC$ carries the
quantum-corrected Drinfeld associator
\begin{equation}\label{eq:ek-ker-av3}
  \Phi_{\mathrm{KZ}}^S - 1
  \;=\; \hbar^2\,\zeta(2)\,[\Omega_{12}, \Omega_{23}]
  + O(\hbar^3),
\end{equation}
whose leading term agrees with the classical KZ associator.

\medskip
\noindent
\textbf{Part C: $S$-locality as an $\Ainf$ relation.}
The braided VOA axioms of the EK quantum VOA are equivalent
to the chiral $\Ainf$ relations for $V_{\mathrm{EK}}$
through the chiral quantum group equivalence
(Theorem~\ref{thm:chiral-qg-equiv},
(I)$\Leftrightarrow$(II)).
In the ordered bar complex
$\Barord(V_{\mathrm{EK}}) = T^c(s^{-1}\overline{V}_{\mathrm{EK}})$,
the $S$-locality axiom~\eqref{eq:ek-s-locality} governs the
reordering of two insertions: the braid group
$\pi_1(\Conf_2^{\mathrm{ord}}(\CC)) = \ZZ$ acts on degree-$2$
bar chains by
$\sigma_1 \cdot [a | b] = S(z) \cdot [b | a]$.
The $S$-locality axiom is the $n = 2$ Stasheff relation,
the hexagon is the $n = 3$ relation, and the QYBE is the
$n = 4$ relation (from the pentagon for $\Phi$ on
$K_4 = \overline{\FM}_4^{\mathrm{ord}}(\RR)$).

\medskip
\noindent
\textbf{Summary.}
\noindent
The EK quantum VOA $(V_{\mathrm{EK}},\, Y,\, S(z))$
provides the first concrete instance of all three structures
of the chiral quantum group equivalence
(Theorem~\ref{thm:chiral-qg-equiv}) exhibited
simultaneously on a single algebra with all axioms verified.
The ordered chiral center carries the full quantum $R$-matrix
$S(z)$ (not just the classical $r$-matrix
$r(z) = \hbar\,\Omega/z$), and the kernel
$\ker(\av_2) \neq 0$ demonstrates that the ordered theory
is strictly richer than the symmetric Beilinson--Drinfeld
chiral homology even at the quantum level.
\end{example}

\begin{example}[Heisenberg: trivial case]
\label{ex:heis-coproduct}
For the Heisenberg algebra $\cH_k$ (class $G$), all three
structures are trivial:
$S(z) = \id$ (the braiding is symmetric);
$m_k^{\mathrm{ch}} = 0$ for $k \geq 3$ (the algebra is
strictly associative in the chiral sense);
$\Delta^{\mathrm{ch}}(J) = J \otimes 1 + 1 \otimes J$ is
cocommutative and strictly coassociative ($\Phi$ acts
trivially). The ordered and symmetric data coincide.
\end{example}

\begin{example}[$\cW_{1+\infty}$ as a chiral quantum group
via the cohomological Hall algebra]
\label{ex:w-infty-coha}
The algebra $\cW_{1+\infty}[\Psi]$ admits a chiral quantum
group structure originating from the cohomological Hall
algebra (CoHA) of the Jordan quiver $Q_{\mathrm{Jor}}$.
The CoHA $\cH_{\mathrm{Jor}}
= \bigoplus_{n \geq 0}
H^{GL_n}_{\mathrm{BM},*}(\cM_n^{\mathrm{nil}})$
(Kontsevich--Soibelman~\cite{KS11}) carries a Hecke product
and a comultiplication from short exact sequences;
Yang--Zhao~\cite{YZ18a} proved this is a bialgebra with
Drinfeld double $D(\cH_{\mathrm{Jor}}) \cong
Y(\mathfrak{gl}_1)$.
Schiffmann--Mellit--Minets--Vasserot~\cite{SMMV23} identified
$\cH_{\mathrm{Jor}} \cong U(\cW_{1+\infty}[\Psi])^+$.

The three structures of
Theorem~\ref{thm:chiral-qg-equiv} on $\cW_{1+\infty}$ are:
\begin{enumerate}[label=\textup{(\roman*)}]
\item \textup{($R$-matrix.)}
  $R_{\mathrm{CoHA}} = q^{\alpha \cdot \beta}$
  ($q = e^{\pi i \Psi}$), the universal $R$-matrix of
  quantum toroidal $\mathfrak{gl}_1$ restricted to the CoHA;
  classical limit $r(z) = \Psi/z$.
\item \textup{($\Ainf$ structure.)} Higher operations
  $m_k^{\mathrm{ch}}$ for $k \geq 3$ encode the obstruction
  to strict commutativity at successive spins; these lie in
  $\ker(\av_k)$.
\item \textup{(Chiral coproduct.)} The CoHA
  comultiplication gives
  $\Delta^{\mathrm{ch}}(J) = J \otimes 1 + 1 \otimes J$
  (cocommutative at spin~$1$) and
  $\Delta^{\mathrm{ch}}(T) = T \otimes 1 + 1 \otimes T
  + \Psi \cdot (J \otimes J)/z^2 + \cdots$
  (spectral corrections from the Virasoro OPE at spin~$2$).
\end{enumerate}

\noindent\textbf{Descent to $\cW_N$.}
Drinfeld--Sokolov reduction gives
$\pi_N \colon \cW_{1+\infty}[\Psi]
\twoheadrightarrow \cW_N$
for each $N \geq 2$, compatible with the coproduct.
The tower
$\cdots \twoheadrightarrow \cW_{N+1}
\twoheadrightarrow \cW_N
\twoheadrightarrow \cdots
\twoheadrightarrow \cW_2 = \mathrm{Vir}_c$
is a projective system of chiral quantum groups.
At $N = 2$: class $M$, $\kappa = c/2$, irregular KZ
connection. At $N = 3$: $\kappa(\cW_3) = 5c/6$ (where
$H_3 - 1 = 5/6$).

\noindent\textbf{$Y$-algebras and holomorphic gravity.}
The Gaiotto--Rap\v{c}\'ak $Y$-algebras
$Y_{N_1,N_2,N_3}[\Psi]$~\cite{GR17} are truncations of
$\cW_{1+\infty}$; the $S_3$-triality is a symmetry of the
chiral coproduct, corresponding to three Hecke
correspondences on spiked-instanton moduli
(Rap\v{c}\'ak--Soibelman--Yang--Zhao~\cite{RSYZ20}).
The $\cD$-module $\cF_n^{\mathrm{ord}}(\cW_{1+\infty})$
has irregular singularities of unbounded Poincar\'e rank
(spin-$s$ contributes pole order $2s - 1$ after $d$-log
absorption), providing the boundary algebra of higher-spin
holomorphic gravity.
\end{example}


% ----------------------------------------------------------------
\subsection{Structural consequences}
\label{subsec:structural-consequences}

The three equivalent descriptions of an $\Eone$-chiral
algebra all project from a single datum: the ordered bar
complex $\Barord(\cA) = T^c(s^{-1}\bar{\cA})$ with its
deconcatenation coproduct. The bar differential encodes
the $\Ainf$ operations, the deconcatenation induces the
chiral coproduct, and the KZ monodromy recovers the vertex
$R$-matrix. This is the structural content of $\Eone$
primacy.

\begin{corollary}[The ordered bar complex encodes all three
structures]
\label{cor:bar-encodes-all}
The ordered bar complex $\Barord(\cA) = T^c(s^{-1}\bar{\cA})$,
equipped with its bar differential $d_{\Barord}$ and
deconcatenation coproduct $\Delta$, is the universal
object from which all three structures in
Theorem~\textup{\ref{thm:chiral-qg-equiv}} are recovered:
\begin{enumerate}[label=\textup{(\roman*)}]
\item The bar differential $d_{\Barord}$ encodes the
  $\Ainf$ structure maps $m_k^{\mathrm{ch}}$
  \textup{(}by the standard bar-cobar correspondence:
  coderivations of $T^c(s^{-1}\bar{\cA})$ biject with
  elements of $\prod_k \End^{\mathrm{ch}}_\cA(k)$\textup{)}.
\item The deconcatenation coproduct $\Delta$ induces,
  via cobar dualization, the chiral coproduct
  $\Delta^{\mathrm{ch}}$.
\item The vertex $R$-matrix $S(z)$ is recovered as the
  monodromy of the KZ connection, which is itself the
  arity-$2$ projection of the Maurer--Cartan element
  $\Theta_\cA$ in $\gEone$.
\end{enumerate}
This is the structural content of the $\Eone$ primacy
principle: the ordered bar complex with deconcatenation
is the primitive object; the $R$-matrix, the chiral
coproduct, and the $\Ainf$ operations are projections of
this single datum.
\end{corollary}

\begin{remark}[Relation to factorization quantum groups]
\label{rem:factorization-qg}
When the hexagon holds strictly, the chiral quantum group
$(\cA, S(z), \Delta^{\mathrm{ch}})$ is a
\emph{factorization quantum group}
(Latyntsev~\cite{Latyntsev23}): its ordered bar complex
defines a factorization coalgebra on the curve. When the
hexagon holds only up to $\Phi$, one obtains a
\emph{quasi-factorization quantum group} with regularized
holonomy data $(R, \Phi)$.
\end{remark}

% ================================================================
\section{Ordered chiral homology on elliptic curves:
$\int_{E_\tau}^{\mathrm{ord}} Y_\hbar(\mathfrak{sl}_2)$}
\label{sec:elliptic-ordered}

The preceding sections develop ordered chiral homology at
genus~$0$: the base curve is $\PP^1$ (or the formal disk
$D$), the bar propagator is $d\log(z_1 - z_2)$, and the KZ
connection governs the flat transport. The five
explicit examples (\S\ref{sec:examples}) treat chiral
algebras from classes $G$, $L$, $C$, $M$ and the genuinely
$\Eone$ Yangian on $D^\times$.

The purpose of this section is to construct the ordered
chiral homology of the Yangian
$Y_\hbar(\mathfrak{sl}_2)$ on an elliptic curve
$E_\tau = \CC / (\ZZ + \tau\ZZ)$. This is a genuinely new
invariant: (i) the base curve has genus~$1$, so the bar
propagator acquires the elliptic function
$\wp_1 = \theta_1'/\theta_1$ in place of $1/z$; (ii) the algebra is genuinely $\Eone$,
so the ordered theory differs from the symmetric; (iii)
the KZ connection is replaced by the KZB connection, and
the $B$-cycle monodromy produces the quantum group
parameter $q = e^{2\pi i \hbar}$.

The construction proceeds arity by arity.

% ----------------------------------------------------------------
\subsection{Setup: the ordered chiral chain complex on
$E_\tau$}
\label{subsec:elliptic-setup}

Fix $\tau \in \HHH = \{\mathrm{Im}(\tau) > 0\}$ and let
$E_\tau = \CC/(\ZZ + \tau\ZZ)$ with coordinate $z$. The
ordered configuration space at arity~$n$ is
\begin{equation}\label{eq:conf-elliptic}
  \Conf_n^{\mathrm{ord}}(E_\tau)
  = \{(z_1, \ldots, z_n) \in E_\tau^n :
  z_i \neq z_j \text{ for } i \neq j\}.
\end{equation}
No ordering on $E_\tau$ itself is needed: ``ordered'' means
the labelling of points is retained (no $\Sigma_n$-quotient).

The bar propagator on $E_\tau$ is
\begin{equation}\label{eq:bar-prop-ell}
  \eta^{(1)}_{ij}
  = d\log E(z_i, z_j)
  = \wp_1(z_i - z_j, \tau)\,d(z_i - z_j),
\end{equation}
where $E(z,w)$ is the prime form on $E_\tau$ and
$\wp_1(z, \tau) = \theta_1'(z|\tau)/\theta_1(z|\tau)$
is the logarithmic derivative of the Jacobi theta function. This is weight~$1$ (AP27: the bar propagator
is always weight~$1$).

The ordered chiral chain complex at arity~$n$ on $E_\tau$
is
\begin{equation}\label{eq:ell-chiral-cx}
  \cC_n^{\mathrm{ord}}(E_\tau, \cA)
  = \Bigl(
  \Omega^*_{\log}\bigl(
  \overline{\FM}_n^{\mathrm{ord}}(E_\tau)\bigr)
  \otimes (s^{-1}\bar{\cA})^{\otimes n},
  \;\; d_{\mathrm{total}}
  \Bigr)
\end{equation}
with
\begin{equation}\label{eq:d-total-ell}
  d_{\mathrm{total}}
  = d_{\mathrm{dR}}
  + d_{\mathrm{bar}}
  + \nabla_{\mathrm{KZB}}.
\end{equation}
The decisive change from genus~$0$ to genus~$1$: the
KZ connection $\nabla_{\mathrm{KZ}} = d + \sum r_{ij}\,dz_{ij}$
is replaced by the KZB connection
\begin{equation}\label{eq:kzb-ell}
  \nabla_{\mathrm{KZB}}
  = d
  - \frac{1}{k + h^\vee}
  \sum_{i < j} \Omega_{ij}\,
  \wp_1(z_{ij}, \tau)\,dz_{ij}
  - \frac{1}{2\pi i(k + h^\vee)}
  \sum_{i < j} \wp(z_{ij}, \tau)\,\Omega_{ij}\,d\tau,
\end{equation}
where $\hbar = 1/(k + h^\vee) = 1/(k+2)$ for
$\mathfrak{sl}_2$
\textup{(}the same convention as
Remark~\textup{\ref{rem:hbar-convention}}\textup{)}.
The $d\tau$ term is absent at genus~$0$; it encodes the
curvature $\kappa(\cA) \cdot \omega_1$ from the Hodge class
$\omega_1 = c_1(\lambda)$ on $\overline{\cM}_{1,n}$.

\begin{definition}[Ordered chiral homology on $E_\tau$]
\label{def:ord-ch-ell}
The \emph{ordered chiral homology of $\cA$ on $E_\tau$} is
\begin{equation}\label{eq:ord-ch-ell}
  \int_{E_\tau}^{\mathrm{ord}} \cA
  \;:=\;
  R\Gamma_{\mathrm{dR}}\bigl(
  \Ran^{\mathrm{ord}}(E_\tau),\;
  \cF^{\mathrm{ord}}(\cA)\bigr)
  \;=\;
  \bigoplus_{n \geq 0}
  H^*\bigl(
  \cC_n^{\mathrm{ord}}(E_\tau, \cA),\;
  d_{\mathrm{total}}\bigr).
\end{equation}
\end{definition}

% ----------------------------------------------------------------
\subsection{Arity $0$: the center}
\label{subsec:ell-arity0}

At arity~$0$ there are no marked points and no
configuration space; the ordered chiral homology reduces
to the chiral center of~$\cA$.

\begin{proposition}[Arity-$0$ ordered chiral homology on
$E_\tau$]
\label{prop:ell-arity0}
For $\cA = Y_\hbar(\mathfrak{sl}_2)$:
\begin{equation}\label{eq:ell-arity0}
  H^*\bigl(\cC_0^{\mathrm{ord}}(E_\tau,
  Y_\hbar(\mathfrak{sl}_2))\bigr)
  = Z(Y_\hbar(\mathfrak{sl}_2))
  = \CC[\operatorname{qdet}\,T(u)].
\end{equation}
The quantum determinant
$\operatorname{qdet}\,T(u) = A(u)\,D(u - \hbar)
- B(u)\,C(u - \hbar)$ generates the center as a
polynomial algebra \textup{(}Molev\textup{)}.
The averaging map $\av_0$ is an isomorphism
\textup{(}there is nothing to symmetrise at arity~$0$\textup{)},
so $\ker(\av_0) = 0$.
\end{proposition}

\begin{remark}
The center is independent of the base curve: it depends
only on the algebra $\cA$ and not on $E_\tau$. This reflects
the arity-$0$ factorisation $\cD$-module being a skyscraper at
the empty configuration.
\end{remark}

% ----------------------------------------------------------------
\subsection{Arity $1$: de~Rham cohomology with Yangian
coefficients}
\label{subsec:ell-arity1}

At arity~$1$,
$\Conf_1^{\mathrm{ord}}(E_\tau) = E_\tau$ and the ordered
chiral chain complex is
\begin{equation}\label{eq:ell-arity1-cx}
  \cC_1^{\mathrm{ord}}(E_\tau, Y_\hbar)
  = \Bigl(\Omega^*(E_\tau) \otimes s^{-1}\bar{Y}_\hbar,
  \;\; d_{\mathrm{dR}} + d_{\mathrm{bar}}\Bigr),
\end{equation}
where $\bar{Y}_\hbar = \ker(\varepsilon)$ is the
augmentation ideal of $Y_\hbar(\mathfrak{sl}_2)$ and
$|s^{-1}v| = |v| - 1$.
At arity~$1$ the KZB connection is trivial (no pairs
$i < j$), so $\nabla_{\mathrm{KZB}} = d$, absorbed into
$d_{\mathrm{dR}}$. The bar differential $d_{\mathrm{bar}}$
acts on each factor of $s^{-1}\bar{Y}_\hbar$ via the
$\Ainf$ structure maps $m_k$ of the chiral algebra.

The de~Rham complex of $E_\tau$ has cohomology
\begin{equation}\label{eq:dr-etau}
  H^0(E_\tau) = \CC, \qquad
  H^1(E_\tau) = \CC^2, \qquad
  H^2(E_\tau) = \CC,
\end{equation}
generated respectively by $1$; the $A$-cycle form
$dz$ and the $B$-cycle form
$d\bar{z}/\mathrm{Im}(\tau)$; and the volume form
$dz \wedge d\bar{z}/(2i\,\mathrm{Im}(\tau))$.
By the K\"unneth theorem:
\begin{equation}\label{eq:ell-arity1-kunneth}
  H^*\bigl(\cC_1^{\mathrm{ord}}\bigr)
  = H^*(E_\tau) \otimes H^*(s^{-1}\bar{Y}_\hbar,
  d_{\mathrm{bar}}).
\end{equation}
The bar cohomology $H^*(s^{-1}\bar{Y}_\hbar, d_{\mathrm{bar}})$
is the arity-$1$ bar cohomology of the chiral algebra.
For $Y_\hbar(\mathfrak{sl}_2)$ viewed through its
RTT presentation, the bar differential at arity~$1$ is
the internal differential of the Yangian (the differential
inherited from the dg structure of the ordered bar
complex of $V_k(\mathfrak{sl}_2)$ at the first page of
the bar spectral sequence). Since $Y_\hbar(\mathfrak{sl}_2)$
is the Koszul dual of $V_k(\mathfrak{sl}_2)$ and the
affine algebra is Koszul (class~$L$), the bar cohomology
at arity~$1$ is concentrated in degree~$1$:
\begin{equation}\label{eq:bar-arity1-yangian}
  H^*(s^{-1}\bar{Y}_\hbar, d_{\mathrm{bar}})
  = s^{-1}\mathfrak{sl}_2
  \;\;\text{(concentrated in cohomological degree $1$).}
\end{equation}
This is the desuspended Lie algebra: the three generators
$\mathcal{E}^{(0)}, \mathcal{F}^{(0)}, H^{(0)}$ in mode
zero, shifted by $s^{-1}$.

\begin{proposition}[Arity-$1$ ordered chiral homology on
$E_\tau$]
\label{prop:ell-arity1}
\begin{equation}\label{eq:ell-arity1-result}
  H^*\bigl(\cC_1^{\mathrm{ord}}(E_\tau,
  Y_\hbar(\mathfrak{sl}_2))\bigr)
  \cong H^*(E_\tau) \otimes s^{-1}\mathfrak{sl}_2
  \cong (s^{-1}\mathfrak{sl}_2)^{\oplus 4},
\end{equation}
with the four copies sitting in cohomological degrees
$1, 2, 2, 3$ \textup{(}from the tensor product of
$H^{0,1,1,2}(E_\tau)$ with the degree-$1$ bar cohomology
$s^{-1}\mathfrak{sl}_2$\textup{)}. The total dimension is
$4 \cdot 3 = 12$.
\end{proposition}

\begin{remark}
At arity~$1$, $\Sigma_1 = \{e\}$, so the averaging map
is the identity and $\ker(\av_1) = 0$. The ordered and
symmetric theories agree at arity~$1$: the non-trivial
ordered data enters only at arity $\geq 2$.
\end{remark}

\begin{remark}[Comparison with genus~$0$]
At genus~$0$, $H^*(D^\times)$ contributes only in degree~$0$
(and degree~$1$ from the circle at infinity), so the
arity-$1$ cohomology is $3$-dimensional. The passage to
genus~$1$ multiplies by $\dim H^*(E_\tau) = 4$: the new
$B$-cycle produces the extra copy.
\end{remark}

% ----------------------------------------------------------------
\subsection{Arity $2$: the elliptic KZB local system}
\label{subsec:ell-arity2}

At arity~$2$, the configuration space is
$\Conf_2^{\mathrm{ord}}(E_\tau) = E_\tau^2 \setminus
\Delta$, and the KZB connection acts non-trivially.
The ordered chiral chain complex is
\begin{equation}\label{eq:ell-arity2-cx}
  \cC_2^{\mathrm{ord}}(E_\tau, Y_\hbar)
  = \Bigl(
  \Omega^*_{\log}\bigl(
  \overline{\FM}_2^{\mathrm{ord}}(E_\tau)\bigr)
  \otimes (s^{-1}\bar{Y}_\hbar)^{\otimes 2},
  \;\;
  d_{\mathrm{dR}} + d_{\mathrm{bar}}
  + \nabla_{\mathrm{KZB}}
  \Bigr).
\end{equation}
The key new ingredient is the elliptic collision residue:
the arity-$2$ classical $r$-matrix on $E_\tau$ is the
Belavin elliptic $r$-matrix
% AP1: kappa(V_k(sl_2)) = dim(sl_2)(k+h^v)/(2h^v)
%   = 3(k+2)/4.  k=0 -> 3/2. k=-2 -> 0.
% PE-1: r-matrix. Family: affine KM sl_2.
%   r(z,tau) = k * r^{Belavin}(z,tau).
%   Level parameter: k. Convention: trace-form.
%   k=0 check: r(z,tau)|_{k=0} = 0. Match: Y.
%   Source: elliptic_rmatrix_shadow.py:genus1_shadow_rmatrix_sl2
%   verdict: ACCEPT
\begin{equation}\label{eq:belavin-r}
  r^{\mathrm{ell}}(z, \tau)
  = k \cdot r^{\mathrm{Belavin}}(z, \tau)
  = k \Bigl(
  \frac{\wp_1(z, \tau)}{2}\,H \otimes H
  + \phi_+(z, \tau)\,E \otimes F
  + \phi_-(z, \tau)\,F \otimes E
  \Bigr),
\end{equation}
where $\phi_\pm(z, \tau) = \theta_1'(0|\tau)\,
\theta_1(z \pm \tfrac{1}{2}|\tau) /
(\theta_1(z|\tau)\,\theta_1(\pm\tfrac{1}{2}|\tau))$
are Belavin's $\phi$-functions. At $k = 0$,
$r^{\mathrm{ell}} = 0$ (the algebra becomes abelian).
The degeneration
$\tau \to i\infty$ sends
$\wp_1(z, \tau) \to \pi\cot(\pi z) \to 1/z + O(z)$
and $\phi_\pm \to 1/z + O(1)$, recovering
$r^{\mathrm{rat}}(z) = k\,\Omega/z$.

The KZB connection at arity~$2$ on the relative coordinate
$w = z_1 - z_2$ is
\textup{(}Bernard~\cite{Bernard88};
Calaque--Enriquez--Etingof~\cite{CEE09}\textup{)}
\begin{equation}\label{eq:kzb-arity2}
  \nabla_{\mathrm{KZB}}^{(2)}
  = d
  - \frac{1}{k + h^\vee}\,
  \Omega\,\wp_1(w, \tau)\,dw
  - \frac{1}{2\pi i(k + h^\vee)}\,
  \wp(w, \tau)\,\Omega\,d\tau.
\end{equation}
The flat sections satisfy
\begin{align}
  \frac{\partial \Phi}{\partial w}
  &= \frac{1}{k + h^\vee}\,
  \wp_1(w, \tau)\,\Omega\,\Phi,
  \label{eq:kzb-w} \\
  \frac{\partial \Phi}{\partial \tau}
  &= \frac{1}{2\pi i(k + h^\vee)}\,
  \wp(w, \tau)\,\Omega\,\Phi.
  \label{eq:kzb-tau}
\end{align}
The flatness of the KZB system follows from the heat
equation for theta functions:
$\partial_\tau\,\wp_1(z, \tau)
= \frac{1}{2\pi i}\,\partial_z\,\wp(z, \tau)$.

\smallskip
\noindent
\textbf{Decomposition by $\mathfrak{sl}_2$-isospin.}
On the fundamental representation $V = \CC^2$, the
Casimir decomposes as before:
\[
  V \otimes V
  = \operatorname{Sym}^2 V \oplus \bigwedge\nolimits^2 V,
  \qquad
  \Omega\big|_{\operatorname{Sym}^2} = +\tfrac{1}{2},
  \quad
  \Omega\big|_{\bigwedge^2} = -\tfrac{3}{2}.
\]
The KZB system decomposes into scalar equations on each
isospin channel. On the symmetric channel:
\begin{equation}\label{eq:kzb-sym-channel}
  \frac{\partial \Phi_{\mathrm{sym}}}{\partial w}
  = \frac{1}{2(k + h^\vee)}\,\wp_1(w, \tau)\,
  \Phi_{\mathrm{sym}}.
\end{equation}

\smallskip
\noindent
\textbf{Monodromy.}
The function $\wp_1(w, \tau) = \theta_1'(w|\tau)/\theta_1(w|\tau)$
is periodic under $w \mapsto w + 1$ and quasi-periodic
under $w \mapsto w + \tau$ with increment $-2\pi i$.
The puncture monodromy (small loop around $w = 0$)
captures the residue:
\begin{equation}\label{eq:puncture-monodromy-ell}
  M_\gamma
  = \exp\!\bigl(2\pi i\,\hbar\,\Omega\bigr)
  \quad \text{where } \hbar = \frac{1}{k + h^\vee}.
\end{equation}
On each channel:
$M_\gamma^{\mathrm{sym}} = e^{\pi i/(k+2)}$,
$M_\gamma^{\mathrm{alt}} = e^{-3\pi i/(k+2)}$.

The $B$-cycle monodromy ($w \mapsto w + \tau$) is the
genuinely new genus-$1$ datum.
The quasi-periodicity $\wp_1(w + \tau) = \wp_1(w) - 2\pi i$
produces:
\begin{equation}\label{eq:b-monodromy-ell}
  M_B
  = \exp\!\bigl(-2\pi i\,\hbar\,\Omega\bigr)
  = M_\gamma^{-1}.
\end{equation}
\begin{remark}[Dynamical parameter]
\label{rem:dynamical-parameter}
The $B$-cycle monodromy~\eqref{eq:b-monodromy-ell}
is the leading-order result from the quasi-periodicity
of $\wp_1$. The full genus-$1$ monodromy representation
involves the \emph{dynamical parameter}
$\lambda \in \fh^*$ \textup{(}Felder~\cite{Felder94};
Felder--Varchenko~\cite{FelderVarchenko96}\textup{)}: the
$B$-cycle monodromy operator depends on the basepoint
through a shift $\lambda \mapsto \lambda - h^{(i)}$,
producing the dynamical Yang--Baxter equation rather than
the ordinary YBE
\textup{(}Etingof--Varchenko~\cite{EV98};
see also~\cite[Thm~5.6]{Lorgat26I}\textup{)}.
The gauge-invariant content is the constraint
$M_\gamma \cdot [M_A, M_B] = \id$
from $\pi_1(E_\tau \setminus \{0\})$.
\end{remark}
On each isospin channel:
\begin{equation}\label{eq:b-monodromy-channels}
  M_B^{\mathrm{sym}} = e^{-\pi i/(k+2)}, \qquad
  M_B^{\mathrm{alt}} = e^{3\pi i/(k+2)}.
\end{equation}

\begin{proposition}[Arity-$2$ ordered chiral homology on
$E_\tau$]
\label{prop:ell-arity2}
The arity-$2$ ordered chiral homology
$H^*(\cC_2^{\mathrm{ord}}(E_\tau, Y_\hbar))$ is the de~Rham
cohomology of the KZB local system on
$\Conf_2^{\mathrm{ord}}(E_\tau)$ with coefficients in
$(s^{-1}\bar{Y}_\hbar)^{\otimes 2}$.
On evaluation modules $V_a \otimes V_b$:
\begin{enumerate}[label=\textup{(\roman*)}]
\item The KZB local system has monodromy
  $(R, q^\Omega)$ around the $A$- and $B$-cycles.
\item The de~Rham cohomology of the KZB local system on
  $E_\tau \setminus \{0\}$ \textup{(}equivalently,
  $\Conf_2^{\mathrm{ord}}(E_\tau)$ in the relative
  coordinate $w = z_1 - z_2$\textup{)} is:
  \begin{equation}\label{eq:ell-arity2-dr}
    H^p_{\mathrm{dR}}(E_\tau \setminus \{0\},\,
    \nabla_{\mathrm{KZB}})
    = \begin{cases}
      \ker(R - I) \cap \ker(q^\Omega - I)
      & p = 0, \\
      \text{\textup{(}see~\eqref{eq:ell-arity2-h1}
      below\textup{)}}
      & p = 1.
    \end{cases}
  \end{equation}
\item For generic $\hbar$ \textup{(}equivalently, generic
  level $k$\textup{)}, $R$ and $q^\Omega$ have distinct
  eigenvalues on each isospin channel, and
  $H^0_{\mathrm{dR}} = 0$: there are no global flat
  sections.
\item The $H^1$ is computed by the Euler characteristic
  of the local system. By the Riemann--Hilbert
  correspondence and the Euler characteristic formula for
  a rank-$d^2$ local system on a genus-$1$ curve with one
  puncture:
  \begin{equation}\label{eq:ell-arity2-h1}
    \dim H^1_{\mathrm{dR}}(E_\tau \setminus \{0\},\,
    \nabla_{\mathrm{KZB}})
    = d^2 \cdot (2 \cdot 1 - 2 + 1)
    = d^2,
  \end{equation}
  where $d = \dim V = 2$, so $\dim H^1 = 4$.
\end{enumerate}
\end{proposition}

\begin{proof}
The Euler characteristic of a local system of rank~$r$ on a
surface of genus~$g$ with $s$ punctures is
$\chi = r(2 - 2g - s)$. For $g = 1$, $s = 1$, $r = d^2 = 4$:
$\chi = 4(2 - 2 - 1) = -4$. Since $H^0 = 0$ for generic
$\hbar$ (the monodromy has no invariant vectors) and
$H^2 = 0$ (the puncture forces vanishing), we have
$\dim H^1 = 4$. The four-dimensional space decomposes as
$3 + 1$ along the isospin channels $\operatorname{Sym}^2 V$
(dimension~$3$) and $\bigwedge^2 V$ (dimension~$1$).
\end{proof}

\begin{remark}[The symmetric shadow]
\label{rem:ell-arity2-shadow}
The averaging map $\av_2$ at arity~$2$ sends the
ordered de~Rham cohomology to $\Sigma_2$-coinvariants.
The $\Sigma_2$-action is the composition of the geometric
swap $w \mapsto -w$ with the $R$-matrix. On the symmetric
channel, the $R$-matrix eigenvalue is
$+\hbar/2 + O(\hbar^2)$; on the antisymmetric channel,
$-3\hbar/2 + O(\hbar^2)$.
% AP1: kappa(V_k(sl_2)) = 3(k+2)/4.
The scalar surviving $\av_2$ is
$\kappa = 3(k + 2)/4$.

The kernel $\ker(\av_2)$ at genus~$1$ is richer than at
genus~$0$: the $B$-cycle monodromy contributes an
additional antisymmetric class (from the
$e^{3\pi i/(k+2)}$ eigenvalue of $M_B$ on
$\bigwedge^2 V$) that is invisible to the
symmetric theory.
\end{remark}

% ----------------------------------------------------------------
\subsection{Arity $n$: elliptic braid group monodromy}
\label{subsec:ell-arity-n}

At arity~$n$, the fundamental group of
$\Conf_n^{\mathrm{ord}}(E_\tau)$ is the \emph{elliptic
braid group} $\mathrm{EB}_n$, which is an extension
\begin{equation}\label{eq:ell-braid}
  1 \to PB_n \to \mathrm{EB}_n \to
  \pi_1(E_\tau)^n \to 1,
\end{equation}
where $PB_n$ is the pure braid group on $n$ strands and
$\pi_1(E_\tau) \cong \ZZ^2$ contributes the $A$- and
$B$-cycle monodromies for each point.

The KZB monodromy representation at arity~$n$ is
\begin{equation}\label{eq:kzb-monodromy-n}
  \rho_n^{\mathrm{KZB}} \colon
  \mathrm{EB}_n \to
  \operatorname{GL}(V^{\otimes n}).
\end{equation}
The ordered chiral homology at arity~$n$ is the de~Rham
cohomology of the KZB local system on the Fulton--MacPherson
compactification:
\begin{equation}\label{eq:ell-arity-n}
  \int_{E_\tau, n}^{\mathrm{ord}}
  Y_\hbar(\mathfrak{sl}_2)
  = H^*_{\mathrm{dR}}\bigl(
  \overline{\FM}_n^{\mathrm{ord}}(E_\tau),\;
  \nabla_{\mathrm{KZB}}^{(n)}\bigr)
  \otimes (s^{-1}\bar{Y}_\hbar)^{\otimes n}.
\end{equation}

\begin{proposition}[Euler characteristic at low arities]
\label{prop:ell-euler-n}
The twisted Euler characteristics of the KZB local system on
$\Conf_n^{\mathrm{ord}}(E_\tau)$, with fibre $V^{\otimes n}$
\textup{(}$V = \CC^2$\textup{)} for generic $\hbar$, are:
\begin{align}
  n = 0&: \quad \chi = \infty
  \;\;\text{\textup{(}$\CC[\operatorname{qdet}]$ is
  infinite-dimensional\textup{)}},
  \label{eq:euler-n0} \\
  n = 1&: \quad \chi = 0
  \;\;\text{\textup{(}trivial local system on $E_\tau$,
  $\chi(E_\tau) = 0$\textup{)}},
  \label{eq:euler-n1} \\
  n = 2&: \quad \chi = -4
  \;\;\text{\textup{(}rank $4$ on $E_\tau \setminus
  \{0\}$\textup{)}}.
  \label{eq:euler-n2}
\end{align}
\end{proposition}

\begin{proof}
At $n = 1$: $\Conf_1(E_\tau) = E_\tau$ and the KZB
connection has no pairwise terms (there are no pairs
$i < j$), so the local system is trivial of rank $d = 2$.
The twisted Euler characteristic is
$\chi = d \cdot \chi(E_\tau) = 2 \cdot 0 = 0$, consistent
with $\dim H^0 = 2$, $\dim H^1 = 4$, $\dim H^2 = 2$
from the K\"unneth computation: $\chi = 2 - 4 + 2 = 0$.

At $n = 2$: the relative coordinate $w = z_1 - z_2$
identifies $\Conf_2^{\mathrm{ord}}(E_\tau)$ (modulo the
$E_\tau$-translation in the centre-of-mass) with
$E_\tau \setminus \{0\}$, a genus-$1$ surface with one
puncture. The KZB local system has rank $d^2 = 4$. The
Euler characteristic of a rank-$r$ local system on a
surface of genus~$g$ with $s$ punctures is
$\chi = r(2 - 2g - s)$: here
$\chi = 4(2 - 2 - 1) = -4$. For generic $\hbar$, the
monodromy has no invariant vectors and
$H^0 = 0$; the puncture forces $H^2 = 0$; so
$\dim H^1 = 4$. This is consistent with
Proposition~\ref{prop:ell-arity2}.
\end{proof}

\begin{remark}[Higher arities]
For $n \geq 3$, $\Conf_n^{\mathrm{ord}}(E_\tau)$ is a
smooth quasi-projective variety of complex dimension~$n$.
Its topological Euler characteristic vanishes:
$\chi(\Conf_n^{\mathrm{ord}}(E_\tau)) = 0$
for all $n \geq 1$ (the product formula
$\chi = \prod_{j=0}^{n-1}(\chi(E_\tau) - j) = 0$ since
$\chi(E_\tau) = 0$). The twisted Euler characteristic of the
KZB local system is also zero:
$\chi(\Conf_n^{\mathrm{ord}}(E_\tau),\,
\cL_{\mathrm{KZB}}) = 0$ for $n \geq 3$.
This does \emph{not} mean the de~Rham cohomology vanishes;
it means the dimensions of the cohomology groups cancel
in alternating sum. The individual Betti numbers with
KZB coefficients require explicit computation of the
elliptic braid group monodromy at each arity.
\end{remark}

% ----------------------------------------------------------------
\subsection{The total ordered chiral homology and the
averaging kernel}
\label{subsec:ell-total}

Assembling all arities, the total ordered chiral homology is
\begin{equation}\label{eq:ell-total}
  \int_{E_\tau}^{\mathrm{ord}}
  Y_\hbar(\mathfrak{sl}_2)
  = \bigoplus_{n \geq 0}
  \int_{E_\tau, n}^{\mathrm{ord}}
  Y_\hbar(\mathfrak{sl}_2).
\end{equation}
At each arity, the averaging map
$\av_n \colon \int_{E_\tau, n}^{\mathrm{ord}}
\to (\int_{E_\tau, n}^{\mathrm{ord}})_{\Sigma_n}$
projects to the symmetric shadow. The kernel
$\ker(\av) = \bigoplus_n \ker(\av_n)$ carries the
genuinely $\Eone$ data that is invisible to the symmetric
theory.

\begin{theorem}[Structure of
$\int_{E_\tau}^{\mathrm{ord}}
Y_\hbar(\mathfrak{sl}_2)$]
\label{thm:ell-ordered-ch}
For generic $\hbar = 1/(k+2)$ with $k \notin \{-2\}
\cup \ZZ_{\leq 0}$:
\begin{enumerate}[label=\textup{(\roman*)}]
\item \textup{(Arity decomposition.)}
  The arity-by-arity structure is:
  \begin{center}
  \renewcommand{\arraystretch}{1.3}
  \begin{tabular}{@{} c l l l @{}}
  \toprule
  $n$ & $\int_{E_\tau, n}^{\mathrm{ord}}
  Y_\hbar$ & $\ker(\av_n)$ & Shadow \\
  \midrule
  $0$ & $\CC[\operatorname{qdet}\,T(u)]$
    & $0$ & $\CC[\operatorname{qdet}]$ \\
  $1$ & $(s^{-1}\mathfrak{sl}_2)^{\oplus 4}$,
    dim $12$ & $0$
    & $(s^{-1}\mathfrak{sl}_2)^{\oplus 4}$ \\
  $2$ & KZB flat sections, dim $4$
    & $\bigwedge^2 V \oplus$ $B$-cycle class
    & $\kappa = 3(k{+}2)/4$ \\
  $3$ & elliptic braid monodromy
    & $\supseteq \Lambda^3(\fg)$
    & traces \\
  $n$ & $\mathrm{EB}_n$-monodromy of KZB
    & matrix-valued
    & scalar traces \\
  \bottomrule
  \end{tabular}
  \end{center}

\item \textup{(Genus-$1$ enhancement of the kernel.)}
  The kernel $\ker(\av_2)$ at genus~$1$ is strictly larger
  than the kernel at genus~$0$: the $B$-cycle monodromy
  $q^{\Omega}$ introduces new antisymmetric classes in the
  $\bigwedge^2 V$ channel that have no genus-$0$ analogue.
  The Drinfeld associator at arity~$3$ is deformed by the
  KZB associator $\Phi_{\mathrm{KZB}}$, which is a genus-$1$
  deformation of $\Phi_{\mathrm{KZ}}$.

\item \textup{(Degeneration.)}
  In the limit $\tau \to i\infty$
  \textup{(}the $B$-cycle collapses and $E_\tau$ degenerates to
  $\CC^\times$\textup{)}, the KZB connection
  degenerates to the KZ connection, the $B$-cycle monodromy
  becomes trivial, and
  \begin{equation}\label{eq:ell-degen}
    \lim_{\tau \to i\infty}
    \int_{E_\tau}^{\mathrm{ord}}
    Y_\hbar(\mathfrak{sl}_2)
    = \int_{D^\times}^{\mathrm{ord}}
    Y_\hbar(\mathfrak{sl}_2),
  \end{equation}
  recovering the genus-$0$ ordered chiral homology of
  \S\textup{\ref{sec:examples}}.

\item \textup{(Modular dependence.)}
  The ordered chiral homology varies with $\tau$ as a section
  of a vector bundle over $\HHH$ with flat connection
  $\nabla_\tau = \partial_\tau - \frac{\hbar}{2\pi i}
  \sum_{i<j} \wp(z_{ij}, \tau)\,\Omega_{ij}$.
  The $\mathrm{SL}_2(\ZZ)$-action on $\tau$ acts on the
  ordered chiral homology through the modular transformation
  of the KZB connection. The kernel $\ker(\av)$ transforms as
  a non-trivial representation of the modular group: the
  $B$-cycle monodromy is exchanged with the $A$-cycle
  monodromy under $\tau \mapsto -1/\tau$: the pair
  $(M_\gamma, M_B) = (q^\Omega, q^{-\Omega})$ is interchanged.
  The parameter $q = e^{2\pi i \hbar}$ itself is
  $\tau$-independent; the modular group acts on the
  monodromy representation, not on $q$.
\end{enumerate}
\end{theorem}

\begin{proof}[Proof sketch]
Parts~(i)--(ii) follow from the arity-by-arity computation
of \S\S\ref{subsec:ell-arity0}--\ref{subsec:ell-arity-n}.

Part~(iii): the Belavin $r$-matrix degenerates to the
rational $r$-matrix as $\tau \to i\infty$
($\wp_1(z, \tau) = \theta_1'(z|\tau)/\theta_1(z|\tau)
\to 1/z$ and $\phi_\pm(z, \tau) \to 1/z$), the $B$-cycle shrinks,
and the KZB flat sections degenerate to KZ flat sections.
The arity-by-arity Euler characteristics are continuous under
this degeneration.

Part~(iv): the modular transformation of the KZB connection is
the standard result of Bernard~\cite{Bernard88}: the $d\tau$
term ensures compatibility with the $\mathrm{SL}_2(\ZZ)$-action,
and the flat sections transform with the modular anomaly
proportional to $\kappa(\cA) = 3(k+2)/4$.
\end{proof}

\begin{remark}[Comparison with the quantum group]
\label{rem:ell-quantum-group}
The object
$\int_{E_\tau}^{\mathrm{ord}} Y_\hbar(\mathfrak{sl}_2)$
computes the conformal blocks of the Yangian on $E_\tau$.
By the Yangian--quantum group deformation
\cite[Theorem~5.5.1]{Lorgat26I},
the genus-$1$ line-operator algebra is
$U_q(\mathfrak{sl}_2)$ with $q = e^{2\pi i \hbar}$.
The ordered chiral homology packages the full KZB monodromy
data into a single invariant: it sees both the $R$-matrix
($A$-cycle, Yangian data) and the quantum group parameter
($B$-cycle, $q$-data) simultaneously. The symmetric shadow
$(\int_{E_\tau}^{\mathrm{ord}})_{\Sigma_n}$ retains only
the scalar traces (characters of $U_q(\mathfrak{sl}_2)$
representations), which are the WZW conformal blocks; the
kernel $\ker(\av)$ carries the matrix-valued monodromy from
which knot invariants on the torus are extracted.
\end{remark}

\begin{remark}[Roots of unity]
\label{rem:ell-roots-of-unity}
At integer level $k \in \ZZ_{>0}$, the quantum group
parameter is $q = e^{2\pi i/(k+2)}$ and
$q^{k+2} = 1$. The KZB local system acquires finite
monodromy, and the ordered chiral homology at each arity
becomes a finite-dimensional representation of the mapping
class group of $E_\tau$ with $n$ marked points. The
$k + 1$ integrable modules $V_0, \ldots, V_k$ of
$U_q(\mathfrak{sl}_2)$ correspond to the truncation of the
arity-by-arity decomposition: at arity $n > k$, the
representation theory constrains certain components to
vanish.
\end{remark}

\begin{remark}[The elliptic $R$-matrix as ordered datum]
\label{rem:ell-r-matrix-ordered}
The Belavin elliptic $r$-matrix
$r^{\mathrm{ell}}(z, \tau) = k \cdot r^{\mathrm{Belavin}}
(z, \tau)$ is the arity-$2$ projection of the genus-$1$
Maurer--Cartan element
$\Theta_\cA^{(1)} \in \gEone$. It satisfies the classical
Yang--Baxter equation with spectral parameter (verified
numerically in \texttt{elliptic\_rmatrix\_shadow.py}),
and its averaging gives the genus-$1$ shadow:
% AP1: kappa = 3(k+2)/4 for sl_2. k=0 -> 3/2.
$\av(r^{\mathrm{ell}}) + \dim(\mathfrak{sl}_2)/2
= 3k/4 + 3/2 = 3(k+2)/4 = \kappa$.
The Sugawara shift $\dim(\fg)/2 = 3/2$ is present at
genus~$1$ exactly as at genus~$0$: the shift comes from
the simple-pole self-contraction through the adjoint
Casimir eigenvalue $2h^\vee = 4$, which is a local
computation independent of the global topology of the
base curve.
\end{remark}


% ================================================================
\section{Conclusion and open problems}
\label{sec:conclusion}

This paper defines ordered chiral homology
$\int_X^{\mathrm{ord}} \cA$ for $\Eone$-chiral algebras
on algebraic curves, constructs the chiral
$\Ethree$-algebra
$\CE^{\mathrm{ch}}_\hbar(\fg)$ as the holomorphic
refinement of the Costello--Francis--Gwilliam perturbative
Chern--Simons algebra, and establishes the chiral quantum
group equivalence (Theorem~\ref{thm:chiral-qg-equiv})
relating vertex $R$-matrices, $\Ainf$-structures, and
chiral coproducts as three faces of a single ordered bar
complex.
Five explicit computations across the shadow
classification (Heisenberg: class~$G$; affine Kac--Moody:
class~$L$; $\beta\gamma$: class~$C$; Virasoro:
class~$M$; Yangian: genuinely $\Eone$) verify the
formality bridge, the regular/irregular dichotomy,
and the Stokes phenomena that distinguish gauge theory
from gravity at the level of the KZ $\cD$-module.
Section~\ref{sec:elliptic-ordered} extends the construction
to genus~$1$, computing
$\int_{E_\tau}^{\mathrm{ord}} Y_\hbar(\mathfrak{sl}_2)$
arity by arity: the Belavin elliptic $r$-matrix replaces
the rational $r$-matrix, the KZB connection replaces the
KZ connection, and the $B$-cycle monodromy produces the
quantum group parameter $q = e^{2\pi i \hbar}$.

The following problems remain open.

\begin{enumerate}[label=\textup{(\alph*)}]

\item \textup{(The $\Ethree$ identification
for affine Kac--Moody.)}
Conjecture~\ref{conj:e3-identification} asserts that the
derived chiral centre
$Z^{\mathrm{der}}_{\mathrm{ch}}(V_k(\fg))$ and the
CFG algebra $\cA^\lambda$ are isomorphic as
$\Ethree$-algebras over the formal deformation space
$\hbar H^3(\fg)[[\hbar]]$.
This paper proves that the deformation spaces agree
(both one-dimensional at each order in $\hbar$ for
simple~$\fg$) and that the formal disk global sections
match at the level of $\Pthree$-brackets
(\S\ref{subsec:e3-explicit}).
What is missing is a chain-level comparison that
identifies the $\Ethree$-structures as families, not
merely fiber by fiber: is there a natural map between the
two $\Ethree$-algebras that is a quasi-isomorphism at
each non-critical level $k \neq -h^\vee$?

\item \textup{(Holomorphic Chern--Simons invariants of
curves.)}
The chiral $\Ethree$-algebra
$\CE^{\mathrm{ch}}_\hbar(\fg)$ on a smooth proper
curve~$X$ of genus~$g$ should produce, via the
higher-genus ordered chiral homology and factorisation on
$\overline{\cM}_{g,n}$, invariants of $(X, \fg, k)$
that refine Reshetikhin--Turaev invariants in the same
way that perturbative Chern--Simons refines Witten's
partition function.
Concretely: does the genus-$g$ ordered chiral homology
$\int_X^{\mathrm{ord}} \CE^{\mathrm{ch}}_\hbar(\fg)$
recover the perturbative Chern--Simons partition function
$Z_{\mathrm{CS}}(X, \fg, k)$ as its $\Sigma_n$-coinvariant
shadow, and what invariants live in the kernel of
the averaging map $\av$?

\item \textup{(The $\cW_{1+\infty}$ chiral quantum group.)}
Example~\ref{ex:w-infty-coha} identifies
$\cW_{1+\infty}[\Psi]$
as a chiral quantum group whose CoHA origin provides the
$R$-matrix, coproduct, and $\Ainf$-structure.
Two questions are unresolved.
First, the CoHA construction of
Schiffmann--Mellit--Minets--Vasserot~\cite{SMMV23}
gives the positive half
$U(\cW_{1+\infty})^+$; does the full chiral quantum group
structure (including the antipode and the Drinfeld double)
admit a CoHA interpretation?
Second, the KZ connection for $\cW_N$ at finite~$N$ has
irregular singularities of Poincar\'e rank $N - 1$;
what are its Stokes multipliers, and do they encode
the higher-spin BPS spectrum?

\item \textup{(Higher-genus ordered chiral homology.)}
Section~\ref{sec:ord-ch-hom} defines ordered chiral
homology at genus~$0$, and
Section~\ref{sec:elliptic-ordered} constructs the
genus-$1$ theory
$\int_{E_\tau}^{\mathrm{ord}} Y_\hbar(\mathfrak{sl}_2)$
explicitly, computing the arity-by-arity decomposition
through arities~$0$, $1$, and~$2$. The curvature
obstruction $d^2_{\mathrm{fib}} = \kappa \cdot \omega_g$
at genus $g \geq 2$, where $\omega_g = c_1(\lambda)$ is
the Hodge class on $\overline{\cM}_g$, requires a
period-corrected differential whose coefficients involve
the $(g \times g)$-matrix $(\mathrm{Im}\, \Omega)^{-1}$
of the Siegel period matrix, not the scalar
$1/\mathrm{Im}(\tau)$ of the genus-$1$ case.
Does the ordered chiral chain complex at genus~$g \geq 2$
define a flat connection over $\overline{\cM}_g$
compatible with the Mumford--Morita class $\kappa_1$,
and does the resulting system of conformal blocks satisfy
factorisation at the boundary divisors of
$\overline{\cM}_g$?

\item \textup{(Trigonometric and elliptic $R$-matrices
beyond $\mathfrak{sl}_2$.)}
Section~\ref{sec:elliptic-ordered} treats the Belavin
elliptic $R$-matrix for $\mathfrak{sl}_2$ on $E_\tau$.
The trigonometric $R$-matrix
$R^{\mathrm{trig}}(u) = q^{\Omega/u}$ (quantum affine
algebras) and the elliptic Belavin $R$-matrix for
$\mathfrak{sl}_N$ ($N \geq 3$) produce richer ordered
chiral homology with qualitatively different behaviour:
higher-order Stokes phenomena, theta-function monodromy,
and flat sections valued in elliptic hypergeometric
integrals.
Does the regular/irregular dichotomy of
Theorem~\ref{thm:main-ord-ch} extend to a trichotomy
(regular, irregular-rational, irregular-elliptic)
correlated with the three types of quantum groups?
For $\mathfrak{sl}_N$, what is the arity-$2$ de~Rham
cohomology of the KZB local system of rank $N^2$ on
$E_\tau \setminus \{0\}$, and how does the isospin
decomposition into irreducibles of
$V^{\otimes 2} = \operatorname{Sym}^2 V \oplus
\bigwedge^2 V$ generalise?

\item \textup{(Categorification.)}
Ben-Zvi--Brochier--Jordan~\cite{BBJ15,BBJ16} integrate
braided monoidal categories over surfaces; ordered chiral
homology integrates chiral algebras over curves.
The expected bridge
(Remark~\ref{rem:bbj},
equation~\eqref{eq:bbj-comparison}) identifies the
Hochschild homology of the categorical invariant
$\int_\Sigma \mathrm{Rep}_q(\fg)$ with the ordered
chiral homology
$\int_X^{\mathrm{ord}} V_k(\fg)$.
Is there a direct functor from the BBJ categorical
factorisation homology to ordered chiral homology
that intertwines the $\Ethree$-structures, and does
the kernel of this functor detect the higher categorical
data (objects of $\int_\Sigma \mathrm{Rep}_q(\fg)$
that are invisible at the chain level)?

\item \textup{(Refinements of the topological enhancement.)}
Khan--Zeng~\cite{KhanZeng25} prove the topological
enhancement at the classical (PVA) level;
Proposition~\ref{prop:khan-zeng-topological} extends
this to the full quantum vertex algebra $V_k(\fg)$
at non-critical level, using the exactness of the
Sugawara element in $V_k(\fg)$.
Two refinements remain open: (a)~does the topological
$\Ethree$-structure on
$\CE^{\mathrm{ch}}_\hbar(\fg)$ extend to a
topological $\Ethree$-structure on the higher-genus
ordered chiral homology
$\int_X^{\mathrm{ord}} \CE^{\mathrm{ch}}_\hbar(\fg)$
for $g \geq 2$, and (b)~at the critical level
$k = -h^\vee$, where the Sugawara element is undefined,
is the resulting holomorphic $\Ethree$-algebra
distinguished from the topological one by a computable
invariant (a secondary characteristic class of the
chiral $\Ethree$-algebra)?

\end{enumerate}


% ################################################################
% ################################################################
%
% [Part II deleted: the correct setting for holomorphic

\begin{thebibliography}{99}

\bibitem{Arnold69}
V.~I.~Arnold,
\emph{The cohomology ring of the group of dyed braids},
Mat. Zametki \textbf{5} (1969), 227--231.

\bibitem{BBJ15}
D.~Ben-Zvi, A.~Brochier, and D.~Jordan,
\emph{Integrating quantum groups over surfaces},
J. Topol. \textbf{11} (2018), 874--917;
arXiv:1501.04652.

\bibitem{BBJ16}
D.~Ben-Zvi, A.~Brochier, and D.~Jordan,
\emph{Quantum character varieties and braided module
categories},
Selecta Math. (N.S.) \textbf{24} (2018), 4711--4748;
arXiv:1606.04769.

\bibitem{BD04}
A.~Beilinson and V.~Drinfeld,
\emph{Chiral algebras},
American Mathematical Society Colloquium Publications,
vol.~51, AMS, Providence, RI, 2004.

\bibitem{Bernard88}
D.~Bernard,
\emph{On the Wess--Zumino--Witten models on the torus},
Nucl. Phys. B \textbf{303} (1988), 77--93.

\bibitem{CEE09}
D.~Calaque, B.~Enriquez, and P.~Etingof,
\emph{Universal KZB equations: the elliptic case},
in Algebra, Arithmetic, and Geometry (Manin Festschrift),
Progress in Math., vol.~269, Birkh\"auser, 2009,
pp.~165--266.

\bibitem{Kadeishvili82}
T.~Kadeishvili,
\emph{The algebraic structure in the homology of an
$A(\infty)$-algebra},
Soobshch. Akad. Nauk Gruzin. SSR \textbf{108} (1982),
249--252.

\bibitem{Kontsevich99}
M.~Kontsevich,
\emph{Operads and motives in deformation quantization},
Lett. Math. Phys. \textbf{48} (1999), 35--72;
arXiv:math/9904055.

\bibitem{Tamarkin03}
D.~Tamarkin,
\emph{Formality of chain operad of little discs},
Lett. Math. Phys. \textbf{66} (2003), 65--72;
arXiv:math/9809164.


\bibitem{Drinfeld90}
V.~G.~Drinfeld,
\emph{Quasi-Hopf algebras},
Leningrad Math. J. \textbf{1} (1990), 1419--1457.


\bibitem{CFG26}
K.~Costello, J.~Francis, and O.~Gwilliam,
\emph{Chern--Simons factorization algebras and knot
polynomials},
arXiv:2602.12412, 2026.

\bibitem{CG17}
K.~Costello and O.~Gwilliam,
\emph{Factorization algebras in quantum field theory},
vols.~1--2, Cambridge University Press, 2017 and 2021.

\bibitem{CWY1}
K.~Costello, E.~Witten, and M.~Yamazaki,
\emph{Gauge theory and integrability, I},
Notices ICCM \textbf{6} (2018), no.~1, 46--119;
arXiv:1709.09993.

\bibitem{CWY2}
K.~Costello, E.~Witten, and M.~Yamazaki,
\emph{Gauge theory and integrability, II},
Notices ICCM \textbf{6} (2018), no.~1, 120--146;
arXiv:1802.01579.

\bibitem{DAgnoloKashiwara16}
A.~D'Agnolo and M.~Kashiwara,
\emph{Riemann--Hilbert correspondence for holonomic
$D$-modules},
Publ. Math. Inst. Hautes \'Etudes Sci. \textbf{123}
(2016), 69--197.

\bibitem{EK96}
P.~Etingof and D.~Kazhdan,
\emph{Quantization of Lie bialgebras, I},
Selecta Math. (N.S.) \textbf{2} (1996), 1--41.

\bibitem{EK00}
P.~Etingof and D.~Kazhdan,
\emph{Quantization of Lie bialgebras, V: Quantum vertex
operator algebras},
Selecta Math. (N.S.) \textbf{6} (2000), 105--130.

\bibitem{EV98}
P.~Etingof and A.~Varchenko,
\emph{Solutions of the quantum dynamical Yang--Baxter
equation and dynamical quantum groups},
Comm. Math. Phys. \textbf{196} (1998), 591--640.

\bibitem{FBZ04}
E.~Frenkel and D.~Ben-Zvi,
\emph{Vertex algebras and algebraic curves},
second ed., Mathematical Surveys and Monographs, vol.~88,
AMS, Providence, RI, 2004.

\bibitem{Francis2013}
J.~Francis,
\emph{The tangent complex and Hochschild cohomology of
$\En$-rings},
Compos. Math. \textbf{149} (2013), 430--480.

\bibitem{GR17}
D.~Gaiotto and M.~Rap\v{c}\'ak,
\emph{Vertex algebras at the corner},
J. High Energy Phys. \textbf{2019}(1), 160;
arXiv:1703.00982.

\bibitem{HTT08}
R.~Hotta, K.~Takeuchi, and T.~Tanisaki,
\emph{$D$-modules, perverse sheaves, and representation
theory},
Birkh\"auser, Boston, 2008.

\bibitem{HA}
J.~Lurie,
\emph{Higher algebra},
available at \texttt{math.ias.edu/\textasciitilde lurie/},
2017.

\bibitem{Kashiwara03}
M.~Kashiwara,
\emph{$D$-modules and microlocal calculus},
Translations of Mathematical Monographs, vol.~217,
AMS, 2003.

\bibitem{KhanZeng25}
A.~Z.~Khan and K.~Zeng,
\emph{Poisson vertex algebras and three-dimensional gauge
theory},
arXiv:2502.13227, 2025.

\bibitem{KS11}
M.~Kontsevich and Y.~Soibelman,
\emph{Cohomological Hall algebra, exponential Hodge
structures and motivic Donaldson--Thomas invariants},
Commun. Number Theory Phys. \textbf{5} (2011), 231--352;
arXiv:1006.2706.

\bibitem{Latyntsev23}
A.~Latyntsev,
\emph{Chiral homology and derived chiral algebras},
arXiv:2312.07274, 2023.

\bibitem{Felder94}
G.~Felder,
\emph{Conformal field theory and integrable systems
associated to elliptic curves},
in Proc. ICM Z\"urich 1994, Birkh\"auser, 1995,
pp.~1247--1255.

\bibitem{FelderVarchenko96}
G.~Felder and A.~Varchenko,
\emph{On representations of the elliptic quantum group
$E_{\tau,\eta}(\mathfrak{sl}_2)$},
Commun. Math. Phys. \textbf{181} (1996), 741--761;
q-alg/9601003.

\bibitem{Lorgat26I}
R.~Lorgat,
\emph{Modular Koszul duality},
in preparation, 2026.

\bibitem{Lorgat26II}
R.~Lorgat,
\emph{$A_\infty$ chiral algebras and $3$D
holomorphic-topological QFT},
in preparation, 2026.

\bibitem{Lorgat26III}
R.~Lorgat,
\emph{Calabi--Yau categories, quantum groups, and BPS
algebras},
in preparation, 2026.

\bibitem{Molev07}
A.~Molev,
\emph{Yangians and classical Lie algebras},
Mathematical Surveys and Monographs, vol.~143,
AMS, Providence, RI, 2007.

\bibitem{RSYZ20}
M.~Rap\v{c}\'ak, Y.~Soibelman, Y.~Yang, and G.~Zhao,
\emph{Cohomological Hall algebras, vertex algebras and
instantons},
Commun. Math. Phys. \textbf{376} (2020), 1803--1873;
arXiv:1810.10402.

\bibitem{Mebkhout84}
Z.~Mebkhout,
\emph{Une \'equivalence de cat\'egories
et une autre \'equivalence de cat\'egories},
Compos. Math. \textbf{51} (1984), 51--62; 63--68.

\bibitem{Sabbah13}
C.~Sabbah,
\emph{Introduction to Stokes structures},
Lecture Notes in Math., vol.~2060, Springer, Berlin, 2013.

\bibitem{Sinha04}
D.~Sinha,
\emph{Manifold-theoretic compactifications of configuration
spaces},
Selecta Math. (N.S.) \textbf{10} (2004), 391--428.

\bibitem{SMMV23}
O.~Schiffmann, A.~Mellit, A.~Minets, and E.~Vasserot,
\emph{Coherent sheaves on surfaces, COHAs and deformed
$W_{1+\infty}$-algebras},
arXiv:2311.13415, 2023.

\bibitem{SV13}
O.~Schiffmann and E.~Vasserot,
\emph{Cherednik algebras, W-algebras and the equivariant
cohomology of the moduli space of instantons on
$\mathbb{A}^2$},
Publ. Math. Inst. Hautes \'Etudes Sci. \textbf{118}
(2013), 213--342; arXiv:1202.2756.

\bibitem{Wasow65}
W.~Wasow,
\emph{Asymptotic expansions for ordinary differential
equations},
Interscience Publishers, New York, 1965;
reprinted Dover, 1987.

\bibitem{YZ18a}
Y.~Yang and G.~Zhao,
\emph{Cohomological Hall algebras and affine quantum
groups},
Selecta Math. (N.S.) \textbf{24} (2018), 1093--1119;
arXiv:1604.01865.

\end{thebibliography}

\end{document}
